% ===========================================================================
%  BOKY MATEMATIKA — PREMIÈRE S
%  Programme d'Études — Enseignement Secondaire Général, Madagascar
%
%  Compilation :  ./compiler.sh eleve | prof | les-deux | propre
%  Moteur      :  XeLaTeX (latexmk -xelatex)
% ===========================================================================
\documentclass[11pt, twoside, openright]{book}

\usepackage{preambule/bmt-style}
\usepackage{preambule/bmt-macros}
\usepackage{preambule/bmt-blocks}


\begin{document}

\frontmatter
\clearpage% --- Couverture -------------------------------------------------------------
\begin{titlepage}
\thispagestyle{empty}
\begin{tikzpicture}[remember picture, overlay]

  % Fond papier
  \fill[bmtPapier] (current page.south west) rectangle (current page.north east);

  % Bloc principal sarcelle
  \fill[bmtSarcelle]
    (current page.north west) rectangle
    ([yshift=-19.4cm]current page.north east);

  % Liseré de série, sur toute la hauteur du bord extérieur
  \fill[bmtRavenala]
    ([xshift=-0.55cm]current page.north east) rectangle
    ([yshift=-19.4cm]current page.north east);

  % --- Composition géométrique : un signe par domaine de l'année ------------
  % Analyse : une courbe et sa tangente en un point
  \begin{scope}[shift={([xshift=1.9cm, yshift=-9.1cm]current page.north west)}]
    \clip (-0.15,-0.2) rectangle (6.9,4.6);
    \draw[bmtPapier, line width=1.5pt, opacity=0.92, smooth, samples=160,
          domain=0.15:6.75]
      plot (\x, {0.55 + 2.9/(1 + exp(-1.35*(\x-3.1)))});
    % la tangente au point d'inflexion
    \draw[bmtOcre, line width=1.1pt, opacity=0.95]
      (1.15, 0.0914375) -- (5.05, 3.9085625);
    \fill[bmtPapier] (3.1, 2.0) circle (0.085);
    \draw[bmtPapier, line width=0.7pt, dash pattern=on 2.4pt off 2.4pt,
          opacity=0.6] (3.1, 0) -- (3.1, 2.0);
  \end{scope}

  % Algèbre : les premiers termes d'une suite, et l'escalier qui les relie
  \begin{scope}[shift={([xshift=10.45cm, yshift=-5.35cm]current page.north west)}]
    \draw[bmtPapier, line width=0.7pt, opacity=0.35] (-0.15,0) -- (3.15,0);
    \draw[bmtOcre, line width=1.1pt]
      (0,0.28) -- (0.72,0.28) -- (0.72,0.62) -- (1.44,0.62) --
      (1.44,1.12) -- (2.16,1.12) -- (2.16,1.78) -- (2.88,1.78);
    \foreach \i/\h in {0/0.28, 1/0.62, 2/1.12, 3/1.78}{
      \fill[bmtPapier] (\i*0.72, \h) circle (0.09);
      \draw[bmtPapier, line width=0.5pt, opacity=0.4]
        (\i*0.72, 0) -- (\i*0.72, \h);}
  \end{scope}

  % Géométrie : un cube en perspective cavalière
  \begin{scope}[shift={([xshift=14.55cm, yshift=-8.35cm]current page.north west)},
                scale=1.0]
    \draw[bmtPapier, line width=1.1pt, opacity=0.9]
      (0,0) -- (2.1,0) -- (2.1,2.1) -- (0,2.1) -- cycle;
    \draw[bmtPapier, line width=1.1pt, opacity=0.9]
      (0.78,0.78) -- (2.88,0.78) -- (2.88,2.88) -- (0.78,2.88) -- cycle;
    \draw[bmtPapier, line width=1.1pt, opacity=0.9]
      (0,0) -- (0.78,0.78)  (2.1,0) -- (2.88,0.78)
      (2.1,2.1) -- (2.88,2.88)  (0,2.1) -- (0.78,2.88);
    \draw[bmtRavenala, line width=1.3pt, -{Stealth[length=4.5pt]}]
      (0,0) -- (1.35,0);
    \draw[bmtRavenala, line width=1.3pt, -{Stealth[length=4.5pt]}]
      (0,0) -- (0,1.35);
    \draw[bmtRavenala, line width=1.3pt, -{Stealth[length=4.5pt]}]
      (0,0) -- (0.62,0.62);
  \end{scope}

  % Traitement des données : une boîte à moustaches
  \begin{scope}[shift={([xshift=10.55cm, yshift=-11.25cm]current page.north west)}]
    \draw[bmtPalissandre, line width=1.1pt] (0,0) -- (0.95,0);
    \draw[bmtPalissandre, line width=1.1pt] (0,-0.24) -- (0,0.24);
    \draw[bmtPapier, line width=1.2pt, opacity=0.92]
      (0.95,-0.42) rectangle (3.05,0.42);
    \draw[bmtOcre, line width=1.4pt] (1.95,-0.42) -- (1.95,0.42);
    \draw[bmtPalissandre, line width=1.1pt] (3.05,0) -- (4.1,0);
    \draw[bmtPalissandre, line width=1.1pt] (4.1,-0.24) -- (4.1,0.24);
  \end{scope}

  % --- Titre ---------------------------------------------------------------
  \node[anchor=north west, text=bmtPapier, align=left,
        font=\sffamily\fontsize{12.5}{15}\selectfont, opacity=0.85]
    at ([xshift=1.95cm, yshift=-13.35cm]current page.north west)
    {BOKY MATEMATIKA};

  \node[anchor=north west, text=bmtPapier, align=left,
        font=\sffamily\bfseries\fontsize{58}{62}\selectfont]
    at ([xshift=1.8cm, yshift=-14.05cm]current page.north west)
    {Première~S};

  \draw[bmtOcre, line width=3pt]
    ([xshift=1.95cm, yshift=-16.75cm]current page.north west) --
    ([xshift=5.65cm, yshift=-16.75cm]current page.north west);

  \node[anchor=north west, text=bmtPapier, align=left, text width=16.4cm,
        font=\sffamily\fontsize{11}{15}\selectfont, opacity=0.92]
    at ([xshift=1.95cm, yshift=-17.35cm]current page.north west)
    {Analyse \enspace\textbullet\enspace Algèbre \enspace\textbullet\enspace Géométrie
     \enspace\textbullet\enspace Traitement des données et probabilités};

  % --- Bas de couverture ---------------------------------------------------
  \node[anchor=north west, text=bmtEncre, align=left, text width=12cm,
        font=\sffamily\fontsize{11}{15}\selectfont]
    at ([xshift=1.9cm, yshift=-20.6cm]current page.north west)
    {\textbf{Brouillon de travail}\\
     D’après le Programme d’Études T11 — version expérimentale\\
     Enseignement Secondaire Général — Madagascar};

  % Marque de l'éditeur
  \draw[bmtSarcelle, line width=1.6pt]
    ([xshift=1.95cm, yshift=-22.75cm]current page.north west) --
    ([xshift=3.35cm, yshift=-22.75cm]current page.north west);
  \node[anchor=north west, text=bmtSarcelle, align=left,
        font=\sffamily\bfseries\fontsize{12}{15}\selectfont]
    at ([xshift=1.9cm, yshift=-23.15cm]current page.north west)
    {ÉDITIONS MADA-BOKY};

  % Les quatre couleurs des quatre parties, en signature de bas de page
  \foreach \c [count=\i from 0] in
      {bmtSarcelle, bmtGrenat, bmtRavenala, bmtPalissandre}{
    \fill[\c] ([xshift={1.9cm + \i*0.95cm}, yshift=2.5cm]current page.south west)
      rectangle ++(0.62cm, 0.62cm);}

  \node[anchor=south west, text=bmtGris, align=left,
        font=\sffamily\fontsize{9}{12}\selectfont]
    at ([xshift=1.9cm, yshift=1.5cm]current page.south west)
    {\ifbmtprof ÉDITION DU PROFESSEUR\else ÉDITION DE L'ÉLÈVE\fi};

\end{tikzpicture}
\end{titlepage}

\clearpage% --- Page de titre ----------------------------------------------------------
\cleardoublepage
\thispagestyle{empty}
\begingroup
\raggedright
\vspace*{3.2cm}

{\sffamily\fontsize{11}{14}\selectfont\color{bmtGris} BOKY MATEMATIKA\par}
\vspace{0.5em}
{\sffamily\bfseries\fontsize{40}{44}\selectfont\color{bmtSarcelle} Première~S\par}
\vspace{0.9em}
{\color{bmtOcre}\rule{4cm}{2.6pt}\par}
\vspace{1.4em}
{\sffamily\fontsize{13}{18}\selectfont
 Analyse \quad Algèbre \quad Géométrie \quad Traitement des données et probabilités\par}

\vspace{2.6cm}
{\normalsize
 Manuel de mathématiques pour la classe de T11, série S, élaboré à partir
 du Programme d'Études conservé dans le dépôt Mada-Boky.\par
 \medskip
 \textbf{Brouillon de travail.} Les contrôles et les points restant à vérifier
 sont consignés dans les rapports du projet.\par}

\vfill

{\color{bmtSarcelle}\rule{1.4cm}{1.6pt}\par}
\vspace{0.5em}
{\sffamily\bfseries\fontsize{13}{17}\selectfont\color{bmtSarcelle} ÉDITIONS MADA-BOKY\par}
\vspace{0.3em}
{\sffamily\footnotesize\color{bmtGris} Antananarivo\par}

\vspace{1.2em}
{\sffamily\footnotesize\color{bmtGris}
 \ifbmtprof Édition du professeur — corrigés et notes de conduite de classe
 \else Édition de l'élève\fi\par}
\endgroup
\cleardoublepage

\clearpage\chapter*{Avant-propos}\markboth{Avant-propos}{Avant-propos}\addcontentsline{toc}{chapter}{Avant-propos}
Ce manuel s’adresse aux élèves de Première scientifique à Madagascar. Il rassemble les quatre domaines du Programme d’Études T11 : analyse, algèbre, géométrie, données et probabilités.
Le cours introduit les outils, les exemples montrent comment les utiliser, puis les exercices demandent de choisir une méthode et de justifier les réponses. Les résultats généraux dont la preuve dépasse le niveau sont signalés comme admis. Une figure accompagne le raisonnement; elle ne remplace pas une démonstration.
Les deux éditions proviennent des mêmes sources. L’édition du professeur ajoute les corrections et des notes pour conduire le travail. Les exercices et les évaluations de ce volume sont originaux; aucune provenance d’annale n’est attribuée sans document vérifié.
\section*{Pour l’élève}
Commencez par les rappels. Si un calcul ou une condition de domaine vous pose problème, reprenez-le avant de poursuivre. Dans un exercice, écrivez ce que vous cherchez et pourquoi chaque transformation est autorisée. Le résultat seul ne permet pas toujours de comprendre une erreur.
\section*{Pour le professeur}
Les parties du livre sont un ordre de rangement. La progression indicative propose de commencer par la logique, les équations et la trigonométrie avant l’étude des fonctions. Les évaluations et les sujets de synthèse sont modulables selon les notions réellement enseignées.
\section*{Références}
La référence de contenu est le Programme d’Études T11 conservé dans le dépôt Mada-Boky, partie scientifique, pages imprimées 352–363. Sa couverture porte « version expérimentale ». Le programme antérieur de Première et la répartition 2025–2026 servent de documents de comparaison. Ce volume est un brouillon de travail : ses contrôles sont consignés dans les rapports associés au projet.

\clearpage\chapter*{Comment utiliser ce livre}\markboth{Comment utiliser ce livre}{Comment utiliser ce livre}\addcontentsline{toc}{chapter}{Comment utiliser ce livre}
\begin{revision}Les rappels indiquent les connaissances nécessaires pour lire le chapitre.\end{revision}
\begin{automatismes}Les questions courtes réactivent les calculs et les représentations déjà appris.\end{automatismes}
\begin{methode}{lire le cours}Repérez les objets d’une définition, puis les hypothèses d’une propriété. Dans une preuve, vérifiez chaque passage. Les exemples sont résolus; essayez de refaire les calculs avant de les lire.\end{methode}
\begin{visualisation}Un graphique se lit avec ses axes, ses unités et son domaine. La légende précise ce qui est représenté.\end{visualisation}
\begin{entrainement}Les exercices portent un numéro automatique à l’intérieur de chaque chapitre. Ils alternent application, choix de méthode, preuve, contre-exemple, modélisation et synthèse. Les énoncés sont identiques dans les deux éditions.\end{entrainement}
\begin{jemeteste}Le bilan de fin de chapitre rappelle ce que vous devez savoir expliquer et utiliser.\end{jemeteste}
\begin{notepourlaclasse}Les corrigés renvoient au numéro de l’exercice par un lien. Les notes aident à repérer les prérequis, les erreurs possibles et les choix de séance. Elles décrivent une intention pédagogique, sans prétendre à une efficacité mesurée en classe.\end{notepourlaclasse}
Chaque partie se termine par une évaluation sur vingt points. Deux sujets de synthèse proposent des reprises cumulatives. La calculatrice y est autorisée; les justifications restent nécessaires.


{\hypersetup{linkcolor=bmtEncre}\tableofcontents}

\mainmatter

% ===========================================================================
\bmtpartie{bmtSarcelle}{I}{Analyse}{Fonctions numériques}{%
  Étudier une fonction, c'est répondre à cinq questions dans le même ordre :
  où est-elle définie, vers quoi tend-elle aux bornes, comment varie-t-elle,
  quels points remarquables porte-t-elle, à quoi ressemble sa courbe. Les
  cinq chapitres qui suivent installent ces cinq questions une à une. À la
  fin, vous les poserez sans y penser.}

\clearpage% ===========================================================================
%  Chapitre 1 — Généralités sur les fonctions numériques
%  Partie I — Analyse
% ===========================================================================
\bmtchapitre{Généralités sur les fonctions numériques}{%
  Lire une image et un antécédent sur une courbe comme sur une expression,
  déterminer un ensemble de définition, établir qu'une fonction est paire,
  impaire ou périodique, et en tirer l'intervalle sur lequel il suffit de
  l'étudier.}{Analyse}

\label{t11s-c01}

En seconde, une fonction était le plus souvent une machine à calculer : on
donnait un nombre, elle en rendait un autre. Cette année, on change de point
de vue. Avant de calculer quoi que ce soit, on commence par délimiter le
terrain — pour quels nombres la machine fonctionne-t-elle ? — puis on cherche
ce que la fonction répète, ce qu'elle reflète, ce qu'elle laisse invariant.
Une fonction paire n'a besoin d'être étudiée que sur la moitié de son
domaine ; une fonction périodique, que sur une seule période. Le chapitre
entier tient dans cette économie.

% ---------------------------------------------------------------------------
\begin{revision}[ce qui doit répondre tout de suite]
Vous aurez besoin, dès les premières pages, des acquis de seconde suivants.

\begin{enumerate}[label=\textbf{\arabic*.}, leftmargin=1.5em, itemsep=0.25em]
  \item Les six fonctions de référence : $x \mapsto ax+b$, $x \mapsto x^{2}$,
        $x \mapsto x^{3}$, $x \mapsto \dfrac{1}{x}$, $x \mapsto \sqrt{x}$ et
        $x \mapsto \abs{x}$, avec leur courbe et leur sens de variation.
  \item Résoudre une inéquation du second degré et dresser un tableau de
        signes.
  \item La condition d'existence de $\sqrt{A}$ : il faut $A \geq 0$. La
        condition d'existence de $\dfrac{1}{B}$ : il faut $B \neq 0$.
  \item Sur le cercle trigonométrique : le radian, les angles remarquables,
        et l'égalité $\cos^{2}x + \sin^{2}x = 1$.
  \item Lire les coordonnées d'un point sur un graphique et savoir ce que
        signifie « le point $A(2\,;\,5)$ appartient à la courbe ».
\end{enumerate}
\end{revision}

\begin{automatismes}\phantomsection\label{t11s-c01-auto}
À faire de tête, sans écrire autre chose que la réponse.

\begin{multicols}{2}
\begin{enumerate}[label=\textbf{\arabic*.}, leftmargin=1.5em, itemsep=0.15em]
  \item Pour quelles valeurs de $x$ a-t-on $x - 3 \geq 0$ ?
  \item Pour quelles valeurs de $x$ a-t-on $x^{2} - 4 = 0$ ?
  \item Calculer $\abs{-7}$, puis $\abs{3 - 10}$.
  \item Calculer $\sqrt{49}$, puis $\sqrt{(-5)^{2}}$.
  \item Résoudre $2x + 1 \neq 0$.
  \item Donner le signe de $-x^{2} - 1$ selon $x$.
  \item Calculer $\cos(0)$, $\sin\!\left(\dfrac{\pi}{2}\right)$,
        $\cos(\pi)$.
  \item Développer $(-x)^{2}$, puis $(-x)^{3}$.
  \item Résoudre $x^{2} = 9$.
  \item Le nombre $-4$ a-t-il une racine carrée ?
\end{enumerate}
\end{multicols}
\end{automatismes}

% ===========================================================================
\section{Fonction, image et antécédent}

\begin{definition}[fonction numérique d'une variable réelle]
Soit $\Df$ une partie de $\R$. Définir une \textbf{fonction numérique} $f$ sur
$\Df$, c'est associer à tout réel $x$ de $\Df$ un unique réel, noté $f(x)$.
On écrit
\[
  f : \begin{array}{ccc}
        \Df & \longrightarrow & \R \\[0.15em]
        x   & \longmapsto     & f(x)
      \end{array}
\]
Le réel $f(x)$ s'appelle l'\textbf{image} de $x$ par $f$. Réciproquement, si
$y$ est un réel et si $x$ est un élément de $\Df$ tel que $f(x) = y$, on dit
que $x$ est un \textbf{antécédent} de $y$ par $f$.
\end{definition}

\begin{remarque}[le mot « unique » n'est pas décoratif]
Un réel de $\Df$ a une image et une seule. En revanche, un réel $y$ peut
avoir plusieurs antécédents, ou n'en avoir aucun : la fonction carré donne
deux antécédents à $4$, et aucun à $-1$.
\end{remarque}

\begin{visualisation}[tout se lit verticalement, puis horizontalement]
\begin{center}
\begin{tikzpicture}[scale=0.95]
  \begin{axis}[
      axis lines = middle, xlabel = {$x$}, ylabel = {$y$},
      xmin = -3.4, xmax = 3.4, ymin = -1.6, ymax = 4.6,
      xtick = {-3,-2,-1,1,2,3}, ytick = {1,2,3,4},
      width = 9.2cm, height = 6.4cm,
      tick label style = {font=\footnotesize},
      label style = {font=\footnotesize},
      axis line style = {bmtGris},
      every tick/.style = {bmtGris},
    ]
    \addplot[bmtSarcelle, very thick, samples = 200, domain = -3.1:3.1]
      {0.55*x^2 - 0.25*x};
    % image de 2
    \draw[bmtOcre, thick, dashed] (axis cs:2,0) -- (axis cs:2,1.7);
    \draw[bmtOcre, thick, dashed] (axis cs:2,1.7) -- (axis cs:0,1.7);
    \fill[bmtOcre] (axis cs:2,1.7) circle (1.9pt);
    % antécédents de 1.7
    \draw[bmtPalissandre, thick, dashed] (axis cs:-1.5454545,0) -- (axis cs:-1.5454545,1.7);
    \fill[bmtPalissandre] (axis cs:-1.5454545,1.7) circle (1.9pt);
    \node[font=\footnotesize, bmtOcre, anchor=west] at (axis cs:2.1,1.9) {$f(2)$};
  \end{axis}
\end{tikzpicture}
\end{center}

\vspace{0.2em}
\noindent
Pour lire l'\emph{image} de $2$ : on part de $2$ sur l'axe des abscisses, on
monte jusqu'à la courbe, on lit sur l'axe des ordonnées. Pour lire les
\emph{antécédents} de $1{,}7$ : on part de $1{,}7$ sur l'axe des ordonnées,
on se déplace horizontalement, et on relève l'abscisse de \emph{chaque} point
d'intersection rencontré. Il y en a deux ici.
\end{visualisation}

\begin{exemple}[une image, des antécédents]
Soit $f$ la fonction définie sur $\R$ par $f(x) = x^{2} - 3x$.

\textbf{Image de $4$.} On remplace : $f(4) = 4^{2} - 3 \times 4 = 16 - 12 = 4$.
L'image de $4$ est $4$.

\textbf{Antécédents de $-2$.} On résout l'équation $f(x) = -2$, c'est-à-dire
\[
  x^{2} - 3x = -2
  \quad\daccord\quad x^{2} - 3x + 2 = 0 .
\]
Le discriminant vaut $\Delta = 9 - 8 = 1 > 0$, d'où les deux solutions
$x_{1} = \dfrac{3-1}{2} = 1$ et $x_{2} = \dfrac{3+1}{2} = 2$. Le réel $-2$
admet donc exactement deux antécédents : $1$ et $2$.

\textbf{Antécédents de $-3$.} On résout $x^{2} - 3x + 3 = 0$, dont le
discriminant vaut $\Delta = 9 - 12 = -3 < 0$. Le réel $-3$ n'a aucun
antécédent par $f$.
\end{exemple}

\begin{erreurclassique}
« Calculer l'antécédent de $5$ » ne se fait pas en remplaçant $x$ par $5$ :
cela donnerait l'\emph{image} de $5$. Chercher un antécédent, c'est toujours
\emph{résoudre une équation}. Le réflexe : dès que l'énoncé dit
« antécédent », écrivez $f(x) = \ldots$ avant toute chose.
\end{erreurclassique}

\begin{entrainement}
\exo\label{t11s-c01-e01} Soit $f$ définie sur $\R$ par $f(x) = 2x^{2} - 5x + 1$. Calculer
$f(0)$, $f(-1)$, $f\!\left(\dfrac{1}{2}\right)$.

\exo\label{t11s-c01-e02} Soit $g$ définie sur $\R \setminus \{2\}$ par
$g(x) = \dfrac{3x+1}{x-2}$. Calculer $g(0)$, $g(3)$, $g(-1)$.

\exo\label{t11s-c01-e03} Avec $f$ de l’exercice~\ref{t11s-c01-e01} : déterminer les antécédents de $1$, puis
ceux de $-3$.

\exo\label{t11s-c01-e04} Avec $g$ de l’exercice~\ref{t11s-c01-e02} : déterminer l'antécédent éventuel de $3$.
Que constatez-vous, et comment l'expliquer ?

\exo\label{t11s-c01-e05} Soit $h$ définie par $h(x) = \sqrt{x+4}$. Déterminer les antécédents
de $0$, de $3$, puis de $-1$.

\exo\label{t11s-c01-e06} Une fonction $f$ vérifie $f(-2) = 7$. Traduire cette égalité par une
phrase contenant le mot « image », puis par une phrase contenant le mot
« antécédent ».
\end{entrainement}

% ===========================================================================
\section{L'ensemble de définition}

\begin{definition}[ensemble de définition]
Lorsqu'une fonction $f$ est donnée par une expression, son \textbf{ensemble de
définition}, noté $\Df_{f}$, est l'ensemble des réels $x$ pour lesquels cette
expression a un sens, c'est-à-dire pour lesquels le calcul de $f(x)$ est
possible dans $\R$.
\end{definition}

\begin{propriete}[les trois obstacles]
Pour les expressions polynomiales, rationnelles et à racines carrées,
les contraintes suivantes déterminent le domaine :
\begin{enumerate}[label=\textbf{\alph*)}, leftmargin=1.8em, itemsep=0.2em]
  \item $x$ annule un \textbf{dénominateur} : le quotient $\dfrac{A}{B}$
        n'existe que si $B \neq 0$ ;
  \item $x$ rend négatif ce qui se trouve \textbf{sous une racine carrée} :
        $\sqrt{A}$ n'existe que si $A \geq 0$ ;
  \item les deux à la fois, lorsqu'une racine figure au dénominateur :
        $\dfrac{1}{\sqrt{A}}$ n'existe que si $A > 0$, l'inégalité devenant
        stricte.
\end{enumerate}
Une fonction polynôme n'est jamais concernée : son ensemble de définition
est $\R$ tout entier.
\end{propriete}

\begin{methode}{déterminer un ensemble de définition}
\begin{enumerate}[label=\textbf{\arabic*.}, leftmargin=1.6em, itemsep=0.2em]
  \item Repérer dans l'expression tous les dénominateurs et toutes les
        racines carrées. Pour un polynôme, conclure : $\Df_f=\R$. Pour une fonction trigonométrique, vérifier aussi son domaine; la tangente exige $\cos x\ne0$.
  \item Écrire une condition par obstacle rencontré : une équation
        $B \neq 0$, une inéquation $A \geq 0$.
  \item Résoudre chaque condition séparément. Un tableau de signes est
        souvent le chemin le plus court pour une inéquation du second degré
        ou pour un quotient.
  \item Conserver les réels qui vérifient \emph{toutes} les conditions à la
        fois : l'ensemble de définition est leur intersection.
  \item Écrire la réponse sous forme d'intervalles ou de réunion
        d'intervalles, jamais sous forme de phrase.
\end{enumerate}
\end{methode}

\begin{exemple}[un quotient]
Soit $f(x) = \dfrac{x+5}{x^{2} - x - 6}$.

Le seul obstacle est le dénominateur : il faut $x^{2} - x - 6 \neq 0$. On
résout d'abord l'équation associée : $\Delta = 1 + 24 = 25$, donc
$x_{1} = \dfrac{1-5}{2} = -2$ et $x_{2} = \dfrac{1+5}{2} = 3$.

Ces deux réels, et eux seuls, sont à exclure :
\[
  \Df_{f} = \R \setminus \{-2\,;\,3\}
          = \intervalleo{-\infty}{-2} \cup \intervalleo{-2}{3}
            \cup \intervalleo{3}{+\infty}.
\]
\end{exemple}

\begin{exemple}[une racine, puis les deux ensemble]
Soit $g(x) = \dfrac{\sqrt{x+1}}{x-4}$.

Deux obstacles cette fois. Sous la racine : $x + 1 \geq 0$, soit
$x \geq -1$. Au dénominateur : $x - 4 \neq 0$, soit $x \neq 4$.

On garde les réels vérifiant les deux conditions :
\[
  \Df_{g} = \intervallefo{-1}{4} \cup \intervalleo{4}{+\infty}.
\]
Remarquez que $-1$ est conservé — la racine y vaut $0$, ce qui est
parfaitement licite — tandis que $4$ est exclu.
\end{exemple}

\begin{erreurclassique}
Écrire $\Df_{g} = \intervallefo{-1}{+\infty}$ en « oubliant » $4$ parce que
la racine, elle, existe en $4$. Chaque obstacle doit être traité pour
lui-même : on ne conclut qu'après avoir croisé toutes les conditions.
\end{erreurclassique}

\begin{entrainement}
\exo\label{t11s-c01-e07} Déterminer l'ensemble de définition de chacune des fonctions suivantes.
\begin{multicols}{2}
\begin{enumerate}[label=\textbf{\alph*)}, leftmargin=1.6em, itemsep=0.1em]
  \item $f(x) = 3x^{3} - x + 7$
  \item $f(x) = \dfrac{1}{x+5}$
  \item $f(x) = \dfrac{2x}{x^{2}-9}$
  \item $f(x) = \sqrt{2x-6}$
  \item $f(x) = \sqrt{5-x}$
  \item $f(x) = \dfrac{1}{\sqrt{x-1}}$
\end{enumerate}
\end{multicols}

\exo\label{t11s-c01-e08} Même question pour
$f(x) = \dfrac{x-1}{x^{2}+1}$ et $g(x) = \dfrac{x-1}{x^{2}-1}$.
Comparer les deux résultats et dire d'où vient la différence.

\exo\label{t11s-c01-e09} Déterminer l'ensemble de définition de
$f(x) = \sqrt{x^{2} - 5x + 4}$.

\exo\label{t11s-c01-e10} Déterminer l'ensemble de définition de
$f(x) = \dfrac{\sqrt{3-x}}{x^{2}-4}$.

\exo\label{t11s-c01-e11} Déterminer l'ensemble de définition de
$f(x) = \dfrac{1}{\abs{x} - 2}$.

\exo\label{t11s-c01-e12} On considère $f(x) = \sqrt{\dfrac{x+2}{x-3}}$. Dresser le tableau de
signes du quotient, puis en déduire $\Df_{f}$.
\end{entrainement}

% ===========================================================================
\section{Sens de variation}

\begin{definition}[fonction croissante, fonction décroissante]
Soit $f$ une fonction définie sur un intervalle $I$.
\begin{itemize}[leftmargin=1.6em, itemsep=0.2em]
  \item $f$ est \textbf{croissante} sur $I$ lorsque, pour tous réels $a$ et
        $b$ de $I$ : $a \leq b \implies f(a) \leq f(b)$.
  \item $f$ est \textbf{décroissante} sur $I$ lorsque, pour tous réels $a$ et
        $b$ de $I$ : $a \leq b \implies f(a) \geq f(b)$.
\end{itemize}
Pour la croissance stricte, $a<b$ implique $f(a)<f(b)$; pour la
décroissance stricte, $a<b$ implique $f(a)>f(b)$. Une fonction \textbf{monotone} sur $I$ est une fonction
qui y est croissante, ou qui y est décroissante.
\end{definition}

\begin{remarque}[une variation se dit toujours \emph{sur un intervalle}]
La fonction inverse est décroissante sur $\intervalleo{-\infty}{0}$ et
décroissante sur $\intervalleo{0}{+\infty}$, mais elle n'est pas décroissante
sur leur réunion. Prenons en effet $a = -1$ et $b = 1$ : on a bien
$a \leq b$, et pourtant $f(a) = -1$ tandis que $f(b) = 1$, de sorte que
$f(a) < f(b)$ — l'inégalité n'a pas été renversée. Dire « $f$ est
décroissante sur $\R^{*}$ » est donc faux. On écrit deux lignes, pas une.
\end{remarque}

\begin{theoreme}[multiplier par une constante]
Soit $f$ une fonction monotone sur un intervalle $I$ et $\lambda$ un réel non
nul. La fonction $\lambda f$ a le même sens de variation que $f$ sur $I$ si
$\lambda > 0$, et le sens contraire si $\lambda < 0$.
\end{theoreme}

\begin{demonstration}
Supposons $f$ croissante sur $I$ et $\lambda > 0$. Soient $a$ et $b$ deux
réels de $I$ tels que $a \leq b$. Par croissance de $f$, on a
$f(a) \leq f(b)$. Multiplier les deux membres d'une inégalité par un réel
strictement positif en conserve le sens, donc
$\lambda f(a) \leq \lambda f(b)$, c'est-à-dire
$(\lambda f)(a) \leq (\lambda f)(b)$ : la fonction $\lambda f$ est croissante
sur $I$.

Si $\lambda < 0$, la même multiplication renverse l'inégalité et donne
$\lambda f(a) \geq \lambda f(b)$ : la fonction $\lambda f$ est alors
décroissante sur $I$. Le cas où $f$ est décroissante se traite mot pour mot
de la même manière.
\end{demonstration}

\begin{visualisation}[le tableau de variation résume la courbe]
\begin{center}
\begin{tikzpicture}[scale=0.95]
  \begin{axis}[
      axis lines = middle, xlabel = {$x$}, ylabel = {$y$},
      xmin = -2.6, xmax = 3.6, ymin = -3.2, ymax = 3.6,
      xtick = {-2,-1,1,2,3}, ytick = {-3,-2,-1,1,2,3},
      width = 8.6cm, height = 6.2cm,
      tick label style = {font=\footnotesize},
      label style = {font=\footnotesize},
      axis line style = {bmtGris}, every tick/.style = {bmtGris},
    ]
    \addplot[bmtSarcelle, very thick, samples = 200, domain = -2.2:3.2]
      {0.4*x^3 - 1.2*x^2 + 1};
    \fill[bmtOcre] (axis cs:0,1) circle (2.2pt);
    \fill[bmtOcre] (axis cs:2,-0.6) circle (2.2pt);
    \node[font=\scriptsize, bmtOcre, anchor=south west] at (axis cs:0.1,1.05) {$1$};
    \node[font=\scriptsize, bmtOcre, anchor=north west] at (axis cs:2.1,-0.65) {$-0{,}6$};
  \end{axis}
\end{tikzpicture}
\end{center}

\vspace{0.3em}
\noindent
La courbe monte, puis descend, puis remonte. Les deux points marqués sont
les endroits où le sens change. Le tableau de variation n'est rien d'autre
que cette lecture, écrite :

\vspace{0.4em}
\begin{center}
\renewcommand{\arraystretch}{1.15}
\begin{tabular}{|c|ccccccc|}
\hline
$x$ & $-\infty$ & & $0$ & & $2$ & & $+\infty$ \\
\hline
      & & & $1$ & & & & \\[-0.35em]
$f$   & & $\nearrow$ & & $\searrow$ & & $\nearrow$ & \\[-0.35em]
      & & & & & $-0{,}6$ & & \\
\hline
\end{tabular}
\end{center}

\vspace{0.3em}
\noindent
Les flèches se lisent de gauche à droite, et chaque flèche vaut pour un
intervalle entier : ici $\intervalleof{-\infty}{0}$, puis
$\intervalle{0}{2}$, puis $\intervallefo{2}{+\infty}$.
\end{visualisation}

\begin{entrainement}
\exo\label{t11s-c01-e13} Donner, sans calcul, le sens de variation de $x \mapsto -3x + 7$ sur
$\R$, puis celui de $x \mapsto 5x^{2}$ sur $\intervallefo{0}{+\infty}$.

\exo\label{t11s-c01-e14} Soit $f(x) = x^{2}$ et $g(x) = -4x^{2}$. Dresser les tableaux de
variation de $f$ et de $g$ sur $\R$ et justifier le second à l'aide du
théorème de multiplication par une constante.

\exo\label{t11s-c01-e15} Soient $a$ et $b$ deux réels de $\intervallefo{0}{+\infty}$ tels que
$a \leq b$. Démontrer que $\sqrt{a} \leq \sqrt{b}$. Qu'en déduire pour la
fonction racine carrée ?

\exo\label{t11s-c01-e16} Démontrer, en revenant à la définition, que la fonction
$x \mapsto \dfrac{1}{x}$ est décroissante sur $\intervalleo{0}{+\infty}$.

\exo\label{t11s-c01-e17} Une fonction $f$ est croissante sur $\intervalle{-3}{5}$ et vérifie
$f(-3) = -2$ et $f(5) = 8$. Que peut-on affirmer de $f(0)$ ? Que ne peut-on
\emph{pas} affirmer ?
\end{entrainement}

% ===========================================================================
\section{Parité : quand la courbe se reflète}

Avant de parler de parité, il faut s'assurer que la question a un sens : si
$-x$ ne se trouve pas dans l'ensemble de définition alors que $x$ s'y trouve,
comparer $f(-x)$ et $f(x)$ est impossible.

\begin{definition}[ensemble symétrique par rapport à zéro]
Une partie $\Df$ de $\R$ est \textbf{symétrique par rapport à zéro} lorsque,
pour tout réel $x$, l'appartenance de $x$ à $\Df$ entraîne celle de $-x$ à
$\Df$.
\end{definition}

\begin{definition}[fonction paire, fonction impaire]
Soit $f$ une fonction dont l'ensemble de définition $\Df_{f}$ est symétrique
par rapport à zéro.
\begin{itemize}[leftmargin=1.6em, itemsep=0.2em]
  \item $f$ est \textbf{paire} lorsque $f(-x) = f(x)$ pour tout $x$ de
        $\Df_{f}$.
  \item $f$ est \textbf{impaire} lorsque $f(-x) = -f(x)$ pour tout $x$ de
        $\Df_{f}$.
\end{itemize}
\end{definition}

\begin{theoreme}[lecture géométrique de la parité]
Le plan est rapporté à un repère orthogonal $\repere$ et $\Cf_{f}$ désigne la
courbe représentative de $f$.
\begin{itemize}[leftmargin=1.6em, itemsep=0.2em]
  \item $f$ est paire \ssi\ $\Cf_{f}$ est symétrique par rapport à l'axe des
        ordonnées.
  \item $f$ est impaire \ssi\ $\Cf_{f}$ est symétrique par rapport à
        l'origine $O$ du repère.
\end{itemize}
\end{theoreme}

\begin{demonstration}[cas d'une fonction paire]
Le point $M$ de coordonnées $(x\,;\,y)$ a pour symétrique par rapport à l'axe
des ordonnées le point $M'$ de coordonnées $(-x\,;\,y)$.

Supposons $f$ paire et prenons un point $M(x\,;\,f(x))$ de $\Cf_{f}$. Son
symétrique est $M'(-x\,;\,f(x))$. Comme $\Df_{f}$ est symétrique par rapport
à zéro, le réel $-x$ appartient à $\Df_{f}$, et la parité donne
$f(-x) = f(x)$. L'ordonnée de $M'$ est donc bien $f(-x)$ : le point $M'$
appartient à $\Cf_{f}$. La courbe est donc globalement invariante par cette
symétrie.

Réciproquement, si $\Cf_{f}$ est symétrique par rapport à l'axe des
ordonnées, alors pour tout $x$ de $\Df_{f}$ le symétrique
$M'(-x\,;\,f(x))$ du point $M(x\,;\,f(x))$ appartient à $\Cf_{f}$ ; son
ordonnée vaut donc $f(-x)$, d'où $f(-x) = f(x)$.
\end{demonstration}

\begin{visualisation}[deux symétries, deux comportements]
\begin{center}
\begin{tikzpicture}[scale=0.9]
  \begin{axis}[
      name = gauche,
      axis lines = middle, xlabel = {$x$},
      xmin = -2.6, xmax = 2.6, ymin = -0.6, ymax = 4.2,
      xtick = {-2,-1,1,2}, ytick = {1,2,3,4},
      width = 6.6cm, height = 5.2cm,
      title = {\footnotesize\sffamily $f$ paire : axe $(Oy)$},
      tick label style = {font=\scriptsize},
      label style = {font=\scriptsize},
      axis line style = {bmtGris}, every tick/.style = {bmtGris},
    ]
    \addplot[bmtSarcelle, very thick, samples = 150, domain = -2.1:2.1]
      {0.85*x^2};
    \draw[bmtOcre, thick, dashed] (axis cs:-1.4,1.666) -- (axis cs:1.4,1.666);
    \fill[bmtOcre] (axis cs:1.4,1.666) circle (1.7pt);
    \fill[bmtOcre] (axis cs:-1.4,1.666) circle (1.7pt);
  \end{axis}
  \begin{axis}[
      name = droite, at = {(gauche.right of south east)}, anchor = left of south west,
      axis lines = middle, xlabel = {$x$},
      xmin = -2.6, xmax = 2.6, ymin = -3.4, ymax = 3.4,
      xtick = {-2,-1,1,2}, ytick = {-3,-2,-1,1,2,3},
      width = 6.6cm, height = 5.2cm,
      title = {\footnotesize\sffamily $f$ impaire : centre $O$},
      tick label style = {font=\scriptsize},
      label style = {font=\scriptsize},
      axis line style = {bmtGris}, every tick/.style = {bmtGris},
    ]
    \addplot[bmtSarcelle, very thick, samples = 150, domain = -2.1:2.1]
      {0.35*x^3};
    \draw[bmtOcre, thick, dashed] (axis cs:-1.6,-1.434) -- (axis cs:1.6,1.434);
    \fill[bmtOcre] (axis cs:1.6,1.434) circle (1.7pt);
    \fill[bmtOcre] (axis cs:-1.6,-1.434) circle (1.7pt);
  \end{axis}
\end{tikzpicture}
\end{center}

\vspace{0.2em}
\noindent
À gauche, les deux points marqués ont la même hauteur : on passe de l'un à
l'autre par un pliage le long de l'axe des ordonnées. À droite, le segment
qui joint les deux points marqués passe par l'origine, et celle-ci en est le
milieu : on passe de l'un à l'autre par un demi-tour autour de $O$.
\end{visualisation}

\begin{methode}{établir qu'une fonction est paire ou impaire}
\begin{enumerate}[label=\textbf{\arabic*.}, leftmargin=1.6em, itemsep=0.2em]
  \item Déterminer $\Df_{f}$ et vérifier qu'il est symétrique par rapport à
        zéro. Si ce n'est pas le cas, conclure immédiatement : la fonction
        n'est ni paire ni impaire, et il n'y a rien d'autre à écrire.
  \item Calculer $f(-x)$ en remplaçant partout $x$ par $-x$, puis simplifier
        \emph{complètement} l'expression obtenue.
  \item Comparer le résultat à $f(x)$ et à $-f(x)$ :
        \begin{itemize}[leftmargin=1.2em, itemsep=0.05em]
          \item si $f(-x) = f(x)$, la fonction est paire ;
          \item si $f(-x) = -f(x)$, la fonction est impaire ;
          \item si aucune des deux égalités n'est vraie pour tout $x$, elle
                n'est ni l'une ni l'autre — et un seul contre-exemple
                numérique pour chacune des deux égalités suffit alors à le prouver.
        \end{itemize}
\end{enumerate}
\end{methode}

\begin{exemple}[les trois issues possibles]
\textbf{a)} $f(x) = 3x^{4} - x^{2} + 5$, définie sur $\R$, qui est symétrique
par rapport à zéro. On calcule
\[
  f(-x) = 3(-x)^{4} - (-x)^{2} + 5 = 3x^{4} - x^{2} + 5 = f(x).
\]
La fonction $f$ est paire. C'était prévisible : tous les exposants de $x$
sont pairs, et la constante $5$ se lit comme $5x^{0}$.

\textbf{b)} $g(x) = \dfrac{x}{x^{2}+1}$, définie sur $\R$. On calcule
\[
  g(-x) = \frac{-x}{(-x)^{2}+1} = \frac{-x}{x^{2}+1}
        = -\,\frac{x}{x^{2}+1} = -g(x).
\]
La fonction $g$ est impaire.

\textbf{c)} $h(x) = x^{2} + x$, définie sur $\R$. On calcule
$h(-x) = x^{2} - x$. Pour conclure, un contre-exemple suffit : $h(1) = 2$ et
$h(-1) = 0$. Comme $h(-1) \neq h(1)$ la fonction n'est pas paire, et comme
$h(-1) \neq -h(1)$ elle n'est pas impaire non plus.
\end{exemple}

\begin{propriete}[axe et centre de symétrie en un point quelconque]
Soit $f$ une fonction et $\Cf_{f}$ sa courbe dans un repère orthogonal
$\repere$. Soient $a$ et $b$ deux réels. Le domaine doit être symétrique par rapport à $a$ : $a+h$ appartient au domaine si et seulement si $a-h$ y appartient. Sous cette condition,
\begin{itemize}[leftmargin=1.6em, itemsep=0.2em]
  \item La droite d'équation $x = a$ est un \textbf{axe de symétrie} de
        $\Cf_{f}$ \ssi, pour tout réel $h$ tel que $a+h$ et $a-h$
        appartiennent à $\Df_{f}$,
        \[ f(a+h) = f(a-h). \]
  \item Le point $\Omega(a\,;\,b)$ est un \textbf{centre de symétrie} de
        $\Cf_{f}$ \ssi, pour tout réel $h$ tel que $a+h$ et $a-h$
        appartiennent à $\Df_{f}$,
        \[ f(a+h) + f(a-h) = 2b. \]
\end{itemize}
\end{propriete}

\begin{demonstration}[cas du centre de symétrie]
Le symétrique du point $M(a+h\,;\,f(a+h))$ par rapport au point
$\Omega(a\,;\,b)$ est le point $M'$ tel que $\Omega$ soit le milieu de
$[MM']$. En écrivant les coordonnées du milieu, $M'$ a pour abscisse
$2a - (a+h) = a - h$ et pour ordonnée $2b - f(a+h)$.

Dire que $\Omega$ est centre de symétrie de $\Cf_{f}$, c'est dire que $M'$
appartient à $\Cf_{f}$, c'est-à-dire que son ordonnée est l'image de son
abscisse :
\[
  2b - f(a+h) = f(a-h),
  \quad\text{soit}\quad
  f(a+h) + f(a-h) = 2b .
\]
\end{demonstration}

\begin{remarque}[la parité est un cas particulier]
En prenant $a = 0$, la première égalité devient $f(h) = f(-h)$ : c'est la
parité. En prenant $a = 0$ et $b = 0$, la seconde devient
$f(h) + f(-h) = 0$ : c'est l'imparité. Il n'y a donc qu'un seul résultat à
retenir, la parité n'en étant que la version centrée à l'origine.
\end{remarque}

\begin{entrainement}
\exo\label{t11s-c01-e18} Étudier la parité de chacune des fonctions suivantes.
\begin{multicols}{2}
\begin{enumerate}[label=\textbf{\alph*)}, leftmargin=1.6em, itemsep=0.1em]
  \item $f(x) = x^{6} - 4x^{2}$
  \item $f(x) = x^{5} + 2x$
  \item $f(x) = \dfrac{1}{x^{2}+3}$
  \item $f(x) = \dfrac{x^{3}}{x^{2}-4}$
  \item $f(x) = \abs{x} - 1$
  \item $f(x) = \sqrt{x} + 1$
\end{enumerate}
\end{multicols}

\exo\label{t11s-c01-e19} Pourquoi la question de la parité ne se pose-t-elle pas pour
$f(x) = \dfrac{1}{x-2}$ ? Répondre en une phrase, en citant la condition qui
n'est pas remplie.

\exo\label{t11s-c01-e20} Soit $f(x) = x^{2} - 6x + 5$. Démontrer que la droite d'équation
$x = 3$ est un axe de symétrie de la courbe de $f$.

\exo\label{t11s-c01-e21} Soit $g(x) = \dfrac{2x-1}{x-1}$, définie sur $\R \setminus \{1\}$.
Démontrer que le point $\Omega(1\,;\,2)$ est un centre de symétrie de la
courbe de $g$.

\exo\label{t11s-c01-e22} Soit $f$ une fonction paire et $g$ une fonction impaire, toutes deux
définies sur $\R$. Que peut-on dire de la parité de $f + g$ ? de $f \times g$ ?
Justifier.

\exo\label{t11s-c01-e23} Démontrer qu'une fonction définie sur $\R$, à la fois paire et impaire,
est nécessairement la fonction nulle.
\end{entrainement}

% ===========================================================================
\section{Périodicité : quand la courbe se répète}

\begin{definition}[fonction périodique]
Soit $f$ une fonction définie sur $\Df_{f}$ et $T$ un réel strictement
positif. La fonction $f$ est \textbf{périodique de période $T$} lorsque, pour
tout $x$ de $\Df_{f}$, les réels $x+T$ et $x-T$ appartiennent à $\Df_{f}$ et
\[
  f(x+T) = f(x).
\]
Le plus petit réel strictement positif vérifiant cette égalité, lorsqu'il
existe, s'appelle la \textbf{période} de $f$.
\end{definition}

\begin{propriete}[répétition et réduction de l'étude]
Si $f$ est périodique de période $T$, alors pour tout entier relatif $k$ et
tout $x$ de $\Df_{f}$ :
\[
  f(x + kT) = f(x).
\]
Il suffit donc d'étudier $f$ sur un intervalle de longueur $T$ — par exemple
$\intervallefo{0}{T}$ — et de reproduire la portion de courbe obtenue par
translations successives de vecteur $T\veci$ pour obtenir la courbe entière.
\end{propriete}

\begin{demonstration}
L’égalité est vraie pour $k=0$. Si elle est vraie pour un entier $k\ge0$, alors
$f(x+(k+1)T)=f((x+kT)+T)=f(x+kT)=f(x)$; une récurrence donne tous les entiers naturels.
En appliquant la définition en $x-T$, on obtient $f(x-T)=f(x)$. Répéter ce déplacement donne les entiers négatifs. Le domaine reste invariant à chaque déplacement.
\end{demonstration}

\begin{exemple}[les fonctions du cercle trigonométrique]
Les fonctions $x \mapsto \cos x$ et $x \mapsto \sin x$ sont définies sur $\R$
et périodiques de période $2\pi$ : un tour complet du cercle
trigonométrique ramène au même point, donc aux mêmes coordonnées.

La fonction $f(x) = \cos(3x)$ est également périodique. Cherchons sa période.
Pour tout réel $x$,
\[
  f\!\left(x + \frac{2\pi}{3}\right)
  = \cos\!\left(3x + 2\pi\right) = \cos(3x) = f(x),
\]
donc $\dfrac{2\pi}{3}$ est une période de $f$.
\end{exemple}

\begin{visualisation}[un motif, puis des copies]
\begin{center}
\begin{tikzpicture}
  \begin{axis}[
      axis lines = middle, xlabel = {$x$},
      xmin = -7.2, xmax = 7.2, ymin = -1.8, ymax = 1.8,
      xtick = {-6.2832, -3.1416, 3.1416, 6.2832},
      xticklabels = {$-2\pi$, $-\pi$, $\pi$, $2\pi$},
      ytick = {-1,1},
      width = 12cm, height = 4.6cm,
      tick label style = {font=\footnotesize},
      label style = {font=\footnotesize},
      axis line style = {bmtGris}, every tick/.style = {bmtGris},
    ]
    \addplot[bmtGrisClair, very thick, samples = 250, domain = -7:7]
      {sin(deg(x))};
    \addplot[bmtSarcelle, very thick, samples = 150, domain = 0:6.2832]
      {sin(deg(x))};
    \draw[bmtOcre, thick, dashed] (axis cs:0,-1.5) -- (axis cs:0,1.5);
    \draw[bmtOcre, thick, dashed] (axis cs:6.2832,-1.5) -- (axis cs:6.2832,1.5);
  \end{axis}
\end{tikzpicture}
\end{center}

\vspace{0.2em}
\noindent
La portion foncée, tracée sur $\intervallefo{0}{2\pi}$, contient toute
l'information. Le reste de la courbe, en clair, s'en déduit par translations
de vecteur $2\pi\veci$ vers la droite et vers la gauche. C'est là tout
l'intérêt de repérer une période : l'étude se fait une fois, le tracé se
répète.
\end{visualisation}

\begin{methode}{réduire l'intervalle d'étude}
Lorsqu'une fonction est à la fois périodique et paire (ou impaire), les deux
propriétés se cumulent et l'intervalle d'étude se réduit deux fois.
\begin{enumerate}[label=\textbf{\arabic*.}, leftmargin=1.6em, itemsep=0.2em]
  \item La périodicité de période $T$ ramène l'étude à un intervalle de
        longueur $T$, centré si possible en zéro :
        $\intervalle{-\frac{T}{2}}{\frac{T}{2}}$.
  \item La parité, ou l'imparité, permet alors de n'étudier que la moitié
        droite : $\intervalle{0}{\frac{T}{2}}$.
  \item On trace la courbe sur ce dernier intervalle, on complète par la
        symétrie donnée par la parité, puis on reproduit le motif obtenu par
        translations de vecteur $T\veci$.
\end{enumerate}
\end{methode}

\begin{exemple}[cumuler les deux réductions]
Soit $f(x) = \cos(2x)$, définie sur $\R$.

\textbf{Périodicité.} Pour tout réel $x$, $f(x + \pi) = \cos(2x + 2\pi)
= \cos(2x) = f(x)$. La fonction est périodique de période $\pi$, et l'on
peut se limiter à $\intervalle{-\frac{\pi}{2}}{\frac{\pi}{2}}$.

\textbf{Parité.} Pour tout réel $x$, $f(-x) = \cos(-2x) = \cos(2x) = f(x)$,
car la fonction cosinus est paire. La fonction $f$ est donc paire, et l'on
peut se limiter à $\intervalle{0}{\frac{\pi}{2}}$.

D'un intervalle infini, on est passé à un intervalle de longueur
$\dfrac{\pi}{2}$.
\end{exemple}

\begin{entrainement}
\exo\label{t11s-c01-e24} Démontrer que la fonction $x \mapsto \sin(4x)$ est périodique et
déterminer sa période.

\exo\label{t11s-c01-e25} Même question pour $x \mapsto \cos\!\left(\dfrac{x}{2}\right)$.

\exo\label{t11s-c01-e26} Soit $f(x) = \sin x + \cos x$. Démontrer que $f$ est périodique de
période $2\pi$. Cette fonction est-elle paire ? impaire ?

\exo\label{t11s-c01-e27} Soit $f(x) = \cos^{2} x$. Démontrer que $\pi$ est une période de $f$,
puis étudier sa parité et en déduire un intervalle d'étude réduit par la parité et la périodicité.

\exo\label{t11s-c01-e28} Une fonction $f$ est périodique de période $3$ et vérifie $f(1) = 5$.
Calculer $f(4)$, $f(7)$ et $f(-2)$.

\exo\label{t11s-c01-e29} Soit $f$ définie sur $\R$, paire et périodique de période $T$.
Démontrer que la droite d'équation $x = \dfrac{T}{2}$ est un axe de symétrie
de la courbe de $f$.
\end{entrainement}

% ===========================================================================
% BANQUE D'EXERCICES DE FIN DE CHAPITRE
% À compléter après le dépouillement des sujets d'examen de fin de première
% (probatoire du Cameroun séries C et D, compositions d'Afrique francophone).
% Blocs attendus : \begin{annales} ... \end{annales} puis \begin{plusloin}.
% ===========================================================================

\begin{jecherche}[une courbe reconstituée]\phantomsection\label{t11s-c01-problem}
On considère une fonction $f$ définie sur $\R$, impaire et périodique de
période $4$. On sait de plus que, pour tout $x$ de $\intervalle{0}{2}$,
$f(x) = x(2-x)$.

\begin{enumerate}[label=\textbf{\arabic*.}, leftmargin=1.6em, itemsep=0.2em]
  \item Calculer $f(0)$, $f(1)$ et $f(2)$.
  \item Déterminer $f(-1)$, puis $f(3)$, puis $f(5)$, en précisant à chaque
        fois laquelle des deux hypothèses est utilisée.
  \item Tracer la courbe de $f$ sur $\intervalle{0}{2}$, puis la compléter
        sur $\intervalle{-4}{6}$.
  \item La droite d'équation $x = 1$ est-elle un axe de symétrie de cette
        courbe ? Justifier.
\end{enumerate}
\end{jecherche}

\begin{jemeteste}\phantomsection\label{t11s-c01-test}
\begin{enumerate}[label=\textbf{\arabic*.}, leftmargin=1.6em, itemsep=0.25em]
  \item Déterminer l'ensemble de définition de
        $f(x) = \dfrac{\sqrt{x+3}}{x^{2}-16}$.
  \item Soit $g(x) = 2x^{2} - 8x + 1$. Déterminer les antécédents de $1$ par
        $g$.
  \item Étudier la parité de $h(x) = \dfrac{x^{4}+1}{x^{2}}$.
  \item Démontrer que la droite d'équation $x = 2$ est un axe de symétrie de
        la courbe de $g$ définie à la question $2$.
  \item Soit $u(x) = \sin(3x)$. Déterminer une période de $u$, puis étudier
        sa parité, et en déduire un intervalle réduit par ces deux propriétés sur lequel il
        suffit d'étudier $u$.
\end{enumerate}
\end{jemeteste}

\begin{notepourlaclasse}
Une organisation possible consiste à répartir le travail en quatre séances : image et antécédent avec l'ensemble de définition, puis le sens de
variation, puis la parité, puis la périodicité avec la réduction de
l'intervalle d'étude.

Le point qui coûte le plus de temps chaque année n'est pas la parité mais la
\emph{vérification préalable de la symétrie de l'ensemble de définition}. Les
élèves calculent $f(-x)$ sans y penser et concluent sur une fonction dont le
domaine ne s'y prête pas. L’exercice~\ref{t11s-c01-e19} sur la parité est fait
pour cela : le faire traiter au tableau, pas en devoir.

Dans une classe à grand effectif, la démonstration du centre de symétrie peut
être menée collectivement sur une figure au tableau, le calcul des
coordonnées du symétrique étant la seule étape réellement nouvelle ; le cas
de l'axe de symétrie est alors laissé en exercice.
\end{notepourlaclasse}

\ifbmtprof\section*{Corrigés et accompagnement}\fi
\begin{corrige}[automatismes]\label{sol-t11s-c01-auto}
Dans l’ordre : $x\ge3$; $x=\pm2$; $7$ et $7$; $7$ et $5$; $x\ne-1/2$; strictement négatif; $1,1,-1$; $x^2,-x^3$; $x=\pm3$; aucune racine carrée réelle de $-4$.
\end{corrige}
\begin{corrige}[exercice~\ref{t11s-c01-e01}]\label{sol-t11s-c01-e01}$f(0)=1$, $f(-1)=8$, $f(1/2)=-1$.
\end{corrige}
\begin{corrige}[exercice~\ref{t11s-c01-e02}]\label{sol-t11s-c01-e02}$g(0)=-1/2$, $g(3)=10$, $g(-1)=2/3$.
\end{corrige}
\begin{corrige}[exercice~\ref{t11s-c01-e03}]\label{sol-t11s-c01-e03}$f(x)=1$ équivaut à $x(2x-5)=0$, donc $0$ ou $5/2$. Pour $-3$, $2x^2-5x+4=0$ a pour discriminant $-7$ : aucun antécédent réel.
\end{corrige}
\begin{corrige}[exercice~\ref{t11s-c01-e04}]\label{sol-t11s-c01-e04}$g(x)=3$ équivaut, pour $x\ne2$, à $3x+1=3x-6$, impossible. Le quotient s’écrit $3+7/(x-2)$ : la valeur $3$ n’est jamais atteinte.
\end{corrige}
\begin{corrige}[exercice~\ref{t11s-c01-e05}]\label{sol-t11s-c01-e05}Le domaine est $[-4;+\infty[$. Pour $0$, $x=-4$; pour $3$, $x+4=9$ donc $x=5$; pour $-1$, impossible car une racine carrée est positive ou nulle.
\end{corrige}
\begin{corrige}[exercice~\ref{t11s-c01-e06}]\label{sol-t11s-c01-e06}$7$ est l’image de $-2$; $-2$ est un antécédent de $7$.
\end{corrige}
\begin{corrige}[exercice~\ref{t11s-c01-e07}]\label{sol-t11s-c01-e07}Dans l’ordre : $\R$; $\R\setminus\{-5\}$; $\R\setminus\{-3;3\}$; $[3;+\infty[$; $]-\infty;5]$; $]1;+\infty[$. Le dernier dénominateur exige $x-1>0$.
\end{corrige}
\begin{corrige}[exercice~\ref{t11s-c01-e08}]\label{sol-t11s-c01-e08}$D_f=\R$ car $x^2+1>0$. $D_g=\R\setminus\{-1;1\}$ : annulation du dénominateur; une simplification ne restitue pas les points exclus.
\end{corrige}
\begin{corrige}[exercice~\ref{t11s-c01-e09}]\label{sol-t11s-c01-e09}$(x-1)(x-4)\ge0$, donc $]-\infty;1]\cup[4;+\infty[$.
\end{corrige}
\begin{corrige}[exercice~\ref{t11s-c01-e10}]\label{sol-t11s-c01-e10}$x\le3$ et $x\ne-2,2$. Domaine $]-\infty;-2[\cup]-2;2[\cup]2;3]$.
\end{corrige}
\begin{corrige}[exercice~\ref{t11s-c01-e11}]\label{sol-t11s-c01-e11}$|x|\ne2$, donc $\R\setminus\{-2;2\}$.
\end{corrige}
\begin{corrige}[exercice~\ref{t11s-c01-e12}]\label{sol-t11s-c01-e12}Le quotient est positif sur $]-\infty;-2[$ et $]3;+\infty[$, nul en $-2$, négatif sur $]-2;3[$, indéfini en $3$. Domaine $]-\infty;-2]\cup]3;+\infty[$.
\end{corrige}
\begin{corrige}[exercice~\ref{t11s-c01-e13}]\label{sol-t11s-c01-e13}La fonction affine de pente $-3$ est strictement décroissante sur $\R$. Le carré, puis sa multiplication par $5>0$, est strictement croissant sur $[0;+\infty[$.
\end{corrige}
\begin{corrige}[exercice~\ref{t11s-c01-e14}]\label{sol-t11s-c01-e14}$x^2$ décroît de $+\infty$ à $0$ sur $]-\infty;0]$, puis croît vers $+\infty$. Multiplication par $-4$ : croissance de $-\infty$ à $0$, puis décroissance vers $-\infty$.
\end{corrige}
\begin{corrige}[exercice~\ref{t11s-c01-e15}]\label{sol-t11s-c01-e15}Si $\sqrt a>\sqrt b\ge0$, le carré conserve l’ordre et donnerait $a>b$, contradiction. Donc $\sqrt a\le\sqrt b$; la racine est croissante sur $[0;+\infty[$ (strictement si $a<b$).
\end{corrige}
\begin{corrige}[exercice~\ref{t11s-c01-e16}]\label{sol-t11s-c01-e16}Pour $0<a<b$, $ab>0$ et $1/a-1/b=(b-a)/(ab)>0$. La fonction inverse est strictement décroissante sur $]0;+\infty[$.
\end{corrige}
\begin{corrige}[exercice~\ref{t11s-c01-e17}]\label{sol-t11s-c01-e17}$-2\le f(0)\le8$. On ne peut déterminer sa valeur, ni affirmer une croissance stricte ou une continuité à partir de ces seules données.
\end{corrige}
\begin{corrige}[exercice~\ref{t11s-c01-e18}]\label{sol-t11s-c01-e18}Dans l’ordre : paire, impaire, paire, impaire (domaine $\R\setminus\{-2;2\}$), paire, ni paire ni impaire (domaine $[0;+\infty[$ non symétrique).
\end{corrige}
\begin{corrige}[exercice~\ref{t11s-c01-e19}]\label{sol-t11s-c01-e19}Le domaine $\R\setminus\{2\}$ n’est pas symétrique par rapport à zéro : $-2$ y appartient mais $2$ non.
\end{corrige}
\begin{corrige}[exercice~\ref{t11s-c01-e20}]\label{sol-t11s-c01-e20}Domaine $\R$, symétrique autour de $3$. $f(3+h)=h^2-4=f(3-h)$; l’axe est $x=3$.
\end{corrige}
\begin{corrige}[exercice~\ref{t11s-c01-e21}]\label{sol-t11s-c01-e21}Domaine $\R\setminus\{1\}$, symétrique autour de $1$. $g(1+h)=2+1/h$ et $g(1-h)=2-1/h$, somme $4$. Centre $(1;2)$.
\end{corrige}
\begin{corrige}[exercice~\ref{t11s-c01-e22}]\label{sol-t11s-c01-e22}La somme n’est en général ni paire ni impaire : $f(x)=1$, $g(x)=x$ donne $1+x$. Elle est paire si $g=0$ et impaire si $f=0$. Le produit est impaire car $(fg)(-x)=f(x)(-g(x))=-(fg)(x)$.
\end{corrige}
\begin{corrige}[exercice~\ref{t11s-c01-e23}]\label{sol-t11s-c01-e23}Les deux identités donnent $f(x)=-f(x)$, donc $2f(x)=0$ et $f(x)=0$ pour tout réel.
\end{corrige}
\begin{corrige}[exercice~\ref{t11s-c01-e24}]\label{sol-t11s-c01-e24}$\sin(4(x+\pi/2))=\sin(4x+2\pi)=\sin(4x)$. Pour prouver que $\pi/2$ est la plus petite période positive, une période $T$ impose en $x=0$ puis en $x=\pi/8$ : $\sin4T=0$ et $\cos4T=1$, donc $4T\in2\pi\Z$.
\end{corrige}
\begin{corrige}[exercice~\ref{t11s-c01-e25}]\label{sol-t11s-c01-e25}Une période est $4\pi$. Si $T>0$ est une période, en $x=0$, $\cos(T/2)=1$, donc $T/2\in2\pi\Z$; la plus petite est $4\pi$.
\end{corrige}
\begin{corrige}[exercice~\ref{t11s-c01-e26}]\label{sol-t11s-c01-e26}Les deux termes ont période $2\pi$, donc leur somme aussi. $f(0)=1$ exclut l’imparité, et $f(\pi/2)=1\ne-1=f(-\pi/2)$ exclut la parité.
\end{corrige}
\begin{corrige}[exercice~\ref{t11s-c01-e27}]\label{sol-t11s-c01-e27}$\cos(x+\pi)=-\cos x$, donc le carré a période $\pi$; il est pair. On peut se limiter à $[0;\pi/2]$. En $x=0$, une période $T$ impose $\cos^2T=1$, donc $T\in\pi\Z$ : $\pi$ est la période minimale positive.
\end{corrige}
\begin{corrige}[exercice~\ref{t11s-c01-e28}]\label{sol-t11s-c01-e28}$4=1+3$, $7=1+2\times3$, $-2=1-3$; les trois valeurs sont $5$.
\end{corrige}
\begin{corrige}[exercice~\ref{t11s-c01-e29}]\label{sol-t11s-c01-e29}$f(T/2+h)=f(h-T/2)=f(T/2-h)$ par périodicité puis parité. Le domaine reste invariant sous $x\mapsto T-x$ : l’axe est $x=T/2$.
\end{corrige}
\begin{corrige}[problème « une courbe reconstituée »]\label{sol-t11s-c01-problem}
$f(0)=0$, $f(1)=1$, $f(2)=0$. Par imparité $f(-1)=-1$; par périodicité $f(3)=f(-1)=-1$ et $f(5)=f(1)=1$.
Sur $[-2;0]$, $f(x)=x(2+x)$ par imparité; sur $[0;2]$, $f(x)=x(2-x)$. Reproduire ce motif de longueur $4$ par translations de $4k$ ($k\in\Z$).
La symétrie du seul arc $[0;2]$ doit être vérifiée sur toute la courbe. Pour $h\in[-1;1]$, $f(1+h)=1-h^2=f(1-h)$. La symétrie s’étend par périodicité et imparité aux deux autres demi-périodes : pour $1\le h\le3$, poser $t=h-1\in[0;2]$, alors $f(1-h)=f(-t)=-t(2-t)$ et $f(1+h)=f(t+2)=f(t-2)=-t(2-t)$. Toute valeur de $h$ se ramène par multiples de $4$ à l’un de ces cas. L’axe $x=1$ est donc bien un axe de symétrie. Le domaine est $\R$.
\end{corrige}
\begin{corrige}[test de fin de chapitre]\label{sol-t11s-c01-test}
1. Domaine $[-3;4[\cup]4;+\infty[$, car $x\ge-3$ et $x\ne\pm4$ (seul $4$ appartient au premier domaine).
2. $2x^2-8x=0$, donc $x=0$ ou $4$.
3. Domaine $\R^*$, symétrique; $h(-x)=h(x)$, donc paire.
4. $g(2+t)=2t^2-7=g(2-t)$; domaine $\R$ : axe $x=2$.
5. $\sin(3x)$ est impaire et de période minimale $2\pi/3$; $[0;\pi/3]$ suffit, puis symétrie centrale et translations.
\end{corrige}

\clearpage\bmtchapitre{Limites et continuité}{Calculer une limite en un point ou à l’infini, distinguer les limites latérales et étudier la continuité.}{Analyse}
\label{t11s-c02}
\begin{revision}Savoir factoriser un polynôme, simplifier un quotient sur son domaine et lire les abscisses et les ordonnées. Revoir les ensembles de définition au chapitre~\ref{t11s-c01}.\end{revision}
\begin{automatismes}\exo\label{t11s-c02-auto}Factoriser $x^2-9$. Calculer $\dfrac{6}{3}$, $\dfrac{1}{10^3}$ et $\sqrt{16}$. Résoudre $x-2>0$.\end{automatismes}
\begin{corrige}[exercice~\ref{t11s-c02-auto}]\label{sol-t11s-c02-auto}$(x-3)(x+3)$; $2$; $0{,}001$; $4$; $x>2$.\end{corrige}

\section{Observer un rapprochement}
\begin{situation}[autour d’une valeur interdite]
\exo\label{t11s-c02-act}Pour $x\ne1$, calculer $\dfrac{x^2-1}{x-1}$ aux valeurs $0{,}9$, $0{,}99$, $1{,}01$ et $1{,}1$. Factoriser le numérateur. Vers quelle valeur les résultats semblent-ils se rapprocher ? Le quotient est-il défini en $1$ ?
\end{situation}
\begin{corrige}[exercice~\ref{t11s-c02-act}]\label{sol-t11s-c02-act}Les valeurs sont $1{,}9$, $1{,}99$, $2{,}01$ et $2{,}1$. Pour $x\ne1$, le quotient vaut $x+1$, donc tend vers $2$. Il reste indéfini en $1$.\end{corrige}
\begin{definition}[limite finie en un point]
Dire que $f(x)$ tend vers $\ell$ lorsque $x$ tend vers $a$, c’est dire que l’on peut rendre $f(x)$ aussi proche que l’on veut de $\ell$ en prenant $x$ suffisamment proche de $a$ dans le domaine, avec $x\ne a$. On note $\lim_{x\to a}f(x)=\ell$. On suppose que des points du domaine autres que $a$ peuvent s’approcher de $a$.
\end{definition}
\begin{remarque}La limite ne dépend pas de la valeur éventuelle $f(a)$. Les limites à gauche ($x<a$) et à droite ($x>a$) doivent coïncider pour qu’une limite bilatérale existe lorsque les deux côtés appartiennent au domaine.\end{remarque}
\begin{exemple}[limites latérales]
Pour $f(x)=\dfrac1{x-2}$, $\lim_{x\to2^-}f(x)=-\infty$ et $\lim_{x\to2^+}f(x)=+\infty$. Ces limites sont différentes : il n’existe pas de limite bilatérale en $2$.
\end{exemple}
\section{Limites à l’infini et opérations}
\begin{definition}[limites infinies]
$f(x)\to+\infty$ signifie que $f(x)$ dépasse tout seuil réel fixé lorsque $x$ est suffisamment proche du point considéré, ou suffisamment grand si $x\to+\infty$. La définition de $-\infty$ utilise des valeurs inférieures à tout seuil fixé.
\end{definition}
\begin{propriete}[règles admises]
Si $f\to\ell$ et $g\to m$ avec $\ell,m\in\R$, alors $f+g\to\ell+m$, $fg\to\ell m$, et $f/g\to\ell/m$ si $m\ne0$. Si $f\ge0$ et $f\to\ell\ge0$, alors $\sqrt f\to\sqrt\ell$. Les polynômes ont pour limite leur valeur en tout réel. À l’infini, $x^n$ ($n\in\N^*$) domine ses termes de degré inférieur. Pour un quotient de polynômes, factoriser les puissances dominantes.
\end{propriete}
\begin{avertissement}Les écritures $0/0$, $\infty/\infty$ et $\infty-\infty$ sont des formes indéterminées : elles ne donnent aucune réponse. $\infty$ n’est pas un nombre sur lequel on calcule comme dans $\R$.\end{avertissement}
\begin{methode}{lever une indétermination}
Au voisinage d’un réel, factoriser puis simplifier sans oublier la valeur exclue; avec des racines, multiplier par l’expression conjuguée. À l’infini, diviser numérateur et dénominateur par la plus grande puissance utile.
\end{methode}
\begin{exemple}
Pour $x\ne4$ et $x\ge0$,
\[
\frac{\sqrt x-2}{x-4}=\frac1{\sqrt x+2},\qquad
\lim_{x\to4}\frac{\sqrt x-2}{x-4}=\frac14.
\]
De même, $\dfrac{3x^2-x+1}{2x^2+5}\to\dfrac32$ quand $x\to\pm\infty$, car on divise par $x^2$.
\end{exemple}
\section{Continuité}
\begin{definition}Une fonction définie en $a$ est continue en $a$ lorsque $\lim_{x\to a}f(x)=f(a)$, la limite étant prise dans son domaine. Elle est continue sur un intervalle si elle est continue en chacun de ses points, avec une limite latérale aux extrémités incluses.\end{definition}
\begin{propriete}[résultats admis]Les fonctions polynômes sont continues sur $\R$. Les quotients de fonctions continues sont continus là où le dénominateur ne s’annule pas. La racine carrée est continue sur $[0;+\infty[$. Les fonctions $\sin$ et $\cos$ sont continues sur $\R$, et $\tan$ sur chacun des intervalles de son domaine.\end{propriete}
\begin{theoreme}[valeurs intermédiaires, admis]Si $f$ est continue sur $[a;b]$, tout réel compris entre $f(a)$ et $f(b)$ est atteint au moins une fois. Si $f$ est de plus strictement monotone, cet antécédent est unique.\end{theoreme}
La continuité assure l’existence, la stricte monotonie assure l’unicité. La seule connaissance des valeurs aux bornes ne suffit pas sans continuité.

% ENRICHISSEMENT C02

\subsection*{Choisir une technique et rédiger le raisonnement}
\begin{methode}{décider avant de calculer}
\begin{enumerate}
\item Écrire le domaine et le point approché. Une limite en une borne peut n’être accessible que d’un côté.
\item Si les opérations portent sur des limites finies et un dénominateur non nul, substituer puis citer la règle utilisée.
\item Pour un dénominateur tendant vers zéro, déterminer son signe de chaque côté; une taille ne suffit pas à déterminer le signe d’une limite infinie.
\item Pour une forme indéterminée, transformer l’expression sur un voisinage où elle est définie. Conclure seulement après cette transformation.
\end{enumerate}
\end{methode}
\begin{exemple}[un quotient dont le signe change]
Pour $r(x)=(x+1)/(x-2)$, le numérateur tend vers $3>0$ en $2$. À gauche, $x-2<0$ et tend vers $0$ : $r(x)\to-\infty$. À droite, le dénominateur est positif : $r(x)\to+\infty$. Écrire seulement « le dénominateur tend vers zéro » ne permettrait pas de choisir entre ces deux réponses.
\end{exemple}
\begin{exemple}[dominer les termes sans oublier le signe]
Pour $x\ne0$,
\[
\frac{2x^3-x+1}{x^2+4}
=x\frac{2-1/x^2+1/x^3}{1+4/x^2}.
\]
Le quotient de droite tend vers $2>0$ aux deux infinis; le facteur $x$ impose $+\infty$ à droite et $-\infty$ à gauche. Diviser toujours par le degré du numérateur aurait aussi fonctionné, mais aurait laissé un dénominateur tendant vers zéro à analyser.
\end{exemple}
\begin{methode}{prouver puis encadrer un zéro}
Pour $F(x)=x^3+x-1$, établir la continuité sur l’intervalle choisi et des signes opposés aux bornes. Pour l’unicité, si $a<b$,
\[
F(b)-F(a)=(b-a)(a^2+ab+b^2+1)>0,
\]
car $a^2+ab+b^2=(a+b/2)^2+3b^2/4\ge0$. Une fois l’unique zéro situé, tester le milieu d’un intervalle et conserver la moitié où les signes restent opposés. Cette dichotomie produit un encadrement, pas une valeur exacte par calculatrice.
\end{methode}
\begin{avertissement}[continuité et valeur au point]
Une limite finie n’impose pas que la fonction soit définie au point. Si elle l’est, sa valeur peut différer de la limite. Pour rendre une fonction continue par prolongement, la valeur ajoutée doit être égale à cette limite; deux limites latérales différentes empêchent tout tel prolongement.
\end{avertissement}

\section{Exercices et problèmes}
\exo[Limites directes]\label{t11s-c02-e01}
Calculer $\lim_{x\to2}(x^2-3x+4)$ et $\lim_{x\to-1}\dfrac{2x+5}{x^2+2}$.
\begin{corrige}[exercice~\ref{t11s-c02-e01}]\label{sol-t11s-c02-e01}
Continuité des expressions et dénominateur $3\ne0$ : les deux limites valent respectivement $2$ et $1$.
\end{corrige}
\exo[Factoriser]\label{t11s-c02-e02}
Calculer $\lim_{x\to3}\dfrac{x^2-9}{x-3}$ et $\lim_{x\to1}\dfrac{x^2-3x+2}{x-1}$. Préciser les valeurs interdites.
\begin{corrige}[exercice~\ref{t11s-c02-e02}]\label{sol-t11s-c02-e02}
Les quotients valent $x+3$ pour $x\ne3$, et $x-2$ pour $x\ne1$. Limites $6$ et $-1$.
\end{corrige}
\exo[Racines]\label{t11s-c02-e03}
Calculer $\lim_{x\to0}\dfrac{\sqrt{x+4}-2}{x}$.
\begin{corrige}[exercice~\ref{t11s-c02-e03}]\label{sol-t11s-c02-e03}
Domaine $[-4;+\infty[\setminus\{0\}$. Le conjugué donne $1/(\sqrt{x+4}+2)$, donc la limite est $1/4$.
\end{corrige}
\exo[Aux deux infinis]\label{t11s-c02-e04}
Calculer aux deux infinis les limites de $-2x^3+5x$ et de $\dfrac{4x^2+1}{x^2-3x}$.
\begin{corrige}[exercice~\ref{t11s-c02-e04}]\label{sol-t11s-c02-e04}
Le premier vaut $x^3(-2+5/x^2)$ : limite $-\infty$ en $+\infty$, $+\infty$ en $-\infty$. Le quotient tend vers $4$ aux deux infinis.
\end{corrige}
\exo[Raccorder]\label{t11s-c02-e05}
Pour $x<1$, $f(x)=x+2$; pour $x\ge1$, $f(x)=mx$. Déterminer $m$ pour que $f$ soit continue sur $\R$.
\begin{corrige}[exercice~\ref{t11s-c02-e05}]\label{sol-t11s-c02-e05}
Chaque branche est continue. En $1$, la limite gauche est $3$, la limite droite et la valeur sont $m$. La condition nécessaire et suffisante est $m=3$.
\end{corrige}
\exo[Un zéro]\label{t11s-c02-e06}
Démontrer que $x^3+x-1=0$ possède au moins une solution dans $]0;1[$. Peut-on prouver l’unicité sans dérivée ?
\begin{corrige}[exercice~\ref{t11s-c02-e06}]\label{sol-t11s-c02-e06}
La fonction est continue; ses valeurs en $0$ et $1$ sont $-1$ et $1$. Le théorème des valeurs intermédiaires donne une solution intérieure. Pour $a<b$, $b^3-a^3>0$ et $b-a>0$, donc $f(b)>f(a)$ : unicité.
\end{corrige}
\exo[Se méfier des tables]\label{t11s-c02-e07}
Une table donne $f(0{,}9)=1{,}9$ et $f(1{,}1)=2{,}1$. Cela prouve-t-il que $\lim_{x\to1}f(x)=2$ ? Justifier par deux fonctions.
\begin{corrige}[exercice~\ref{t11s-c02-e07}]\label{sol-t11s-c02-e07}
Non. $f(x)=x+1$ satisfait la table et a pour limite $2$. $g(x)=x+1+100(x-0{,}9)(x-1{,}1)$ satisfait la même table, mais sa limite en $1$ vaut $1$. Deux valeurs ne prouvent pas une limite.
\end{corrige}

\exo[Encadrer sans confondre exact et approché]\label{t11s-c02-p01}
On pose $F(x)=x^3+x-1$ sur $\R$.
\begin{enumerate}
\item Prouver que $F$ a un unique zéro $\alpha$ dans $]0;1[$, sans utiliser de dérivée.
\item Calculer exactement $F(0{,}68)$ et $F(0{,}69)$, puis encadrer $\alpha$.
\item Donner une valeur approchée avec une erreur au plus $0{,}005$. Expliquer pourquoi cela ne fournit pas une solution exacte.
\end{enumerate}
\begin{corrige}[exercice~\ref{t11s-c02-p01}]\label{sol-t11s-c02-p01}
Le polynôme est continu, $F(0)=-1$, $F(1)=1$; le TVI donne un zéro intérieur. La différence factorisée ci-dessus est positive pour $a<b$, donc $F$ est strictement croissante et ce zéro est unique. On calcule $F(0{,}68)=-0{,}005568<0$ et $F(0{,}69)=0{,}018509>0$. Donc $0{,}68<\alpha<0{,}69$. Le milieu $0{,}685$ a une erreur strictement inférieure à $0{,}005$. Il n’est pas présenté comme racine exacte : l’encadrement garantit une distance, pas une égalité.
\end{corrige}
\exo[Peut-on réparer un saut par une seule valeur ?]\label{t11s-c02-p02}
Pour $x\ne1$, on définit $s(x)=\dfrac{x-1}{|x-1|}$. Calculer ses limites latérales en $1$. Existe-t-il un réel $k$ tel que la fonction obtenue en posant $s(1)=k$ soit continue sur $\R$ ? Justifier.
\begin{corrige}[exercice~\ref{t11s-c02-p02}]\label{sol-t11s-c02-p02}
Pour $x<1$, $|x-1|=1-x$ et $s(x)=-1$; pour $x>1$, $s(x)=1$. Les limites latérales sont $-1$ et $1$. La continuité demanderait simultanément $k=-1$ et $k=1$, impossible. Changer une valeur au point ne modifie pas ces deux limites.
\end{corrige}

\begin{jemeteste}Je sais distinguer valeur et limite, traiter un quotient $0/0$, vérifier un raccord et séparer existence et unicité d’une solution.\end{jemeteste}
\begin{notepourlaclasse}Introduire les limites par des tableaux, puis justifier les conclusions algébriquement. Les règles générales et le théorème des valeurs intermédiaires sont admis; ne pas présenter les tableaux numériques comme des preuves.\end{notepourlaclasse}

\clearpage\bmtchapitre{Branches infinies et asymptotes}{Reconnaître une asymptote verticale, horizontale ou oblique et étudier la position relative d’une courbe et de sa droite asymptote.}{Analyse}
\label{t11s-c03}
\begin{revision}Limites latérales et à l’infini; signe d’un quotient; équation réduite d’une droite.\end{revision}
\begin{automatismes}\exo\label{t11s-c03-auto}Calculer $\dfrac{x^2-1}{x-1}$ pour $x\ne1$. Donner le signe de $1/(x-2)$ sur chacun des intervalles de son domaine.\end{automatismes}
\begin{corrige}[exercice~\ref{t11s-c03-auto}]\label{sol-t11s-c03-auto}Le quotient vaut $x+1$. Le second est négatif pour $x<2$ et positif pour $x>2$.\end{corrige}

\section{Une courbe qui se rapproche d’une droite}
\begin{definition}[asymptotes verticales et horizontales]
La droite $x=a$ est une asymptote verticale si au moins une limite latérale de $f$ en $a$ vaut $+\infty$ ou $-\infty$. La droite $y=\ell$ est une asymptote horizontale en $+\infty$ (respectivement $-\infty$) si $\lim_{x\to+\infty}f(x)=\ell$ (respectivement en $-\infty$).
\end{definition}
\begin{exemple}Pour $f(x)=2+\dfrac1{x-1}$, la droite $x=1$ est asymptote verticale et $y=2$ est asymptote horizontale aux deux infinis. $f(x)-2=1/(x-1)$ donne une courbe sous $y=2$ à gauche de $1$, au-dessus à droite.\end{exemple}
\begin{visualisation}[l’écart vertical tend vers zéro]
\begin{center}\begin{tikzpicture}\begin{axis}[width=10cm,height=6cm,axis lines=middle,xmin=-4,xmax=6,ymin=-2,ymax=6,xlabel={$x$},ylabel={$y$},samples=100]
\addplot[bmtSarcelle,thick,domain=-4:0.75]{2+1/(x-1)};
\addplot[bmtSarcelle,thick,domain=1.25:6]{2+1/(x-1)};
\addplot[bmtOcre,dashed,domain=-4:6]{2};
\draw[bmtOcre,dashed] (axis cs:1,-2)--(axis cs:1,6);
\end{axis}\end{tikzpicture}\end{center}
Les deux branches sont tracées séparément : relier leurs extrémités créerait un segment qui n’appartient pas à la courbe.
\end{visualisation}
\section{Asymptote oblique}
\begin{definition}La droite $y=ax+b$, avec $a\ne0$, est une asymptote oblique en l’infini considéré si $f(x)-(ax+b)$ tend vers $0$.\end{definition}
\begin{propriete}[méthode des deux limites]
Si $\lim f(x)/x=a\in\R$ et $\lim(f(x)-ax)=b\in\R$, alors $y=ax+b$ est une asymptote; elle est oblique si $a\ne0$, horizontale si $a=0$.
\end{propriete}
\begin{demonstration}La seconde limite signifie exactement $f(x)-(ax+b)\to0$. Réciproquement, cette relation donne $f(x)/x=a+b/x+o(1)/x\to a$. La notation $o(1)$ signifie ici une quantité tendant vers zéro.\end{demonstration}
\begin{methode}{utiliser la division euclidienne}
Pour une fraction rationnelle dont le numérateur a un degré de plus que le dénominateur, effectuer la division. Si $f(x)=ax+b+r(x)$ avec $r(x)\to0$, la droite cherchée est $y=ax+b$. Étudier le signe de $r$ pour situer la courbe.
\end{methode}
\begin{exemple}
\[
\frac{x^2+1}{x-1}=x+1+\frac2{x-1}\qquad(x\ne1).
\]
L’asymptote oblique aux deux infinis est $y=x+1$. La courbe est en dessous pour $x<1$, au-dessus pour $x>1$. $x=1$ est également une asymptote verticale.
\end{exemple}
\section{Toutes les branches n’ont pas d’asymptote}
\begin{remarque}Si $f(x)\to\pm\infty$ et $|f(x)/x|\to+\infty$, on parle d’une branche parabolique de direction l’axe des ordonnées. Si $f(x)\to\pm\infty$ et $f(x)/x\to0$, la direction est l’axe des abscisses. Si $f(x)/x\to a\ne0$ mais $f(x)-ax\to\pm\infty$, la branche a pour direction $y=ax$, sans asymptote oblique. Ce vocabulaire décrit une direction, pas nécessairement une parabole.\end{remarque}
Pour $x^2$ en $+\infty$, $f(x)/x=x\to+\infty$. Pour $\sqrt x$, $f(x)/x=1/\sqrt x\to0$ alors que $f(x)\to+\infty$ : aucune asymptote horizontale.

\section{Exercices et problèmes}
\exo[Horizontale]\label{t11s-c03-e01}
Pour $f(x)=\dfrac{3x+2}{x-2}$, déterminer les asymptotes et la position relative à l’asymptote horizontale.
\begin{corrige}[exercice~\ref{t11s-c03-e01}]\label{sol-t11s-c03-e01}
$f(x)=3+8/(x-2)$; asymptotes $x=2$ et $y=3$. Courbe en dessous pour $x<2$ et au-dessus pour $x>2$.
\end{corrige}
\exo[Division]\label{t11s-c03-e02}
Pour $g(x)=\dfrac{2x^2-3x+4}{x-1}$, effectuer la division puis déterminer l’asymptote oblique.
\begin{corrige}[exercice~\ref{t11s-c03-e02}]\label{sol-t11s-c03-e02}
$g(x)=2x-1+3/(x-1)$. L’écart tend vers zéro aux deux infinis : asymptote $y=2x-1$; verticale $x=1$.
\end{corrige}
\exo[Trou ou asymptote]\label{t11s-c03-e03}
La droite $x=1$ est-elle asymptote de $h(x)=\dfrac{x^2-1}{x-1}$ ? Expliquer.
\begin{corrige}[exercice~\ref{t11s-c03-e03}]\label{sol-t11s-c03-e03}
Non : $h(x)=x+1$ pour $x\ne1$, et la limite en $1$ vaut $2$. La courbe est une droite privée d’un point.
\end{corrige}
\exo[Une racine]\label{t11s-c03-e04}
Calculer $\lim_{x\to+\infty}(\sqrt{x^2+1}-x)$ puis donner une asymptote. Recommencer en $-\infty$.
\begin{corrige}[exercice~\ref{t11s-c03-e04}]\label{sol-t11s-c03-e04}
En $+\infty$, le conjugué donne $1/(\sqrt{x^2+1}+x)\to0$, donc $y=x$. En $-\infty$, $\sqrt{x^2+1}+x=1/(\sqrt{x^2+1}-x)\to0$, donc $y=-x$. Ne pas remplacer $\sqrt{x^2}$ par $x$ si $x<0$.
\end{corrige}
\exo[Lire un écart]\label{t11s-c03-e05}
On sait $f(x)=2x-3+\dfrac{x}{x^2+1}$. Donner l’asymptote aux deux infinis et tous ses points d’intersection avec la courbe.
\begin{corrige}[exercice~\ref{t11s-c03-e05}]\label{sol-t11s-c03-e05}
L’écart tend vers zéro : $y=2x-3$. Il s’annule uniquement en $x=0$, d’où le point $(0;-3)$. Une courbe peut couper son asymptote.
\end{corrige}
\exo[Classer]\label{t11s-c03-e06}
Décrire les branches en $+\infty$ de $x^2$, $\sqrt{x}$ et $x+\sqrt{x}$, en utilisant les quotients par $x$.
\begin{corrige}[exercice~\ref{t11s-c03-e06}]\label{sol-t11s-c03-e06}
Quotients : $x\to+\infty$ (direction verticale), $1/\sqrt x\to0$ (direction horizontale), $1+1/\sqrt x\to1$ mais écart à $y=x$ infini (direction $y=x$, sans asymptote).
\end{corrige}
\begin{jemeteste}Je calcule une limite avant de nommer une asymptote; je ne confonds pas direction et droite asymptote; le signe de l’écart donne la position relative.\end{jemeteste}
\begin{notepourlaclasse}Faire comparer le trou de la fonction simplifiable et la branche infinie. La recherche par division est prioritaire; réserver le vocabulaire des branches sans asymptote à une lecture guidée.\end{notepourlaclasse}

\clearpage\bmtchapitre{Dérivation et tangentes}{Calculer un nombre dérivé, dériver les fonctions usuelles et utiliser le signe de la dérivée pour déterminer les variations.}{Analyse}
\label{t11s-c04}
\begin{revision}Limites, continuité, développement et factorisation; pente d’une droite.\end{revision}
\begin{automatismes}\exo\label{t11s-c04-auto}Développer $(a+h)^2-a^2$. Calculer la pente de la droite passant par $(1;2)$ et $(3;6)$.\end{automatismes}
\begin{corrige}[exercice~\ref{t11s-c04-auto}]\label{sol-t11s-c04-auto}$2ah+h^2$; pente $(6-2)/(3-1)=2$.\end{corrige}

\section{Du taux de variation au nombre dérivé}
\begin{situation}
\exo\label{t11s-c04-act}Pour $f(x)=x^2$, calculer la pente de la sécante entre les points d’abscisses $2$ et $2+h$ ($h\ne0$). Que devient cette pente quand $h\to0$ ?
\end{situation}
\begin{corrige}[exercice~\ref{t11s-c04-act}]\label{sol-t11s-c04-act}Le taux vaut $((2+h)^2-4)/h=4+h$, et tend vers $4$.\end{corrige}
\begin{definition}Si $a$ est intérieur au domaine et si $\lim_{h\to0}\dfrac{f(a+h)-f(a)}h$ existe et est un réel, $f$ est dérivable en $a$ et cette limite est $f'(a)$. La fonction dérivée associe à chaque point où ce nombre existe sa valeur $f'(a)$.\end{definition}
\begin{propriete}[tangente]Si $f$ est dérivable en $a$, la tangente au point $(a;f(a))$ a pour équation $y=f'(a)(x-a)+f(a)$.\end{propriete}
\begin{demonstration}La droite passe par $(a;f(a))$ et a pour coefficient directeur la limite des pentes des sécantes, $f'(a)$. La forme point-pente donne l’équation.\end{demonstration}
\begin{propriete}La dérivabilité en un point entraîne la continuité en ce point. La réciproque est fausse.\end{propriete}
\begin{demonstration}Pour $h\ne0$, $f(a+h)-f(a)=h\dfrac{f(a+h)-f(a)}h$. Le second facteur tend vers un réel, donc le produit tend vers $0$. Ainsi $f(a+h)\to f(a)$. Pour $|x|$ en $0$, les taux sont $-1$ à gauche et $1$ à droite : continuité, mais pas dérivabilité.\end{demonstration}
\section{Dérivées usuelles et opérations}
\begin{center}\begin{tabular}{lll}\toprule
Fonction & Dérivée & Domaine de dérivation\\\midrule
$c$ & $0$ & $\R$\\
$x^n$, $n\in\N^*$ & $nx^{n-1}$ & $\R$\\
$1/x$ & $-1/x^2$ & $\R^*$\\
$\sqrt x$ & $1/(2\sqrt x)$ & $]0;+\infty[$\\
$\sin x$ & $\cos x$ & $\R$\\
$\cos x$ & $-\sin x$ & $\R$\\
$\tan x$ & $1/\cos^2x$ & $\cos x\ne0$\\\bottomrule
\end{tabular}\end{center}
Les formules trigonométriques supposent des angles en radians. La fonction racine carrée est continue en $0$, mais n’y admet pas de nombre dérivé fini : son taux vaut $1/\sqrt h$ pour $h>0$.
\begin{propriete}[opérations]
Pour $u,v$ dérivables et $\lambda\in\R$,
\[
\begin{aligned}
(u+v)'&=u'+v',& (\lambda u)'&=\lambda u',\\
(uv)'&=u'v+uv',& \left(\frac uv\right)'&=\frac{u'v-uv'}{v^2}\quad(v\ne0).
\end{aligned}
\]
Si $u$ est dérivable et $F$ dérivable aux valeurs de $u$, $(F\circ u)'=(F'\circ u)u'$. Les formules générales sont admises, sauf la justification du produit ci-dessous.
\end{propriete}
\begin{demonstration}[produit]
Ajouter et retrancher $u(a+h)v(a)$ dans la différence du produit donne
\[
\frac{u(a+h)v(a+h)-u(a)v(a)}h
=u(a+h)\frac{v(a+h)-v(a)}h
+v(a)\frac{u(a+h)-u(a)}h.
\]
La continuité de $u$ et les limites des taux donnent $u(a)v'(a)+v(a)u'(a)$.
\end{demonstration}
\begin{exemple}Pour $f(x)=\dfrac{x^2+1}{x-1}$, $x\ne1$,
\[
f'(x)=\frac{2x(x-1)-(x^2+1)}{(x-1)^2}
=\frac{x^2-2x-1}{(x-1)^2}.
\]
Pour $g(x)=\sqrt{3x+2}$, $g'(x)=3/(2\sqrt{3x+2})$ sur $]-2/3;+\infty[$.
\end{exemple}
\section{Variations et extrema}
\begin{theoreme}[admis]Sur un intervalle, une fonction dérivable dont la dérivée est positive ou nulle est croissante; si elle est négative ou nulle, elle est décroissante. Une dérivée strictement positive donne une croissance stricte. Si elle est nulle partout sur l’intervalle, la fonction est constante.\end{theoreme}
\begin{remarque}Un zéro isolé de $f'$ n’impose pas un extremum : $x^3$ a une dérivée nulle en $0$ et reste croissante. Un changement de signe de $+$ vers $-$ donne un maximum local; de $-$ vers $+$, un minimum local. Un extremum au bord d’un intervalle ne nécessite pas une dérivée nulle.\end{remarque}
\begin{methode}{dresser un tableau de variation}Déterminer le domaine; calculer et factoriser la dérivée; étudier son signe sur chaque intervalle; calculer les valeurs aux points critiques et les limites aux bornes; tracer les flèches de variation.\end{methode}

% ENRICHISSEMENT C04

\subsection*{Organiser un calcul puis interpréter le résultat}
\begin{methode}{dériver une expression composée}
Repérer la fonction extérieure et l’expression intérieure. Pour $(u(x))^n$, dériver la puissance puis multiplier par $u'(x)$; pour $\sqrt{u(x)}$, exiger $u(x)>0$ avant d’utiliser la formule; pour un quotient, écrire le numérateur de la dérivée avant de le développer. Toujours annoncer le domaine de dérivation.
\end{methode}
\begin{exemple}[deux étapes qui ne doivent pas être confondues]
Pour $F(x)=\sqrt{2x+3}$, le domaine de définition est $[-3/2;+\infty[$. La formule de dérivation donne, pour $x>-3/2$,
\[
F'(x)=\frac{2}{2\sqrt{2x+3}}=\frac1{\sqrt{2x+3}}.
\]
L’extrémité est dans le domaine de définition mais pas dans ce domaine de dérivation. Pour $G(x)=(x^2+1)^3$, les deux fonctions sont dérivables sur $\R$ et $G'(x)=3(x^2+1)^2\times2x$.
\end{exemple}
\begin{methode}{passer du signe à un optimum}
Après avoir calculé $f'$, déterminer son signe sur le domaine physique du problème. Un changement de $+$ à $-$ prouve un maximum local; pour un maximum global, comparer tous les intervalles et les éventuelles bornes incluses. Donner la valeur maximale, l’endroit où elle est atteinte et les unités. Un point où la dérivée est nulle est seulement un candidat.
\end{methode}
\begin{exemple}[un optimum et une autre preuve]
Pour un rectangle de périmètre $20$ m, si $x$ est un côté en mètres, l’autre vaut $10-x$ et $0<x<10$. L’aire $A(x)=10x-x^2$ a une dérivée $10-2x$, qui change de signe en $5$. Le maximum vaut $25\ \mathrm{m}^2$, pour un carré. La forme $A(x)=25-(x-5)^2$ fournit une seconde preuve globale et précise l’égalité.
\end{exemple}
\begin{remarque}[une tangente est une droite]
Pour rédiger sa recherche, calculer séparément les coordonnées du point et la pente, puis utiliser la forme point-pente. Comparer la courbe à cette droite signifie étudier le signe de la différence des ordonnées; regarder seulement le dessin ne prouve pas ce signe.
\end{remarque}

\section{Exercices et problèmes}
\exo[Par la définition]\label{t11s-c04-e01}
À partir du taux de variation, montrer que la dérivée de $x\mapsto3x^2-2x+1$ en $a$ vaut $6a-2$.
\begin{corrige}[exercice~\ref{t11s-c04-e01}]\label{sol-t11s-c04-e01}
La différence vaut $6ah+3h^2-2h$. Le taux est $6a+3h-2$; sa limite est $6a-2$.
\end{corrige}
\exo[Dériver]\label{t11s-c04-e02}
Dériver $x^3-4x+2$, $(2x-1)^4$, $(x+1)/(x-2)$ et $\sin(3x)$. Préciser les domaines.
\begin{corrige}[exercice~\ref{t11s-c04-e02}]\label{sol-t11s-c04-e02}
$3x^2-4$ sur $\R$; $8(2x-1)^3$ sur $\R$; $-3/(x-2)^2$ sur $\R\setminus\{2\}$; $3\cos(3x)$ sur $\R$.
\end{corrige}
\exo[Tangente]\label{t11s-c04-e03}
Pour $f(x)=x^2-3x$, donner l’équation de la tangente en $x=2$. Étudier la position de la courbe par rapport à cette tangente.
\begin{corrige}[exercice~\ref{t11s-c04-e03}]\label{sol-t11s-c04-e03}
$f(2)=-2$, $f'(2)=1$, donc $y=x-4$. L’écart $f(x)-(x-4)=(x-2)^2\ge0$; la courbe est au-dessus, égalité seulement en $2$.
\end{corrige}
\exo[Variations]\label{t11s-c04-e04}
Étudier les variations de $f(x)=x^3-3x$ sur $\R$.
\begin{corrige}[exercice~\ref{t11s-c04-e04}]\label{sol-t11s-c04-e04}
$f'=3(x-1)(x+1)$. Croissance sur $]-\infty;-1]$, décroissance sur $[-1;1]$, croissance sur $[1;+\infty[$. $f(-1)=2$, $f(1)=-2$; limites $-\infty,+\infty$.
\end{corrige}
\exo[Continuité et dérivabilité]\label{t11s-c04-e05}
Pour $g(x)=|x-2|$, étudier la continuité et la dérivabilité en $2$, puis donner les pentes de part et d’autre.
\begin{corrige}[exercice~\ref{t11s-c04-e05}]\label{sol-t11s-c04-e05}
Continue car $|x-2|\to0=g(2)$. Taux $|h|/h$ : $-1$ à gauche, $1$ à droite; pas dérivable en $2$.
\end{corrige}
\exo[Optimiser]\label{t11s-c04-e06}
Un rectangle a un périmètre de $40$ m. Exprimer son aire en fonction d’un côté $x$, préciser l’intervalle et déterminer l’aire maximale.
\begin{corrige}[exercice~\ref{t11s-c04-e06}]\label{sol-t11s-c04-e06}
L’autre côté vaut $20-x$, et $0<x<20$. $A=20x-x^2$; $A'=20-2x$, positive avant $10$, négative après. Maximum $100\ \mathrm{m}^2$ pour le carré de côté $10$ m.
\end{corrige}
\exo[Raccord lisse]\label{t11s-c04-e07}
$f(x)=x^2$ pour $x\le1$, et $f(x)=ax+b$ pour $x>1$. Trouver $a,b$ pour que $f$ soit continue et dérivable en $1$.
\begin{corrige}[exercice~\ref{t11s-c04-e07}]\label{sol-t11s-c04-e07}
Continuité : $a+b=1$. Taux gauche $2+h\to2$; une fois la continuité imposée, taux droit $a$. Dérivabilité : $a=2$, puis $b=-1$.
\end{corrige}

\exo[Une tangente commune ?]\label{t11s-c04-p01}
Pour $f(x)=x^2$ et $g(x)=2x^2-2x+1$, trouver les points d’intersection des courbes. Calculer leurs tangentes au point commun. Ces courbes ont-elles seulement la même valeur ou aussi la même tangente ? Déterminer leur position relative sur $\R$.
\begin{corrige}[exercice~\ref{t11s-c04-p01}]\label{sol-t11s-c04-p01}
$g-f=(x-1)^2$. L’unique intersection est $(1;1)$. Les dérivées $2x$ et $4x-2$ valent toutes deux $2$ en $1$; les deux tangentes sont $y=2x-1$. Comme $(x-1)^2\ge0$, la courbe de $g$ est au-dessus de celle de $f$, strictement sauf en $1$. Le contact commun est ainsi justifié par valeur et pente, pas seulement par une intersection.
\end{corrige}
\exo[Un enclos contre un mur]\label{t11s-c04-p02}
On dispose de $30$ m de grillage pour trois côtés d’un enclos rectangulaire; le quatrième côté est un mur et ne demande pas de grillage. On note $x>0$ la longueur, en mètres, de chacun des deux côtés perpendiculaires au mur. Exprimer la troisième longueur, le domaine et l’aire $A(x)$. Déterminer les dimensions qui maximisent cette aire, par dérivation puis par une forme canonique.
\begin{corrige}[exercice~\ref{t11s-c04-p02}]\label{sol-t11s-c04-p02}
La longueur parallèle au mur vaut $30-2x$, donc $0<x<15$. L’aire est $A=30x-2x^2$ et $A'=30-4x$. La croissance puis décroissance autour de $x=7{,}5$ donne le maximum $112{,}5\ \mathrm{m}^2$, avec côtés $7{,}5$ m et $15$ m. Aussi $A=112{,}5-2(x-7{,}5)^2$ prouve que ce maximum global est unique. Le rectangle optimal n’est pas un carré, car le grillage ne couvre que trois côtés.
\end{corrige}

\begin{jemeteste}Je contrôle le domaine de dérivation, je distingue continuité et dérivabilité, et je justifie une variation par le signe de la dérivée.\end{jemeteste}
\begin{notepourlaclasse}Faire écrire le taux avant d’introduire la notation. Le carré et la valeur absolue fournissent des exemples contrastés. Pour une optimisation, faire préciser unités et domaine avant la dérivation.\end{notepourlaclasse}

\clearpage\bmtchapitre{Étude complète et représentation graphique}{Assembler domaines, limites, dérivée, points remarquables et positions relatives pour tracer une courbe justifiée.}{Analyse}
\label{t11s-c05}
\begin{revision}Les quatre chapitres précédents; signe d’un polynôme; équations de droites.\end{revision}
\begin{automatismes}\exo\label{t11s-c05-auto}Factoriser $x^2-2x-3$. Calculer $f(0)$ pour $f(x)=x+1+2/(x-1)$.\end{automatismes}
\begin{corrige}[exercice~\ref{t11s-c05-auto}]\label{sol-t11s-c05-auto}$(x-3)(x+1)$; $f(0)=-1$.\end{corrige}

\section{Organiser l’étude}
\begin{methode}{une étude en six étapes}
Domaine et éventuelle réduction par symétrie/périodicité; limites et branches infinies; dérivée et variations; points d’intersection avec les axes; tangentes ou points particuliers; tracé cohérent avec l’ensemble des résultats. Pour comparer deux courbes, étudier le signe de leur différence sur leur domaine commun.
\end{methode}
\begin{exemple}[une fraction rationnelle]
Soit $f(x)=\dfrac{x^2+1}{x-1}=x+1+\dfrac2{x-1}$.
Le domaine est $\R\setminus\{1\}$. Les limites en $1^-$ et $1^+$ sont $-\infty$ et $+\infty$. Aux deux infinis, l’asymptote est $y=x+1$.
\[
f'(x)=1-\frac2{(x-1)^2}
=\frac{(x-1)^2-2}{(x-1)^2}.
\]
Elle s’annule en $a=1-\sqrt2$ et $b=1+\sqrt2$. Elle est positive en dehors de $[a;b]$ et négative à l’intérieur, en séparant les deux intervalles du domaine. Les valeurs sont $f(a)=2-2\sqrt2$ et $f(b)=2+2\sqrt2$.
\begin{center}\begin{tabular}{c|ccccccc}
$x$&$-\infty$&&$a$&&$1$&&\\
$f'$&&$+$&$0$&$-$&non défini&&\\
$f$&$-\infty$&$\nearrow$&$2-2\sqrt2$&$\searrow$&$-\infty$&&\\\hline
$x$&$1$&&$b$&&$+\infty$&&\\
$f'$&non défini&$-$&$0$&$+$&&&\\
$f$&$+\infty$&$\searrow$&$2+2\sqrt2$&$\nearrow$&$+\infty$&&
\end{tabular}\end{center}
$f(0)=-1$. Le numérateur $x^2+1$ ne s’annule jamais, donc la courbe ne coupe pas l’axe des abscisses. Pour tout $h\ne0$, $f(1+h)+f(1-h)=4$, et le domaine est symétrique par rapport à $1$ : $(1;2)$ est centre de symétrie.
\end{exemple}
\begin{visualisation}[deux branches, quatre variations]
\begin{center}\begin{tikzpicture}\begin{axis}[width=10cm,height=7cm,axis lines=middle,xmin=-5,xmax=7,ymin=-6,ymax=10,xlabel={$x$},ylabel={$y$},samples=160]
\addplot[bmtSarcelle,thick,domain=-5:0.73]{x+1+2/(x-1)};
\addplot[bmtSarcelle,thick,domain=1.27:7]{x+1+2/(x-1)};
\addplot[bmtOcre,dashed,domain=-5:7]{x+1};
\draw[bmtOcre,dashed](axis cs:1,-6)--(axis cs:1,10);
\end{axis}\end{tikzpicture}\end{center}
Le maximum de la branche gauche et le minimum de la branche droite ne sont pas les mêmes points. L’asymptote verticale sépare le domaine.
\end{visualisation}
\section{Point d’inflexion et tangente}
\begin{definition}[point d’inflexion]Un point d’inflexion est un point où la courbe change de convexité. Pour les fonctions deux fois dérivables étudiées ici, un changement de signe de $f''$ de part et d’autre fournit un critère suffisant (admis). La convexité désigne l’allure courbée vers le haut, la concavité vers le bas.\end{definition}
\begin{exemple}Pour $g(x)=x^3$, $g''(x)=6x$ change de signe en $0$. L’origine est un point d’inflexion, et la tangente $y=0$ est traversée par la courbe. Pour $x^4$, $g''(0)=0$ mais $g''$ reste positive : pas d’inflexion en $0$.\end{exemple}
\begin{avertissement}Annuler $f''$ ne suffit pas. Une courbe apparemment proche d’une droite ne prouve pas une asymptote. Une esquisse est un résumé de résultats déjà justifiés.\end{avertissement}
\section{Comparer deux fonctions}
Si $f-g>0$ sur un intervalle, $\Cf_f$ est au-dessus de $\Cf_g$; si $f-g<0$, elle est en dessous; les zéros de la différence donnent les intersections. Les unités des axes peuvent être différentes; préciser le repère pour les arguments métriques.

% ENRICHISSEMENT C05

\subsection*{Lire une étude comme un ensemble de contraintes}
\begin{methode}{construire la courbe dans le bon ordre}
Tracer d’abord les axes et les asymptotes. Placer ensuite les points connus, en particulier extrema et intersections. Dessiner séparément chaque intervalle du domaine, en respectant les sens de variation et la position par rapport aux droites. Revenir au tableau pour contrôler chaque branche : la calculatrice peut aider à visualiser, mais ne remplace aucune justification.
\end{methode}
\begin{exemple}[compter des intersections sans formule des racines]
Pour $p(x)=x^3-3x$, la dérivée est $3(x-1)(x+1)$. Les intervalles de monotonie stricte sont $]-\infty;-1]$, $[-1;1]$ et $[1;+\infty[$, avec valeurs remarquables $2$ et $-2$. Une horizontale $y=k$ avec $-2<k<2$ rencontre chaque branche exactement une fois : la continuité donne l’existence et la monotonie l’unicité sur chaque branche. Ainsi on peut compter trois solutions sans les écrire en formule explicite.
\end{exemple}
\begin{methode}{confronter une équation au domaine}
Pour $h(x)=(2x+1)/(x-1)$, résoudre $h(x)=c$ exige $x\ne1$. Multiplier alors par $x-1$ donne $(2-c)x=-1-c$. Si $c=2$, on obtient $0=-3$, impossible. Sinon, $x=(1+c)/(c-2)$; ce nombre n’est jamais $1$, car cela imposerait $1=-2$. Le niveau de l’asymptote horizontale n’est pas atteint dans cet exemple, mais cette propriété ne vaut pas pour toute asymptote.
\end{methode}
\begin{avertissement}[les points remarquables n’ont pas le même rôle]
Un extremum local concerne les valeurs voisines; un extremum global concerne tout le domaine. Un point d’inflexion marque un changement de convexité et peut avoir une tangente non horizontale. Une valeur interdite ne devient pas un point de la courbe même si sa limite est finie.
\end{avertissement}

\section{Exercices et problèmes}
\exo[Une parabole]\label{t11s-c05-e01}
Étudier complètement $f(x)=x^2-4x+3$ : domaine, limites, variations, axe et intersections avec les axes. Tracer.
\begin{corrige}[exercice~\ref{t11s-c05-e01}]\label{sol-t11s-c05-e01}
Domaine $\R$, limites $+\infty$ aux deux infinis. $f'=2x-4$; décroissance jusqu’à $2$, croissance ensuite; minimum $-1$. Forme $(x-2)^2-1$, axe $x=2$. Zéros $1,3$; ordonnée à l’origine $3$.
\end{corrige}
\exo[Courbe et droite]\label{t11s-c05-e02}
Comparer $f(x)=x^2-4x+3$ et $g(x)=x-1$.
\begin{corrige}[exercice~\ref{t11s-c05-e02}]\label{sol-t11s-c05-e02}
$f-g=x^2-5x+4=(x-1)(x-4)$. Intersections $(1;0)$ et $(4;3)$. Parabole au-dessus à l’extérieur de $[1;4]$, en dessous sur $]1;4[$.
\end{corrige}
\exo[Un quotient simple]\label{t11s-c05-e03}
Étudier $h(x)=\dfrac{2x+1}{x-1}$ : limites, variations, asymptotes, intersections et centre de symétrie.
\begin{corrige}[exercice~\ref{t11s-c05-e03}]\label{sol-t11s-c05-e03}
$h=2+3/(x-1)$, domaine $x\ne1$; limites $2$ aux infinis, $-\infty$ et $+\infty$ en $1^-,1^+$. $h'=-3/(x-1)^2<0$ sur chaque intervalle. Asymptotes $x=1,y=2$. Intersections $(-1/2;0)$ et $(0;-1)$; centre $(1;2)$ par la somme $h(1+t)+h(1-t)=4$.
\end{corrige}
\exo[Cubicité]\label{t11s-c05-e04}
Étudier $p(x)=x^3-3x$, sa symétrie, ses extrema et son point d’inflexion. Donner sa tangente en $0$.
\begin{corrige}[exercice~\ref{t11s-c05-e04}]\label{sol-t11s-c05-e04}
Domaine $\R$, impaire. $p'=3(x-1)(x+1)$ : maximum local $(-1;2)$, minimum local $(1;-2)$. Limites $-\infty,+\infty$. $p''=6x$ change de signe en $0$; tangente $y=-3x$. Zéros $-\sqrt3,0,\sqrt3$.
\end{corrige}
\exo[Lecture argumentée]\label{t11s-c05-e05}
Un élève affirme : « $f'(a)=0$, donc $f(a)$ est un maximum ». Donner deux contre-exemples différents.
\begin{corrige}[exercice~\ref{t11s-c05-e05}]\label{sol-t11s-c05-e05}
$x^2$ en $0$ a un minimum et dérivée nulle; $x^3$ en $0$ n’a aucun extremum et dérivée nulle. Le signe de part et d’autre est nécessaire pour le critère utilisé.
\end{corrige}
\exo[Oscillation]\label{t11s-c05-e06}
Étudier $u(x)=2\cos x$ sur $[0;\pi]$, puis expliquer comment obtenir sa courbe sur $\R$.
\begin{corrige}[exercice~\ref{t11s-c05-e06}]\label{sol-t11s-c05-e06}
$u'=-2\sin x\le0$ et strictement négative sur $]0;\pi[$. Décroît de $2$ à $-2$; zéro $\pi/2$. La parité complète $[-\pi;\pi]$ et la période $2\pi$ permet les translations.
\end{corrige}

\exo[Le nombre de solutions dépend d’un niveau]\label{t11s-c05-p01}
On reprend $p(x)=x^3-3x$ sur $\R$. À partir de son tableau de variation, déterminer le nombre de solutions distinctes de $p(x)=k$ selon le réel $k$. Traiter séparément $k=-2$ et $k=2$; pour ces deux niveaux, factoriser l’équation et donner les solutions exactes.
\begin{corrige}[exercice~\ref{t11s-c05-p01}]\label{sol-t11s-c05-p01}
Pour $k<-2$ ou $k>2$, il y a une solution; pour $-2<k<2$, il y en a trois. En $k=2$, $x^3-3x-2=(x-2)(x+1)^2$ : deux solutions distinctes $-1,2$. En $k=-2$, $x^3-3x+2=(x+2)(x-1)^2$ : deux solutions distinctes $-2,1$. Les extrema appartiennent à deux intervalles adjacents; il ne faut pas les compter deux fois comme points distincts. Le comptage se justifie par continuité et monotonie sur les branches, pas seulement par le degré trois.
\end{corrige}
\exo[Auditer une description de courbe]\label{t11s-c05-p02}
Pour $f(x)=(x^2+1)/(x-1)$, un compte rendu annonce : « la courbe coupe l’axe des abscisses; ses deux branches se rejoignent en $x=1$; elle possède un maximum global en $1-\sqrt2$ ». Réfuter chacune des trois affirmations à l’aide de l’expression ou des résultats du cours et écrire une description corrigée.
\begin{corrige}[exercice~\ref{t11s-c05-p02}]\label{sol-t11s-c05-p02}
Le numérateur $x^2+1$ est strictement positif, donc aucun zéro. La valeur $1$ est interdite et les limites latérales sont infinies; les branches ne se rejoignent pas. Le point d’abscisse $1-\sqrt2$ est un maximum local de la branche gauche, mais la branche droite tend vers $+\infty$ en $1^+$, donc pas de maximum global. Description : deux branches séparées par $x=1$, aucune intersection avec l’axe horizontal, maximum local à gauche et minimum local à droite, asymptote oblique $y=x+1$.
\end{corrige}

\begin{jemeteste}Pour vérifier mon tracé, je confronte chaque branche aux limites, chaque extremum au tableau et chaque intersection à son équation.\end{jemeteste}
\begin{notepourlaclasse}Faire d’abord une étude commune puis demander une esquisse individuelle. La seconde dérivée sert ici au critère d’inflexion; ne pas exiger un cours général de convexité.\end{notepourlaclasse}

\clearpage\begingroup
\chapter*{Évaluation : analyse}\markboth{Évaluation : analyse}{Évaluation : analyse}\addcontentsline{toc}{chapter}{Évaluation : analyse}
\renewcommand{\thebmtexo}{DS1.\arabic{bmtexo}}
\renewcommand{\theHbmtexo}{DS1.\arabic{bmtexo}}
\setcounter{bmtexo}{0}
\begin{devoir}[2 h — 20 points]
Calculatrice autorisée. Justifier les réponses et préciser les domaines utiles. Les questions sont indépendantes sauf indication. 
\end{devoir}
\exo[Limites et raccord]\label{t11s-ds1-e01}\bareme{6 points}
\begin{enumerate}
\item Calculer $\lim_{x\to2}\dfrac{x^2-4}{x-2}$ (2 points).
\item Pour $x<0$, $f(x)=x^2+1$; pour $x\ge0$, $f(x)=ax+b$. Déterminer $b$ pour la continuité, puis $a$ pour la dérivabilité en $0$ (4 points).
\end{enumerate}
\begin{corrige}[exercice~\ref{t11s-ds1-e01}]\label{sol-t11s-ds1-e01}
1. Pour $x\ne2$, quotient $x+2$, limite $4$ (factorisation 1, limite 1).
2. Continuité : limite gauche $1$, valeur et limite droite $b$, donc $b=1$ (2). Avec ce choix, taux gauche $h\to0$ et taux droit $a$, donc $a=0$ (2).
\end{corrige}
\exo[Étude d’une fraction rationnelle]\label{t11s-ds1-e02}\bareme{8 points}
Soit $g(x)=\dfrac{x^2-2x+2}{x-1}$.
\begin{enumerate}
\item Domaine et décomposition $g(x)=x-1+\dfrac1{x-1}$ (1 point).
\item Limites, asymptotes et position relative à l’oblique (3 points).
\item Dérivée, tableau de variation et extrema (3 points).
\item Déterminer un centre de symétrie et donner une esquisse (1 point).
\end{enumerate}
\begin{corrige}[exercice~\ref{t11s-ds1-e02}]\label{sol-t11s-ds1-e02}
1. Domaine $x\ne1$; numérateur $(x-1)^2+1$.
2. Limites en $1^-:-\infty$, $1^+:+\infty$; aux infinis $-\infty,+\infty$. Asymptotes $x=1,y=x-1$; écart $1/(x-1)$, sous l’oblique à gauche, au-dessus à droite.
3. $g'=1-1/(x-1)^2=x(x-2)/(x-1)^2$. Croît sur $]-\infty;0]$, décroît sur $[0;1[$ et $]1;2]$, croît sur $[2;+\infty[$. $g(0)=-2$, $g(2)=2$.
4. $g(1+h)+g(1-h)=0$ avec domaine symétrique : centre $(1;0)$. Esquisse à deux branches conforme aux limites et extrema. Les points critiques et chaque branche doivent être indiqués.
\end{corrige}
\exo[Aire et tangente]\label{t11s-ds1-e03}\bareme{6 points}
Un rectangle a périmètre $24$ m.
\begin{enumerate}\item Exprimer l’aire en fonction d’un côté $x$ et préciser le domaine (2 points).
\item Déterminer l’aire maximale (2 points).
\item Donner la tangente à la courbe de l’aire en $x=4$ et étudier la position de cette courbe par rapport à elle (2 points).\end{enumerate}
\begin{corrige}[exercice~\ref{t11s-ds1-e03}]\label{sol-t11s-ds1-e03}
$A(x)=x(12-x)$, $0<x<12$ (1+1). $A'=12-2x$, maximum en $6$, aire $36\ \mathrm{m}^2$ (1+1). $A(4)=32,A'(4)=4$, tangente $y=4x+16$; écart $-(x-4)^2\le0$, courbe en dessous (1+1).
\end{corrige}
\endgroup


% ===========================================================================
\bmtpartie{bmtGrenat}{II}{Algèbre}{Raisonnement et calcul}{%
  Cette partie s'ouvre sur la logique et se ferme sur les suites, et ce n'est
  pas un ordre arbitraire. Entre les deux, on résout : des équations, des
  systèmes, des questions de divisibilité. Chaque fois, la même exigence —
  dire ce que l'on cherche, dire ce que l'on sait, dire ce qui permet de
  passer de l'un à l'autre.}

\clearpage\bmtchapitre{Logique et raisonnement}{Formuler et nier des propositions, distinguer réciproque et contraposée et construire une démonstration ou un contre-exemple.}{Algèbre}
\label{t11s-c06}
\begin{revision}Comparer des nombres, reconnaître les entiers pairs et résoudre une équation élémentaire.\end{revision}
\begin{automatismes}\exo\label{t11s-c06-auto}Écrire trois entiers pairs. L’affirmation « $3^2=6$ » est-elle vraie ? Résoudre $x^2=4$ dans $\R$.\end{automatismes}
\begin{corrige}[exercice~\ref{t11s-c06-auto}]\label{sol-t11s-c06-auto}Par exemple $0,2,4$; affirmation fausse; solutions $-2$ et $2$.\end{corrige}

\section{Propositions et connecteurs}
\begin{definition}Une proposition est un énoncé ayant une valeur de vérité. $P\wedge Q$ signifie « $P$ et $Q$ »; $P\vee Q$ signifie « $P$ ou $Q$ », au moins une étant vraie (ou inclusif). $\neg P$ est la négation. L’implication $P\Rightarrow Q$ est fausse seulement lorsque $P$ est vraie et $Q$ fausse. L’équivalence $P\Leftrightarrow Q$ signifie les deux implications.\end{definition}
\begin{center}\begin{tabular}{cc|cccc}\toprule
$P$&$Q$&$P\wedge Q$&$P\vee Q$&$P\Rightarrow Q$&$P\Leftrightarrow Q$\\\midrule
V&V&V&V&V&V\\V&F&F&V&F&F\\F&V&F&V&V&F\\F&F&F&F&V&V\\\bottomrule
\end{tabular}\end{center}
\begin{propriete}[négation et contraposée]
$\neg(P\wedge Q)\Leftrightarrow(\neg P\vee\neg Q)$;
$\neg(P\vee Q)\Leftrightarrow(\neg P\wedge\neg Q)$.
La contraposée de $P\Rightarrow Q$ est $\neg Q\Rightarrow\neg P$ : elle est équivalente à l’implication initiale. La réciproque $Q\Rightarrow P$ ne l’est pas en général.
\end{propriete}
\begin{demonstration}Les deux implications $P\Rightarrow Q$ et $\neg Q\Rightarrow\neg P$ sont fausses exactement dans le cas $P$ vraie, $Q$ fausse; elles ont donc la même table. Les lois de négation se vérifient sur les quatre lignes.\end{demonstration}
\begin{exemple}« Si un entier est divisible par $6$, il est divisible par $3$ » est vrai. Sa réciproque est fausse : $3$ est divisible par $3$, mais pas par $6$. Sa contraposée, « s’il n’est pas divisible par $3$, il n’est pas divisible par $6$ », est vraie.\end{exemple}
\section{Quantifier et nier}
\begin{definition}$\forall x\in E$ signifie « pour tout $x$ de $E$ »; $\exists x\in E$ signifie « il existe au moins un $x$ de $E$ ». Le domaine $E$ fait partie de l’énoncé.\end{definition}
\begin{propriete}
$\neg(\forall x\in E, P(x))\Leftrightarrow\exists x\in E,\neg P(x)$.
De même, $\neg(\exists x\in E,P(x))\Leftrightarrow\forall x\in E,\neg P(x)$.
\end{propriete}
\begin{demonstration}Dire que tous les éléments vérifient $P$ est faux exactement lorsqu’au moins un élément ne la vérifie pas. Dire qu’il en existe un est faux lorsque tous échouent.\end{demonstration}
\begin{exemple}$\forall x\in\R,\exists y\in\R,y>x$ est vrai : choisir $y=x+1$. $\exists y\in\R,\forall x\in\R,y>x$ est faux : pour le $y$ proposé, prendre $x=y$. L’ordre des quantificateurs change le sens.\end{exemple}
\section{Choisir une preuve}
\begin{methode}{trois voies}
Preuve directe : partir des hypothèses et atteindre la conclusion. Contraposée : supposer la conclusion fausse et montrer l’hypothèse fausse. Absurde : supposer l’hypothèse vraie et la conclusion fausse, puis obtenir une contradiction. Pour réfuter une affirmation universelle, un contre-exemple suffit; pour la prouver, des exemples ne suffisent pas.
\end{methode}
\begin{exemple}[une preuve directe]Si $n=2k$ avec $k\in\Z$, alors $n^2=4k^2=2(2k^2)$ est pair. Pour la réciproque, si $n$ est impair, $n=2k+1$, et $n^2=2(2k^2+2k)+1$ est impair : contraposée.\end{exemple}

% ENRICHISSEMENT C06

\subsection*{Du langage courant à une preuve contrôlable}
\begin{methode}{identifier ce qui est nécessaire ou suffisant}
Dans $P\Rightarrow Q$, vérifier $P$ suffit pour conclure $Q$; vérifier $Q$ ne permet pas en général de remonter à $P$. On dit que $P$ est une condition suffisante de $Q$, et $Q$ une condition nécessaire de $P$. Pour une équivalence, prouver séparément les deux sens, sauf si chaque transformation intermédiaire est elle-même équivalente.
\end{methode}
\begin{exemple}[une condition qui n’est pas réciproque]
Être divisible par $4$ suffit pour qu’un entier soit pair. Être pair est nécessaire pour être divisible par $4$, mais ne suffit pas : $6$ est pair sans être divisible par $4$. Un seul contre-exemple réfute la réciproque; dix exemples divisibles par $4$ ne la prouvent pas.
\end{exemple}
\begin{methode}{nier une phrase complète}
Fixer le domaine; remplacer chaque « pour tout » par « il existe » et inversement, dans l’ordre; nier le prédicat final. La négation de $<$ est $\ge$, celle de « et » est « ou ». Une négation affirme exactement que l’énoncé initial est faux, sans modifier ses hypothèses ou son domaine.
\end{methode}
\begin{exemple}[nier deux quantificateurs]
La négation de $\forall x\in\R,\exists y\in\R,y^2=x$ est $\exists x\in\R,\forall y\in\R,y^2\ne x$. La première est fausse et la seconde vraie : prendre $x=-1$, car tous les carrés réels sont positifs ou nuls. Écrire seulement « il existe $y$ tel que $y^2\ne x$ » aurait laissé un autre sens.
\end{exemple}
\begin{notepourlaclasse}Faire expliciter les hypothèses avant les calculs. Un contre-exemple doit satisfaire toutes les hypothèses et contredire la conclusion; changer le domaine n’est pas réfuter la proposition initiale.\end{notepourlaclasse}

\section{Exercices et problèmes}
\exo[Nier]\label{t11s-c06-e01}
Nier : « $x>2$ et $x<5$ »; « tout entier est pair »; « il existe un réel dont le carré est $-1$ ».
\begin{corrige}[exercice~\ref{t11s-c06-e01}]\label{sol-t11s-c06-e01}
$x\le2$ ou $x\ge5$; il existe un entier non pair; pour tout réel $x$, $x^2\ne-1$.
\end{corrige}
\exo[Implications]\label{t11s-c06-e02}
Pour « si $x>2$, alors $x^2>4$ » ($x\in\R$), écrire la réciproque et la contraposée; déterminer leur vérité.
\begin{corrige}[exercice~\ref{t11s-c06-e02}]\label{sol-t11s-c06-e02}
Initiale vraie. Réciproque « si $x^2>4$, alors $x>2$ » fausse avec $x=-3$. Contraposée « si $x^2\le4$, alors $x\le2$ » vraie.
\end{corrige}
\exo[Un ou inclusif]\label{t11s-c06-e03}
Un élève écrit : « $ab=0$ signifie qu’exactement un des deux facteurs est nul ». Corriger et donner un contre-exemple.
\begin{corrige}[exercice~\ref{t11s-c06-e03}]\label{sol-t11s-c06-e03}
Dans $\R$, $ab=0\Leftrightarrow a=0$ ou $b=0$; les deux peuvent être nuls. Prendre $a=b=0$.
\end{corrige}
\exo[Ordre des quantificateurs]\label{t11s-c06-e04}
Comparer $\forall n\in\N,\exists m\in\N,m=n+1$ et $\exists m\in\N,\forall n\in\N,m=n+1$. Nier chacune.
\begin{corrige}[exercice~\ref{t11s-c06-e04}]\label{sol-t11s-c06-e04}
La première est vraie; la seconde fausse (les choix $n=0$ et $1$ imposeraient $m=1$ et $2$). Négations : $\exists n,\forall m,m\ne n+1$; et $\forall m,\exists n,m\ne n+1$.
\end{corrige}
\exo[Une équivalence]\label{t11s-c06-e05}
Montrer pour $x\in\R$ que $|x|=x$ si et seulement si $x\ge0$.
\begin{corrige}[exercice~\ref{t11s-c06-e05}]\label{sol-t11s-c06-e05}
Si $x\ge0$, définition de la valeur absolue : $|x|=x$. Si $|x|=x$, le membre gauche est positif ou nul, donc $x\ge0$. Les deux sens sont nécessaires.
\end{corrige}
\exo[Faux ou mal formulé]\label{t11s-c06-e06}
« Pour tout réel $x$, $1/x>0$ ». Identifier d’abord le défaut de domaine, puis réfuter la version sur $\R^*$.
\begin{corrige}[exercice~\ref{t11s-c06-e06}]\label{sol-t11s-c06-e06}
L’expression n’existe pas en $0$; le domaine doit exclure $0$. Sur $\R^*$, $x=-1$ donne $1/x=-1<0$.
\end{corrige}
\exo[Divisibilité]\label{t11s-c06-e07}
Démontrer que la somme de deux entiers impairs est paire. La réciproque « si la somme est paire, les deux nombres sont impairs » est-elle vraie ?
\begin{corrige}[exercice~\ref{t11s-c06-e07}]\label{sol-t11s-c06-e07}
$(2k+1)+(2\ell+1)=2(k+\ell+1)$ est pair. Réciproque fausse : $2+4=6$ avec deux entiers pairs.
\end{corrige}

\exo[Un domaine change la vérité]\label{t11s-c06-p01}
Considérer $P$ : « pour tout réel $x$, $x^2\ge x$ ». Écrire sa négation, réfuter $P$ et déterminer tous les réels qui satisfont $x^2<x$. La proposition obtenue en remplaçant $\R$ par $\Z$ est-elle vraie ? Démontrer votre réponse.
\begin{corrige}[exercice~\ref{t11s-c06-p01}]\label{sol-t11s-c06-p01}
La négation est $\exists x\in\R,x^2<x$. Le réel $1/2$ convient. L’inéquation $x(x-1)<0$ donne $]0;1[$. Cet intervalle ne contient aucun entier; tout entier est soit $n\le0$, soit $n\ge1$, et dans les deux cas $n(n-1)\ge0$. La proposition sur $\Z$ est vraie. Un contre-exemple réel non entier ne réfute pas cette dernière.
\end{corrige}
\exo[Réparer une conclusion trop forte]\label{t11s-c06-p02}
Pour $x,y\in\R$, un élève affirme : « si $xy>0$, alors $x>0$ et $y>0$ ». Donner un contre-exemple. Remplacer la conclusion par une condition équivalente à $xy>0$, puis justifier les deux sens. Écrire la contraposée de l’affirmation erronée et dire si elle est vraie.
\begin{corrige}[exercice~\ref{t11s-c06-p02}]\label{sol-t11s-c06-p02}
$x=y=-1$ donne un produit $1>0$ avec deux facteurs négatifs. La réparation est : « les deux facteurs sont strictement positifs ou tous deux strictement négatifs ». Dans chacun de ces cas le produit est positif; réciproquement, un produit positif exclut un facteur nul et des signes opposés, donc impose ces deux cas. La contraposée erronée est « si $x\le0$ ou $y\le0$, alors $xy\le0$ »; le même couple la réfute. Une contraposée est équivalente à l’implication initiale, même si celle-ci est fausse.
\end{corrige}

\begin{jemeteste}Je précise le domaine, je nie chaque quantificateur et je prouve les deux sens d’une équivalence.\end{jemeteste}
\begin{notepourlaclasse}Insister sur la distinction entre une implication logique et un raisonnement causal. Ne pas récompenser une succession d’exemples comme preuve d’un universel.\end{notepourlaclasse}

\clearpage\bmtchapitre{Boucles et organigrammes}{Exécuter et construire une boucle, produire un organigramme et justifier la terminaison d’un algorithme simple.}{Algèbre}
\label{t11s-c07}
\begin{revision}Affectation, tests, reste de division euclidienne et calcul de sommes simples.\end{revision}
\begin{automatismes}\exo\label{t11s-c07-auto}Calculer le reste de $17$ par $5$. Si $a=3$, puis $a\leftarrow a+2$, quelle est sa valeur ? Calculer $1+2+3+4$.\end{automatismes}
\begin{corrige}[exercice~\ref{t11s-c07-auto}]\label{sol-t11s-c07-auto}Reste $2$; nouvelle valeur $5$; somme $10$.\end{corrige}

\section{Exécuter sans sauter une instruction}
\begin{definition}Une affectation $a\leftarrow E$ remplace la valeur de $a$ par la valeur actuelle de $E$. Elle ne constitue pas une égalité à résoudre. Une boucle « pour » parcourt une liste d’indices; une boucle « tant que » répète tant que sa condition est vraie. Le test de « tant que » précède chaque tour, y compris le premier.\end{definition}
\begin{exemple}[additionner les premiers entiers]
\begin{tabular}{l}
Lire un entier $n\ge0$\\
$S\leftarrow0$\\
Pour $k$ allant de $1$ à $n$\\
\quad $S\leftarrow S+k$\\
Fin pour\\
Afficher $S$
\end{tabular}
Pour $n=4$, les valeurs successives de $S$ sont $0,1,3,6,10$. Si $n=0$, aucun tour n’est effectué et la somme vide vaut $0$.
\end{exemple}
\begin{propriete}[invariant de la somme]Après le tour d’indice $k$, $S=1+\cdots+k$.\end{propriete}
\begin{demonstration}Avant le premier tour, $S=0$. Si la somme est correcte au tour précédent, l’ajout de $k$ donne celle des $k$ premiers entiers. Une récurrence justifie chaque tour; après $n$ tours, le résultat annoncé est obtenu.\end{demonstration}
\section{Un organigramme avec retour}
L’organigramme suivant calcule la même somme. Le losange est un test; les rectangles sont des affectations; les flèches indiquent l’ordre d’exécution.
\begin{visualisation}[le test décide de poursuivre]
\begin{center}\begin{tikzpicture}[node distance=9mm, every node/.style={font=\small}, process/.style={draw,rectangle,align=center,minimum width=37mm,minimum height=8mm},test/.style={draw,diamond,aspect=2,inner sep=2pt}]
\node[process](init){Lire $n\in\N$\\$S\leftarrow0$, $k\leftarrow1$};
\node[test,below=of init](test){$k\le n$ ?};
\node[process,below=of test](add){$S\leftarrow S+k$\\$k\leftarrow k+1$};
\node[process,right=20mm of test](end){Afficher $S$};
\draw[->](init)--(test);\draw[->](test)--node[right]{oui}(add);
\draw[->](test)--node[above]{non}(end);
\draw[->](add.west)--++(-12mm,0)|-(test.west);
\end{tikzpicture}\end{center}
Le compteur augmente d’une unité et finit par dépasser $n$. Le nombre de tours est fini.
\end{visualisation}
\section{Boucle conditionnelle et terminaison}
\begin{exemple}[atteindre un seuil]
Avec $u\leftarrow1$, $n\leftarrow0$, tant que $u<20$, faire $u\leftarrow2u$ puis $n\leftarrow n+1$. Les valeurs de $u$ sont $1,2,4,8,16,32$; la sortie est $n=5$. Après $n$ tours, $u=2^n$. La suite des puissances dépasse $20$ : l’arrêt est justifié.
\end{exemple}
\begin{avertissement}Si le corps ne modifie pas ce qui intervient dans le test, la boucle peut être infinie. Pour permuter deux valeurs, $a\leftarrow b$ puis $b\leftarrow a$ ne convient pas : la première valeur de $a$ a été perdue. Utiliser une variable temporaire.\end{avertissement}
\begin{methode}{contrôler un algorithme}Préciser les entrées autorisées; dresser une table d’exécution; expliquer ce qui reste vrai après chaque tour (invariant); expliquer pourquoi le test finit par être faux (terminaison); interpréter la sortie et les cas limites.\end{methode}

% ENRICHISSEMENT C07

\subsection*{Rédiger une table, une preuve et un test}
\begin{exemple}[chaque ligne utilise les valeurs actuelles]
On initialise $a=7,b=3$, puis on exécute $a\leftarrow a-b$ et $b\leftarrow a+b$.
\begin{center}\begin{tabular}{l|cc}\toprule
État&$a$&$b$\\\midrule
Initial&7&3\\Après la première affectation&4&3\\Après la seconde&4&7\\\bottomrule
\end{tabular}\end{center}
Dans la seconde ligne de calcul, $a$ vaut déjà $4$. Remplacer les deux variables simultanément par leurs anciennes valeurs décrirait un autre algorithme.
\end{exemple}
\begin{methode}{séparer résultat correct et arrêt}
Un invariant explique ce que contiennent les variables si l’algorithme atteint une étape. Il ne garantit pas que l’algorithme s’arrête. Pour une boucle bornée, compter les tours; pour une boucle conditionnelle, trouver une quantité qui finit par franchir le seuil. Tester aussi le cas où la condition est fausse dès le départ.
\end{methode}
\begin{exemple}[dénombrer les termes déjà ajoutés]
Pour calculer $u_0+\cdots+u_N$ quand $u_{n+1}=2u_n+1$, initialiser $u=u_0,S=u_0$. Au tour $k=1,\ldots,N$, mettre d’abord $u\leftarrow2u+1$, puis $S\leftarrow S+u$. Après le tour $k$, $u=u_k$ et $S=u_0+\cdots+u_k$. En $N=0$, il n’y a pas de tour et la sortie est déjà correcte. Inverser les deux mises à jour compterait deux fois $u_0$ et oublierait le dernier terme.
\end{exemple}
\begin{avertissement}[atteindre une valeur ou s’en approcher]
Une suite peut tendre vers un seuil sans l’atteindre à un rang fini. Un test d’égalité exacte ne doit pas être remplacé tacitement par un test approché. Pour demander une précision, fixer une tolérance strictement positive dans l’énoncé.
\end{avertissement}

\section{Exercices et problèmes}
\exo[Suivre les valeurs]\label{t11s-c07-e01}
Exécuter pour $n=3$ : $u\leftarrow2$; pour $k=1$ à $n$, $u\leftarrow3u-1$. Donner les valeurs et la sortie.
\begin{corrige}[exercice~\ref{t11s-c07-e01}]\label{sol-t11s-c07-e01}
Valeurs $2,5,14,41$. La sortie est $41$ après exactement trois mises à jour.
\end{corrige}
\exo[Somme de carrés]\label{t11s-c07-e02}
Écrire un algorithme calculant $1^2+\cdots+n^2$ pour $n\in\N$. Tester $n=0$ et $n=3$.
\begin{corrige}[exercice~\ref{t11s-c07-e02}]\label{sol-t11s-c07-e02}
Initialiser $S=0$; pour $k=1$ à $n$, affecter $S\leftarrow S+k^2$; afficher $S$. Sorties $0$ et $14$. Invariant : somme des carrés des indices déjà parcourus.
\end{corrige}
\exo[Échanger]\label{t11s-c07-e03}
Avec $a=4,b=7$, exécuter $a\leftarrow b$, $b\leftarrow a$. Réparer.
\begin{corrige}[exercice~\ref{t11s-c07-e03}]\label{sol-t11s-c07-e03}
On obtient $a=7,b=7$. Réparation : $t\leftarrow a$, $a\leftarrow b$, $b\leftarrow t$; résultat $a=7,b=4$.
\end{corrige}
\exo[Premier seuil]\label{t11s-c07-e04}
Exécuter $u\leftarrow3,n\leftarrow0$; tant que $u\le100$, $u\leftarrow2u$, $n\leftarrow n+1$. Quelle différence ferait le test $u<96$ ?
\begin{corrige}[exercice~\ref{t11s-c07-e04}]\label{sol-t11s-c07-e04}
Valeurs $3,6,12,24,48,96,192$; sortie $n=6$. Avec $u<96$, arrêt à $96$ et $n=5$. Le signe strict change le traitement du seuil atteint.
\end{corrige}
\exo[Une boucle sans fin]\label{t11s-c07-e05}
On initialise $u=1$ et on répète $u\leftarrow(u+1)/2$ tant que $u<2$. La boucle s’arrête-t-elle en calcul exact ?
\begin{corrige}[exercice~\ref{t11s-c07-e05}]\label{sol-t11s-c07-e05}
Non : $u=1$ reste invariant, car $(1+1)/2=1$. Le test reste vrai. La table d’un seul tour suggère l’erreur et l’invariant la prouve.
\end{corrige}
\exo[Dessiner]\label{t11s-c07-e06}
Transformer l’algorithme de la somme des carrés en organigramme avec un test avant la mise à jour. Indiquer les deux branches et l’actualisation de $k$.
\begin{corrige}[exercice~\ref{t11s-c07-e06}]\label{sol-t11s-c07-e06}
Lire $n$; initialiser $S=0,k=1$; tester $k\le n$. Oui : $S\leftarrow S+k^2$, $k\leftarrow k+1$, retour au test. Non : afficher $S$ et terminer. Les tests $n=0,3$ donnent $0,14$.
\end{corrige}

\exo[Une somme avec invariant]\label{t11s-c07-p01}
Pour un entier d’entrée $n\ge0$, on initialise $S=0,k=0$. Tant que $k<n$, on effectue, dans cet ordre, $k\leftarrow k+1$ puis $S\leftarrow S+2k$; on affiche $S$.
\begin{enumerate}
\item Dresser la table pour $n=3$ et traiter séparément $n=0$.
\item Prouver qu’après chaque tour $S=2+4+\cdots+2k$; justifier l’arrêt et interpréter la sortie.
\item Que calculerait-on si l’on ajoutait $2k$ avant d’augmenter $k$ ? Tester $n=3$.
\end{enumerate}
\begin{corrige}[exercice~\ref{t11s-c07-p01}]\label{sol-t11s-c07-p01}
Les états $(k;S)$ sont $(0;0),(1;2),(2;6),(3;12)$. Pour $n=0$, aucun tour, sortie $0$. Initialement la somme est vide; chaque tour augmente $k$ puis ajoute exactement le nouvel entier pair $2k$, ce qui conserve la description. Après $n$ tours, le test est faux et la sortie est la somme des $n$ premiers entiers pairs positifs. Avec l’ordre inversé, on ajoute $0,2,\ldots,2(n-1)$, donc la sortie pour $n=3$ vaut $6$, pas $12$.
\end{corrige}
\exo[La limite ne suffit pas à arrêter]\label{t11s-c07-p02}
En calcul exact, initialiser $u=1,n=0$ et répéter $u\leftarrow u/2$, $n\leftarrow n+1$ tant que $u>0$. Cette boucle termine-t-elle ? Modifier seulement le seuil pour demander $u\le0{,}01$, puis calculer le premier rang d’arrêt et justifier sa minimalité.
\begin{corrige}[exercice~\ref{t11s-c07-p02}]\label{sol-t11s-c07-p02}
Après $n$ tours, $u=2^{-n}>0$; aucun rang fini n’arrête la première boucle. Pour la seconde, tester tant que $u>0{,}01$. L’arrêt est atteint pour $n=7$, car $1/64>1/100$ et $1/128<1/100$. Les termes décroissent strictement; aucun rang plus petit que $7$ ne convient. La réponse concerne le calcul exact, sans hypothèse cachée sur les arrondis d’une machine.
\end{corrige}

\begin{jemeteste}Je vérifie l’initialisation, le premier et le dernier tour, le compteur et le sens du test.\end{jemeteste}
\begin{notepourlaclasse}Faire d’abord exécuter à la main. Le pseudo-code suffit; les invariants sont des justifications accessibles, sans cours de complexité. Relier ensuite aux algorithmes d’Euclide et de suites.\end{notepourlaclasse}

\clearpage\bmtchapitre{Équations et inéquations}{Résoudre les équations du second degré avec paramètre, les cubiques à racine connue, les équations irrationnelles et celles avec valeur absolue.}{Algèbre}
\label{t11s-c08}
\begin{revision}Identités remarquables, discriminant, tableaux de signes, domaines des racines carrées.\end{revision}
\begin{automatismes}\exo\label{t11s-c08-auto}Factoriser $x^2-5x+6$. Résoudre $|x|=3$. Pour quels réels $x$ la racine $\sqrt{x-1}$ existe-t-elle ?\end{automatismes}
\begin{corrige}[exercice~\ref{t11s-c08-auto}]\label{sol-t11s-c08-auto}$(x-2)(x-3)$; $x=-3$ ou $3$; $x\ge1$.\end{corrige}

\section{Second degré et paramètre}
\begin{propriete}Pour $a\ne0$, $ax^2+bx+c=0$ a un discriminant $\Delta=b^2-4ac$. Si $\Delta<0$, aucune racine réelle; si $\Delta=0$, racine double $-b/(2a)$; si $\Delta>0$, deux racines $(-b\pm\sqrt\Delta)/(2a)$. Si $\Delta>0$, le polynôme a le signe de $a$ à l’extérieur des racines et le signe opposé entre elles. Si $\Delta=0$, il a le signe de $a$ sauf à la racine, où il est nul; si $\Delta<0$, il a toujours le signe de $a$.\end{propriete}
\begin{demonstration}La forme canonique est $a[(x+b/(2a))^2-\Delta/(4a^2)]$. Résoudre un carré égal à $\Delta/(4a^2)$ donne les trois cas; pour le signe, factoriser $a(x-x_1)(x-x_2)$ dans le cas de deux racines.\end{demonstration}
\begin{propriete}[somme et produit]Si $x_1,x_2$ sont les racines, leur somme vaut $-b/a$ et leur produit $c/a$. Réciproquement deux nombres de somme $S$ et de produit $P$ sont racines de $X^2-SX+P$.\end{propriete}
\begin{demonstration}Développer $a(X-x_1)(X-x_2)$ et identifier les coefficients.\end{demonstration}
\begin{exemple}$(m-1)x^2-2mx+m+1=0$. Pour $m=1$, l’équation est linéaire : $x=1$. Pour $m\ne1$, $\Delta=4$ et les racines sont $1$ et $(m+1)/(m-1)$. Ne pas appliquer la formule avant de vérifier le coefficient du carré.\end{exemple}
\section{Degré trois et équations bicarrées}
\begin{propriete}[racine connue]Si $P(r)=0$, le polynôme $P$ est divisible par $x-r$. Pour une cubique, le quotient est du second degré.\end{propriete}
\begin{demonstration}La division donne $P(x)=(x-r)Q(x)+c$, avec reste constant. En $x=r$, $c=P(r)=0$.\end{demonstration}
\begin{exemple}[division euclidienne]
Pour $P(x)=x^3-2x^2-x+2$, on vérifie $P(1)=0$. La division s’effectue ainsi :
\[
\begin{array}{r|l}
x^3-2x^2-x+2&x-1\\
\underline{-(x^3-x^2)}&x^2-x-2\\
-x^2-x+2&\\
\underline{-(-x^2+x)}&\\
-2x+2&\\
\underline{-(-2x+2)}&\\
0&
\end{array}
\]
Donc $P=(x-1)(x^2-x-2)=(x-1)(x-2)(x+1)$. Une équation bicarrée se ramène au second degré en posant $X=x^2\ge0$; chaque solution $X>0$ fournit $x=\pm\sqrt X$.
\end{exemple}
\section{Racines carrées : conserver les conditions}
\begin{propriete}
$\sqrt{A(x)}=B(x)$ équivaut à $A(x)\ge0$, $B(x)\ge0$ et $A(x)=B(x)^2$.
Pour $\sqrt A\le B$, il faut $A\ge0,B\ge0,A\le B^2$.
Pour $\sqrt A>B$, si $B<0$, tous les points du domaine conviennent; si $B\ge0$, comparer $A$ à $B^2$.
\end{propriete}
\begin{demonstration}Le carré est strictement croissant sur $[0;+\infty[$. Les équivalences ne valent donc qu’après avoir établi la positivité des deux membres.\end{demonstration}
\begin{exemple}$\sqrt{x+2}=x$ exige $x\ge0$. Le carré donne $(x-2)(x+1)=0$; seule $2$ respecte la condition. La solution $-1$ de l’équation au carré est rejetée.\end{exemple}
\section{Valeurs absolues et tableaux de signes}
\begin{propriete}Pour $k\ge0$, $|u|\le k\Leftrightarrow-k\le u\le k$, et $|u|=k\Leftrightarrow u=k$ ou $u=-k$. Pour une différence de valeurs absolues, séparer les intervalles où chaque expression garde un signe constant.\end{propriete}
\begin{methode}{résoudre une inéquation quotient}Exclure les zéros du dénominateur; factoriser numérateur et dénominateur; dresser le tableau des signes; choisir les intervalles et inclure uniquement les zéros autorisés pour une inégalité large.\end{methode}

% ENRICHISSEMENT C08

\subsection*{Conserver l’équivalence à chaque étape}
\begin{methode}{un paramètre ne se divise pas sans discussion}
Isoler d’abord les valeurs qui annulent un coefficient servant de diviseur. Résoudre ces cas dans l’équation originale; résoudre ensuite le cas général. À la fin, préciser si les solutions coïncident et si elles respectent d’éventuelles contraintes du contexte. Une formule non définie pour une valeur du paramètre ne signifie pas automatiquement « aucune solution ».
\end{methode}
\begin{exemple}[le degré peut changer]
L’équation $mx^2-x=0$ se factorise en $x(mx-1)=0$. Pour $m=0$, elle devient $-x=0$ et admet $0$. Pour $m\ne0$, elle admet $0$ et $1/m$. On ne pouvait pas écrire $1/m$ avant d’examiner $m=0$.
\end{exemple}
\begin{methode}{rédiger une inéquation irrationnelle}
Avec $\sqrt{A(x)}\ge B(x)$, commencer par $A(x)\ge0$. Si $B(x)\le0$, l’inégalité est automatiquement vraie sur ce domaine. Si $B(x)>0$, mettre au carré donne une comparaison équivalente. Réunir ensuite les solutions des cas; ne pas seulement lister les racines de l’équation frontière.
\end{methode}
\begin{exemple}[le signe de l’autre membre décide]
Pour $\sqrt{x+1}\ge x-1$, le domaine est $x\ge-1$. Si $x\le1$, le membre droit est négatif ou nul, donc l’inégalité est vraie. Si $x>1$, le carré donne $x+1\ge(x-1)^2$, soit $x(x-3)\le0$; avec $x>1$, on garde $]1;3]$. La réunion est $[-1;3]$. Les deux bornes satisfont effectivement l’inégalité.
\end{exemple}
\begin{avertissement}[simplifier n’ajoute pas une valeur]
Si un facteur disparaît dans un quotient, sa valeur interdite reste exclue pour l’expression initiale. Dans un tableau de signes, marquer d’abord tous les pôles, puis décider des bornes de l’ensemble de solutions.
\end{avertissement}

\section{Exercices et problèmes}
\exo[Paramètre et degré]\label{t11s-c08-e01}
Résoudre $(m-1)x^2-(m+1)x+2=0$ pour tout $m\in\R$.
\begin{corrige}[exercice~\ref{t11s-c08-e01}]\label{sol-t11s-c08-e01}
Factorisation $(x-1)((m-1)x-2)$. Pour $m=1$, seule solution $1$. Pour $m\ne1$, solutions $1$ et $2/(m-1)$, qui coïncident pour $m=3$.
\end{corrige}
\exo[Somme et produit]\label{t11s-c08-e02}
Un rectangle a un périmètre de $26$ m et une aire de $40\ \mathrm{m}^2$. Trouver ses dimensions.
\begin{corrige}[exercice~\ref{t11s-c08-e02}]\label{sol-t11s-c08-e02}
Les côtés ont somme $13$ et produit $40$. Ils sont racines de $X^2-13X+40=(X-5)(X-8)$. Dimensions $5$ m et $8$ m.
\end{corrige}
\exo[Cubique]\label{t11s-c08-e03}
Un énoncé propose : « $1$ est racine de $P(x)=x^3-4x^2+x+6$; diviser par $x-1$ ». Contrôler cette affirmation, trouver une racine parmi $-1,1,2$ et résoudre $P(x)=0$.
\begin{corrige}[exercice~\ref{t11s-c08-e03}]\label{sol-t11s-c08-e03}
$P(1)=4$, donc $1$ n’est pas racine : l’énoncé contient une fausse hypothèse. Vérifier au lieu de diviser aveuglément. Une racine est $2$; $P=(x-2)(x-3)(x+1)$, racines $-1,2,3$.
\end{corrige}
\exo[Bicarrée]\label{t11s-c08-e04}
Résoudre $x^4-5x^2+4=0$.
\begin{corrige}[exercice~\ref{t11s-c08-e04}]\label{sol-t11s-c08-e04}
Poser $X=x^2$. $(X-1)(X-4)=0$; les deux valeurs sont positives. Solutions $-2,-1,1,2$.
\end{corrige}
\exo[Racine et inégalité]\label{t11s-c08-e05}
Résoudre $\sqrt{x+2}\le x$ et $\sqrt{x+2}>x$ dans $\R$.
\begin{corrige}[exercice~\ref{t11s-c08-e05}]\label{sol-t11s-c08-e05}
Pour la première : $x\ge0$ et $x+2\le x^2$, donc $x\ge2$. Pour la seconde : domaine $x\ge-2$; si $x<0$ tout convient; si $x\ge0$, $x+2>x^2$ donne $0\le x<2$. Réponse $[-2;2[$.
\end{corrige}
\exo[Deux distances]\label{t11s-c08-e06}
Résoudre $|x-1|+|x+2|=5$.
\begin{corrige}[exercice~\ref{t11s-c08-e06}]\label{sol-t11s-c08-e06}
Sur $x<-2$, somme $-2x-1$, d’où $x=-3$. Sur $[-2;1]$, somme $3$, aucune solution. Sur $x>1$, somme $2x+1$, d’où $x=2$. Solutions $\{-3;2\}$.
\end{corrige}
\exo[Quotient]\label{t11s-c08-e07}
Résoudre $\dfrac{x-2}{x+1}\ge0$.
\begin{corrige}[exercice~\ref{t11s-c08-e07}]\label{sol-t11s-c08-e07}
Domaine $x\ne-1$. Positif sur $]-\infty;-1[$ et $]2;+\infty[$; nul en $2$. Réponse $]-\infty;-1[\cup[2;+\infty[$.
\end{corrige}
\exo[Une vérification indispensable]\label{t11s-c08-e08}
Un élève résout $\sqrt{2x+3}=x-1$ et conserve toutes les racines de $x^2-4x-2=0$. Réparer.
\begin{corrige}[exercice~\ref{t11s-c08-e08}]\label{sol-t11s-c08-e08}
Il faut $x\ge1$. Les racines au carré sont $2\pm\sqrt6$; seule $2+\sqrt6$ respecte la condition. Substituer confirme la solution; les équivalences avec les conditions prouvent l’exhaustivité.
\end{corrige}
\exo[Inégalités de distances]\label{t11s-c08-e09}
Résoudre $|2x-3|\le5$, puis $|x+1|>2$, dans $\R$. Représenter les ensembles de solutions sur une droite graduée.
\begin{corrige}[exercice~\ref{t11s-c08-e09}]\label{sol-t11s-c08-e09}
La première équivaut à $-5\le2x-3\le5$, donc $-2\le2x\le8$ et $-1\le x\le4$. La solution est $[-1;4]$, avec les deux bornes incluses. La seconde équivaut à $x+1<-2$ ou $x+1>2$, donc $x<-3$ ou $x>1$. La solution est $]-\infty;-3[\cup]1;+\infty[$, avec $-3$ et $1$ exclus. Sur la droite, colorer le segment fermé dans le premier cas, les deux demi-droites ouvertes dans le second.
\end{corrige}

\exo[Des racines doivent respecter une contrainte]\label{t11s-c08-p01}
Pour $m\in\R$, résoudre $x^2-2mx+m^2-4=0$. Déterminer les valeurs de $m$ pour lesquelles les deux solutions sont strictement positives, puis celles pour lesquelles exactement une solution est strictement positive. Traiter les valeurs frontières.
\begin{corrige}[exercice~\ref{t11s-c08-p01}]\label{sol-t11s-c08-p01}
L’équation est $(x-m)^2=4$, donc deux racines distinctes $m-2$ et $m+2$. Elles sont toutes deux positives si $m>2$. Exactement une est positive si $m-2\le0<m+2$, soit $-2<m\le2$. Pour $m=-2$, les racines sont $-4,0$ : aucune strictement positive. Pour $m=2$, elles sont $0,4$ : exactement une. Pour $m<-2$, aucune n’est positive. Les valeurs frontières dépendent du mot « strictement ».
\end{corrige}
\exo[Une simplification et un trou]\label{t11s-c08-p02}
Résoudre $\dfrac{(x-2)(x+1)}{x-2}\le0$ dans $\R$, en conservant le domaine initial. Un élève annonce $x\le-1$ puis ajoute $x=2$ « car le numérateur est nul ». Expliquer l’erreur et dire si sa première partie est correcte.
\begin{corrige}[exercice~\ref{t11s-c08-p02}]\label{sol-t11s-c08-p02}
Le domaine exclut $2$. Sur ce domaine, le quotient vaut $x+1$, donc l’inéquation donne $x\le-1$. Cette demi-droite ne contient pas $2$ et est bien la solution complète. En $2$, le quotient est indéfini; un zéro du numérateur ne suffit pas lorsque le dénominateur est aussi nul. Il ne faut pas ajouter cette valeur.
\end{corrige}

\begin{jemeteste}Je traite les pertes de degré, je vérifie toute racine annoncée, je sépare les signes avant de mettre au carré.\end{jemeteste}
\begin{notepourlaclasse}Le troisième exercice est volontairement une détection d’erreur : le présenter comme tel pour éviter une injonction impossible. Les cubiques et équations irrationnelles avec paramètre ne sont pas requises.\end{notepourlaclasse}

\clearpage\bmtchapitre{Systèmes linéaires et pivot de Gauss}{Résoudre un système jusqu’à quatre inconnues et interpréter l’unicité, l’incompatibilité ou une famille de solutions.}{Algèbre}
\label{t11s-c09}
\begin{revision}Équations linéaires à deux inconnues; substitution; opérations sur les égalités.\end{revision}
\begin{automatismes}\exo\label{t11s-c09-auto}Résoudre $2x=6$ et $x+y=5$ avec $x=3$. Une égalité $0=1$ peut-elle être vraie ?\end{automatismes}
\begin{corrige}[exercice~\ref{t11s-c09-auto}]\label{sol-t11s-c09-auto}$x=3$, $y=2$; $0=1$ est impossible.\end{corrige}

\section{Transformer un système sans changer ses solutions}
\begin{definition}Une solution d’un système est un tuple qui vérifie toutes les équations simultanément. Les coefficients sont fixés; les inconnues désignent les quantités à déterminer.\end{definition}
\begin{propriete}[opérations élémentaires]On peut échanger deux équations, multiplier une équation par un réel non nul, ou ajouter à une équation un multiple d’une autre. Ces opérations conservent exactement les solutions.\end{propriete}
\begin{demonstration}Chaque opération est réversible : échanger de nouveau, diviser par le coefficient non nul, ou soustraire le même multiple. Les systèmes avant et après sont donc équivalents.\end{demonstration}
\begin{methode}{éliminer puis remonter}Choisir un coefficient non nul comme pivot, échanger des lignes si nécessaire, éliminer cette inconnue dans les lignes suivantes; recommencer avec les inconnues restantes. Quand le système est triangulaire, résoudre la dernière équation puis remonter. Ne jamais diviser par un pivot nul.\end{methode}
\begin{exemple}[quatre inconnues]
\[
\begin{cases}
x+y+z+t=10,\\x-y+z+t=6,\\x+y-z+t=4,\\x+y+z-t=2.
\end{cases}
\]
Soustraire la première ligne aux trois autres donne $-2y=-4$, $-2z=-6$, $-2t=-8$. Ainsi $y=2,z=3,t=4$, puis $x=1$. Vérification : les quatre membres gauches valent $10,6,4,2$.
\end{exemple}
\section{Reconnaître les trois issues}
\begin{exemple}[incompatibilité et liberté]
Le système $x+y=2$, $2x+2y=5$ donne après élimination $0=1$ : aucune solution. Si la seconde équation est $2x+2y=4$, elle devient $0=0$, donc n’ajoute aucune contrainte; les solutions sont $(2-s;s)$ pour $s\in\R$.
\end{exemple}
\begin{avertissement}Une ligne $0=0$ ne signifie pas que toutes les inconnues sont nulles. Une ligne impossible rend tout le système incompatible. Une famille à paramètre doit conserver les autres équations.\end{avertissement}
\section{Modéliser}
Pour un problème concret, nommer les inconnues avec leurs unités, traduire chaque donnée en équation et contrôler ensuite les contraintes (positivité, intégralité, quantités disponibles). Une solution réelle algébrique peut ne pas convenir au contexte.

% ENRICHISSEMENT C09

\subsection*{Voir toutes les étapes du pivot}
\begin{exemple}[un système triangulaire à trois inconnues]
Pour $x+y+z=6$, $2x-y+z=3$, $x+2y-z=2$, utiliser la première ligne comme pivot. Les opérations $L_2\leftarrow L_2-2L_1$ et $L_3\leftarrow L_3-L_1$ donnent
\[
\begin{cases}x+y+z=6,\\-3y-z=-9,\\y-2z=-4.\end{cases}
\]
Puis $L_3\leftarrow3L_3+L_2$ donne $-7z=-21$. On remonte : $z=3$, $y=2$, $x=1$. L’opération combinée signifie multiplier d’abord la troisième ligne par $3\ne0$, puis lui ajouter la deuxième; elle est réversible. La substitution finale contrôle les trois données originales.
\end{exemple}
\begin{methode}{présenter une famille sans perdre les contraintes}
Si une inconnue reste libre, la nommer par un paramètre $s$ et exprimer toutes les autres en fonction de $s$. Indiquer son domaine. Dans une situation de prix ou d’effectifs, ajouter ensuite les contraintes de positivité ou d’intégralité; elles peuvent restreindre la famille obtenue dans $\R$.
\end{methode}
\begin{exemple}[des prix positifs dans une famille]
Si deux prix, en milliers d’ariary, vérifient seulement $a+b=12$, les solutions réelles sont $(12-s;s)$. Pour deux prix strictement positifs, on impose $0<s<12$. Une seule équation ne détermine pas les deux prix, même si le total est connu. Ajouter la différence $a-b=4$ impose $s=4$ et donc $(a;b)=(8;4)$.
\end{exemple}
\begin{remarque}[contrôle de l’unicité]
Un tuple vérifié est une solution, mais cette vérification seule n’exclut pas d’autres solutions. L’équivalence des systèmes et la remontée imposant chaque inconnue établissent l’unicité. Une égalité $0=0$ doit être interprétée; ce n’est pas un nouveau pivot.
\end{remarque}

\section{Exercices et problèmes}
\exo[Éliminer]\label{t11s-c09-e01}
Résoudre $x+y+z=6$, $2x-y+z=3$, $x+2y-z=2$.
\begin{corrige}[exercice~\ref{t11s-c09-e01}]\label{sol-t11s-c09-e01}
De la première, $z=6-x-y$. La seconde donne $x-2y=-3$, la troisième $2x+3y=8$. D’où $y=2,x=1,z=3$. Substitution vérifie $6,3,2$.
\end{corrige}
\exo[Quatre variables]\label{t11s-c09-e02}
Résoudre $x+y=3$, $y+z=5$, $z+t=7$, $x+y+z+t=10$.
\begin{corrige}[exercice~\ref{t11s-c09-e02}]\label{sol-t11s-c09-e02}
$x=3-y$, $z=5-y$, $t=2+y$; la dernière devient $10=10$. Famille $(3-s;s;5-s;2+s)$, $s\in\R$, aucune unicité.
\end{corrige}
\exo[Incompatible]\label{t11s-c09-e03}
Résoudre $x+2y=4$, $2x+4y=7$.
\begin{corrige}[exercice~\ref{t11s-c09-e03}]\label{sol-t11s-c09-e03}
La deuxième moins deux fois la première donne $0=-1$. Aucune solution.
\end{corrige}
\exo[Point de vente]\label{t11s-c09-e04}
Une boutique vend quatre lots A, B, C, D. Acheter A et B coûte $9000$ ariary, B et C coûte $11000$ ariary, C et D coûte $13000$ ariary, A et deux lots D coûte $20000$ ariary. Nommer leurs prix $a,b,c,d$ en milliers d’ariary, traduire les données par un système puis les déterminer.
\begin{corrige}[exercice~\ref{t11s-c09-e04}]\label{sol-t11s-c09-e04}
Le modèle est $a+b=9$, $b+c=11$, $c+d=13$, $a+2d=20$. Donc $a=9-b$, $c=11-b$, $d=2+b$. La dernière : $9-b+4+2b=20$, donc $b=7$, $a=2$, $c=4$, $d=9$. Prix $2000,7000,4000,9000$ ariary, tous positifs.
\end{corrige}
\exo[Pivot nul]\label{t11s-c09-e05}
Résoudre $y+z=3$, $x+y=4$, $x+z=5$. Expliquer pourquoi il faut choisir une autre ligne si l’on souhaite un premier pivot en $x$.
\begin{corrige}[exercice~\ref{t11s-c09-e05}]\label{sol-t11s-c09-e05}
La première ligne a coefficient nul en $x$; échanger avec la deuxième. Ajouter les deux dernières relations puis utiliser la première donne $2x=6$, donc $x=3,y=1,z=2$.
\end{corrige}
\exo[Un paramètre]\label{t11s-c09-e06}
Étudier selon $m$ le système $x+y=2$, $x+my=3$.
\begin{corrige}[exercice~\ref{t11s-c09-e06}]\label{sol-t11s-c09-e06}
Soustraction : $(m-1)y=1$. Si $m=1$, impossible. Sinon $y=1/(m-1)$ et $x=(2m-3)/(m-1)$, solution unique.
\end{corrige}

\exo[Un modèle d’effectifs avec contrainte entière]\label{t11s-c09-p01}
Un club compte $30$ membres répartis entre élèves et adultes. La cotisation est de $2\,000$ ariary par élève et $5\,000$ ariary par adulte; le total reçu est $90\,000$ ariary. Définir deux inconnues, écrire et résoudre le système. Vérifier intégralité, positivité et les deux données. Le même club pourrait-il recevoir exactement $91\,000$ ariary avec ces tarifs et ces $30$ membres ?
\begin{corrige}[exercice~\ref{t11s-c09-p01}]\label{sol-t11s-c09-p01}
Si $e,a$ sont les effectifs, $e+a=30$ et $2e+5a=90$, avec montants en milliers d’ariary. Soustraire deux fois la première donne $3a=30$, donc $a=10,e=20$. Ils sont entiers positifs, leur somme est $30$ et $2\times20+5\times10=90$. Avec total $91$, on aurait $3a=31$, soit $a=31/3$, impossible pour un effectif. Une solution réelle ne suffit pas au contexte.
\end{corrige}
\exo[Discuter un système selon le paramètre]\label{t11s-c09-p02}
Étudier dans $\R^2$, selon $m\in\R$, le système $x+y=2$, $2x+my=m$. Donner tous les couples solutions, y compris le cas exceptionnel; vérifier ce cas dans les deux équations.
\begin{corrige}[exercice~\ref{t11s-c09-p02}]\label{sol-t11s-c09-p02}
Éliminer $x$ donne $(m-2)y=m-4$. Pour $m=2$, l’égalité $0=-2$ est impossible, donc aucune solution. Pour $m\ne2$, $y=(m-4)/(m-2)$ et $x=m/(m-2)$, solution unique. En $m=2$, les équations imposeraient simultanément $x+y=2$ et $x+y=1$.
\end{corrige}

\begin{jemeteste}Je peux expliquer pourquoi chaque opération conserve les solutions et vérifier le tuple dans toutes les équations originales.\end{jemeteste}
\begin{notepourlaclasse}Faire exécuter les opérations sur lignes avec leur notation explicite. La famille de solutions est importante : quatre équations ne garantissent pas quatre pivots ni l’unicité.\end{notepourlaclasse}

\clearpage\bmtchapitre{Divisibilité, nombres premiers, PGCD et PPCM}{Décomposer un entier en facteurs premiers, utiliser Euclide et résoudre des problèmes de regroupement et de périodicité.}{Algèbre}
\label{t11s-c10}
\begin{revision}Division euclidienne, multiples et diviseurs positifs; boucles du chapitre algorithmique.\end{revision}
\begin{automatismes}\exo\label{t11s-c10-auto}Écrire $47=6q+r$ avec $0\le r<6$. Donner les diviseurs positifs de $12$.\end{automatismes}
\begin{corrige}[exercice~\ref{t11s-c10-auto}]\label{sol-t11s-c10-auto}$q=7,r=5$; $1,2,3,4,6,12$.\end{corrige}

\section{Diviseurs et nombres premiers}
\begin{definition}Pour $a,b\in\Z$, $a\mid b$ signifie qu’il existe $k\in\Z$ avec $b=ak$. Un entier naturel $p\ge2$ est premier s’il a exactement deux diviseurs positifs : $1$ et $p$. Le nombre $1$ n’est pas premier.\end{definition}
\begin{propriete}Si $a\mid b$ et $a\mid c$, alors $a\mid (ub+vc)$ pour tous entiers $u,v$. Si $n\ge2$ est composé, il possède un diviseur premier inférieur ou égal à $\sqrt n$.\end{propriete}
\begin{demonstration}Écrire $b=ak,c=a\ell$ donne $ub+vc=a(uk+v\ell)$. Si $n=de$ avec $1<d\le e<n$, alors $d^2\le n$; $d$ possède un diviseur premier (en prenant son plus petit diviseur supérieur à $1$), lui-même au plus $\sqrt n$.\end{demonstration}
\begin{theoreme}[décomposition, admis]Tout entier $n\ge2$ possède une décomposition en produit de nombres premiers, unique à l’ordre des facteurs près.\end{theoreme}
\begin{exemple}$360=2^3\times3^2\times5$. Pour tester $97$, il suffit de vérifier les nombres premiers $2,3,5,7$, car $\sqrt{97}<10$. Aucun ne divise $97$, donc $97$ est premier.\end{exemple}
\section{PGCD et algorithme d’Euclide}
\begin{definition}Pour $a,b\in\N^*$, $\pgcd(a,b)$ est le plus grand diviseur positif commun. Ils sont premiers entre eux lorsque ce PGCD vaut $1$. On étend par $\pgcd(a,0)=\pgcd(0,a)=a$ pour $a>0$; $\pgcd(0,0)$ n’est pas utilisé ici.\end{definition}
\begin{propriete}Si $a=bq+r$, les diviseurs communs de $a,b$ sont ceux de $b,r$. Ainsi $\pgcd(a,b)=\pgcd(b,r)$.\end{propriete}
\begin{demonstration}Un diviseur de $a,b$ divise $r=a-bq$. Un diviseur de $b,r$ divise $a=bq+r$. Les ensembles sont les mêmes.\end{demonstration}
\begin{exemple}$252=198+54$; $198=3\times54+36$; $54=36+18$; $36=2\times18$. Le dernier reste non nul est $18$, donc $\pgcd(252,198)=18$.\end{exemple}
\begin{methode}{Euclide en boucle}Entrées $a,b\in\N$, non tous deux nuls. Tant que $b\ne0$, calculer $r\leftarrow a\bmod b$, puis $a\leftarrow b$, $b\leftarrow r$. Afficher $a$. Les restes non nuls décroissent strictement, donc la boucle termine. Les diviseurs communs sont conservés; la sortie est le PGCD.\end{methode}
\section{PPCM}
\begin{definition}Pour $a,b\in\N^*$, $\ppcm(a,b)$ est le plus petit multiple positif commun.\end{definition}
\begin{propriete}Dans les décompositions premières, le PGCD prend les exposants minimaux et le PPCM les exposants maximaux (un facteur absent a exposant $0$). Ainsi $\pgcd(a,b)\ppcm(a,b)=ab$.\end{propriete}
\begin{demonstration}Un diviseur commun ne peut avoir un exposant supérieur au minimum; le produit des minima convient. Un multiple commun doit avoir au moins les maxima; leur produit est le plus petit. Pour chaque premier, minimum plus maximum égale la somme des deux exposants.\end{demonstration}
\begin{exemple}$72=2^3 3^2$, $90=2\,3^2 5$ : PGCD $18$, PPCM $360$. Pour trois entiers, calculer successivement le PGCD ou PPCM de deux, puis avec le troisième.\end{exemple}

\section{Exercices et problèmes}
\exo[Premier ou composé]\label{t11s-c10-e01}
Déterminer si $91$ et $101$ sont premiers. Justifier les tests suffisants.
\begin{corrige}[exercice~\ref{t11s-c10-e01}]\label{sol-t11s-c10-e01}
$91=7\times13$, composé. $\sqrt{101}<11$ : tester $2,3,5,7$, aucun ne divise $101$, premier.
\end{corrige}
\exo[Décompositions]\label{t11s-c10-e02}
Décomposer $180$ et $252$, puis calculer leur PGCD et PPCM.
\begin{corrige}[exercice~\ref{t11s-c10-e02}]\label{sol-t11s-c10-e02}
$180=2^2 3^2 5$, $252=2^2 3^2 7$; PGCD $36$, PPCM $1260$. Produit $36\times1260=180\times252$.
\end{corrige}
\exo[Euclide]\label{t11s-c10-e03}
Calculer le PGCD de $1071$ et $462$ par Euclide, puis simplifier $1071/462$.
\begin{corrige}[exercice~\ref{t11s-c10-e03}]\label{sol-t11s-c10-e03}
$1071=2\times462+147$; $462=3\times147+21$; $147=7\times21$. PGCD $21$, fraction $51/22$.
\end{corrige}
\exo[Lots identiques]\label{t11s-c10-e04}
On dispose de $84$ cahiers et $126$ stylos. Former le plus grand nombre de lots identiques sans reste, avec au moins un objet de chaque sorte. Donner leur composition.
\begin{corrige}[exercice~\ref{t11s-c10-e04}]\label{sol-t11s-c10-e04}
Le nombre de lots divise les deux effectifs; son maximum est $\pgcd(84,126)=42$. Chaque lot contient $2$ cahiers et $3$ stylos.
\end{corrige}
\exo[Deux cycles]\label{t11s-c10-e05}
Deux opérations ont lieu tous les $12$ et $18$ jours. Elles ont lieu ensemble au jour $0$. Quand coïncident-elles ensuite pour la première fois ?
\begin{corrige}[exercice~\ref{t11s-c10-e05}]\label{sol-t11s-c10-e05}
Le premier jour positif doit être multiple des deux durées : PPCM $36$ jours. Les coïncidences sont aux jours $36k$.
\end{corrige}
\exo[Premiers entre eux]\label{t11s-c10-e06}
Deux entiers premiers entre eux sont-ils nécessairement premiers ? Deux nombres premiers distincts sont-ils premiers entre eux ?
\begin{corrige}[exercice~\ref{t11s-c10-e06}]\label{sol-t11s-c10-e06}
Première affirmation fausse : $8$ et $9$ ont PGCD $1$ et sont composés. Seconde vraie : un diviseur commun supérieur à $1$ de deux premiers imposerait leur égalité.
\end{corrige}
\exo[Trois entiers]\label{t11s-c10-e07}
Calculer le PGCD et le PPCM de $24,36,54$. Peut-on utiliser directement la formule PGCD fois PPCM égale le produit des trois nombres ?
\begin{corrige}[exercice~\ref{t11s-c10-e07}]\label{sol-t11s-c10-e07}
PGCD $6$; PPCM $2^3 3^3=216$. $6\times216=1296\ne24\times36\times54$ : la formule donnée concerne deux entiers, pas trois.
\end{corrige}
\begin{jemeteste}Je distingue diviseur et multiple; un regroupement maximal demande souvent le PGCD, un retour commun minimal le PPCM.\end{jemeteste}
\begin{notepourlaclasse}Faire relier Euclide à l’organigramme et justifier l’arrêt par la décroissance des restes. Bézout et congruences générales ne sont pas requis ici.\end{notepourlaclasse}

\clearpage\bmtchapitre{Systèmes de numération}{Lire une écriture en base entière et convertir un entier naturel d’une base à une autre.}{Algèbre}
\label{t11s-c11}
\begin{revision}Puissances, division euclidienne et somme; chiffres du système décimal.\end{revision}
\begin{automatismes}\exo\label{t11s-c11-auto}Calculer $2^5$, $3^3$, et le quotient et le reste de $23$ par $2$.\end{automatismes}
\begin{corrige}[exercice~\ref{t11s-c11-auto}]\label{sol-t11s-c11-auto}$32$, $27$; quotient $11$, reste $1$.\end{corrige}

\section{La valeur dépend de la position}
\begin{definition}Dans une base entière $b\ge2$, les chiffres sont $0,1,\ldots,b-1$. L’écriture $(a_k\cdots a_1a_0)_b$ représente $\sum_{i=0}^k a_i b^i$. Pour un entier non nul, le premier chiffre n’est pas nul; zéro s’écrit $(0)_b$. Pour les bases supérieures à dix, on peut convenir de lettres $A=10$, $B=11$, etc.\end{definition}
\begin{exemple}$(1011)_2=1\times8+0\times4+1\times2+1=11$. $(203)_4=2\times16+0\times4+3=35$. Dans la base $4$, le chiffre $4$ est interdit.\end{exemple}
\section{Convertir par divisions successives}
\begin{methode}{de la base dix à la base $b$}Diviser l’entier par $b$, noter le reste, recommencer avec le quotient jusqu’à quotient nul. Lire les restes du dernier au premier.\end{methode}
\begin{exemple}Pour $45$ en base $2$ :
\begin{center}\begin{tabular}{r|rr}\toprule
Nombre & Quotient & Reste\\\midrule
45&22&1\\22&11&0\\11&5&1\\5&2&1\\2&1&0\\1&0&1\\\bottomrule
\end{tabular}\end{center}
Donc $45=(101101)_2$. Vérification : $32+8+4+1=45$.
\end{exemple}
\begin{propriete}[existence et unicité]Chaque entier naturel possède une écriture unique en base $b\ge2$ sans zéro initial, sauf l’écriture de zéro.\end{propriete}
\begin{demonstration}La première division euclidienne impose le dernier chiffre, qui est l’unique reste dans $[0;b-1]$. Le quotient impose les suivants de la même manière. Pour un entier positif, les quotients successifs décroissent jusqu’à zéro. Les divisions fournissent donc une écriture finie et forcent tous ses chiffres.\end{demonstration}
\section{Changer de base et calculer}
Passer par la base dix évite de confondre les valeurs des chiffres. On peut aussi additionner directement dans la base choisie, avec une retenue dès que la somme atteint la base.
\begin{exemple}En binaire, $(101)_2+(11)_2=(1000)_2$, car $5+3=8$. Une retenue binaire correspond à deux unités du rang précédent, et non à dix.\end{exemple}

\section{Exercices et problèmes}
\exo[Vers dix]\label{t11s-c11-e01}
Convertir $(11001)_2$, $(2102)_3$ et $(2A)_{16}$ en base dix.
\begin{corrige}[exercice~\ref{t11s-c11-e01}]\label{sol-t11s-c11-e01}
$16+8+1=25$; $2\times27+1\times9+2=65$; $2\times16+10=42$.
\end{corrige}
\exo[Depuis dix]\label{t11s-c11-e02}
Écrire $53$ en base $2$ et en base $5$, puis vérifier.
\begin{corrige}[exercice~\ref{t11s-c11-e02}]\label{sol-t11s-c11-e02}
$53=32+16+4+1$, donc $(110101)_2$; $53=2\times25+0\times5+3$, donc $(203)_5$.
\end{corrige}
\exo[Chiffre interdit]\label{t11s-c11-e03}
L’écriture $(102)_2$ est-elle valide ? Quelle est la plus petite base autorisant $(274)_b$ ?
\begin{corrige}[exercice~\ref{t11s-c11-e03}]\label{sol-t11s-c11-e03}
Non, le chiffre $2$ est interdit en base $2$. La seconde exige $b>7$ : plus petite base $8$.
\end{corrige}
\exo[Chercher une base]\label{t11s-c11-e04}
Trouver l’entier $b$ tel que $(23)_b=17$ en base dix. Vérifier la validité des chiffres.
\begin{corrige}[exercice~\ref{t11s-c11-e04}]\label{sol-t11s-c11-e04}
$2b+3=17$, donc $b=7$. Les chiffres $2$ et $3$ sont inférieurs à $7$, écriture valide.
\end{corrige}
\exo[Addition binaire]\label{t11s-c11-e05}
Calculer $(1011)_2+(110)_2$ directement en base $2$, puis contrôler en base dix.
\begin{corrige}[exercice~\ref{t11s-c11-e05}]\label{sol-t11s-c11-e05}
Résultat $(10001)_2$; contrôle $11+6=17$. Poser les unités l’une sous l’autre et reporter chaque retenue.
\end{corrige}
\exo[Bornes]\label{t11s-c11-e06}
Montrer qu’un entier positif écrit avec $k$ chiffres en base $b$ est compris entre $b^{k-1}$ et $b^k-1$.
\begin{corrige}[exercice~\ref{t11s-c11-e06}]\label{sol-t11s-c11-e06}
Minimum : premier chiffre $1$, autres $0$, soit $b^{k-1}$. Maximum : tous chiffres $b-1$, soit $(b-1)(1+b+\cdots+b^{k-1})=b^k-1$. La somme se vérifie en multipliant par $b$ puis soustrayant.
\end{corrige}
\begin{jemeteste}Je précise la base, je contrôle que chaque chiffre est autorisé et je lis les restes dans l’ordre inverse des divisions.\end{jemeteste}
\begin{notepourlaclasse}Utiliser des groupements concrets en paquets de deux ou cinq. Les écritures fractionnaires dans une base ne sont pas requises; traiter uniquement les entiers naturels.\end{notepourlaclasse}

\clearpage\bmtchapitre{Suites numériques}{Définir et représenter une suite, déterminer ses variations et sa limite, utiliser les suites arithmétiques et géométriques et calculer leurs sommes.}{Algèbre}
\label{t11s-c12}
\begin{revision}Fonctions usuelles, puissances, boucles et limites à l’infini.\end{revision}
\begin{automatismes}\exo\label{t11s-c12-auto}Calculer $2^4$, $1/2^3$ et $3+4+5+6$.\end{automatismes}
\begin{corrige}[exercice~\ref{t11s-c12-auto}]\label{sol-t11s-c12-auto}$16$, $1/8$, $18$.\end{corrige}

\section{Une fonction des entiers naturels}
\begin{definition}Une suite réelle $(u_n)$ est une fonction définie sur $\N$ (ou sur les entiers à partir d’un rang précisé). Une formule explicite donne $u_n$ directement; une relation de récurrence donne le terme suivant à partir des précédents et nécessite des termes initiaux. On représente les points $(n;u_n)$ : seuls les indices entiers appartiennent au domaine.\end{definition}
\begin{exemple}$u_n=n^2+1$ donne $1,2,5,10$ aux rangs $0,1,2,3$. Avec $v_0=1$ et $v_{n+1}=2v_n+1$, on obtient $1,3,7,15$. Ce sont deux modes de génération différents.\end{exemple}
\begin{definition}Une suite est croissante si $u_{n+1}\ge u_n$ pour chaque rang; décroissante avec l’inégalité opposée. Elle converge vers $\ell$ si ses termes se rapprochent de $\ell$ autant que voulu à partir d’un rang suffisamment grand. Sinon elle est divergente; une divergence peut être infinie ou oscillante.\end{definition}
\begin{methode}{étudier les variations}Calculer $u_{n+1}-u_n$. Pour des termes strictement positifs, comparer aussi $u_{n+1}/u_n$ à $1$. Un quotient n’est pas utilisable si un terme est nul; la positivité est nécessaire pour conserver le sens de l’inégalité.\end{methode}
\section{Suites arithmétiques}
\begin{definition}$u_{n+1}=u_n+r$ pour un réel constant $r$ définit une suite arithmétique de raison $r$.\end{definition}
\begin{propriete}Pour $n\ge p$, $u_n=u_p+(n-p)r$. La somme de $k\ge1$ termes à partir de $u_p$ vaut
\[
 u_p+\cdots+u_{p+k-1}=\frac{k(u_p+u_{p+k-1})}{2}.
\]
\end{propriete}
\begin{demonstration}Ajouter $r$ exactement $n-p$ fois donne le terme général, formalisable par récurrence. Pour la somme, écrire les termes dans les deux ordres et additionner : chaque paire vaut premier plus dernier; deux fois la somme vaut $k$ fois cette quantité.\end{demonstration}
Si $r>0$, la suite croît et tend vers $+\infty$; si $r<0$, elle décroît et tend vers $-\infty$; si $r=0$, elle est constante.
\section{Suites géométriques}
\begin{definition}$u_{n+1}=qu_n$, pour un réel constant $q$, définit une suite géométrique. Pour $q=0$, tous les termes après le terme initial sont nuls.\end{definition}
\begin{propriete}Pour $q\ne0$, $u_n=u_pq^{n-p}$. Pour $q=0$, cette formule vaut pour $n>p$ et $u_p$ reste le terme initial, sans utiliser $0^0$. Si $q\ne1$,
\[
 u_p+\cdots+u_{p+k-1}=u_p\frac{1-q^k}{1-q}.
\]
Si $q=1$, la somme vaut $ku_p$. La formule de somme convient aussi à $q=0$.
\end{propriete}
\begin{demonstration}Le terme général résulte de multiplications successives. Pour $q=0$, la somme de $k\ge1$ termes vaut $u_p$, ce qui coïncide avec la formule annoncée. Pour $q\ne0$, considérer la somme $S=u_p+u_pq+\cdots+u_pq^{k-1}$, soustraire $qS$ donne $(1-q)S=u_p(1-q^k)$. Séparer $q=1$ avant la division.\end{demonstration}
\begin{propriete}[limites, admises pour les puissances]
Si $|q|<1$, $q^n\to0$. Si $q>1$, $q^n\to+\infty$. Si $q=-1$, les termes alternent; si $q<-1$, leur module croît sans limite et leurs signes alternent. Ainsi, pour une suite géométrique non nulle, la convergence dépend de ces cas; pour $q=1$, la suite est constante. Si $u_0=0$, elle est nulle quel que soit $q$.
\end{propriete}
\section{Un modèle de récurrence}
\begin{exemple}[apport régulier]
Si $u_{n+1}=0{,}8u_n+10$, la valeur fixe vérifie $L=0{,}8L+10$, donc $L=50$. Poser $v_n=u_n-50$ donne $v_{n+1}=0{,}8v_n$, d’où $u_n=50+(u_0-50)(0{,}8)^n$ et $u_n\to50$.
\end{exemple}
\begin{remarque}Pour $u_{n+1}=au_n+b$, si $a\ne1$, la valeur fixe est $b/(1-a)$. Si $a=1$, la suite est arithmétique; il ne faut pas diviser par $1-a$. Ce changement de suite est une méthode d’exercice, pas une nouvelle formule à apprendre sans preuve.\end{remarque}

% ENRICHISSEMENT C12

\subsection*{Passer d’une récurrence à un modèle explicite}
\begin{methode}{compter les termes avant la somme}
De $u_p$ à $u_q$ inclus, il y a $q-p+1$ termes, et non $q-p$. Identifier la nature de la suite et ses paramètres avant de choisir la formule. Si les termes se composent d’une constante et d’une géométrique, sommer séparément ces deux parties.
\end{methode}
\begin{exemple}[les deux contributions d’une somme]
Pour $u_n=10+6(1/2)^n$, la somme $u_0+\cdots+u_3$ comporte quatre termes. La partie constante vaut $40$; la partie géométrique vaut $6(1-(1/2)^4)/(1-1/2)=45/4$. Au total $S=205/4=51{,}25$. Contrôle direct : $16+13+11{,}5+10{,}75$ donne le même résultat.
\end{exemple}
\begin{methode}{justifier le décalage et les variations}
Pour $u_{n+1}=au_n+b$ avec $a\ne1$, résoudre $L=aL+b$, puis calculer réellement $v_{n+1}=u_{n+1}-L$ pour obtenir $av_n$. Cette identité, avec le terme initial, donne la formule. Étudier ensuite $u_{n+1}-u_n$ en conservant le signe de $u_0-L$ : une raison entre $0$ et $1$ ne rend pas automatiquement la suite décroissante.
\end{methode}
\begin{exemple}[approcher la même limite par deux côtés]
Avec $u_{n+1}=u_n/2+5$, $L=10$ et $u_n=10+(u_0-10)(1/2)^n$. Si $u_0>10$, les termes sont au-dessus de $10$ et décroissent; si $u_0<10$, ils sont en dessous et croissent. Si $u_0=10$, la suite est constante. Dans les trois cas, la limite est $10$.
\end{exemple}
\begin{avertissement}[représentation et indice]
Le terme $u_0$ est le premier terme d’une suite commençant à zéro. Relier les points pour guider l’œil ne transforme pas le domaine discret en intervalle réel. Dans une question de premier rang, vérifier le rang précédent pour établir la minimalité.
\end{avertissement}

\section{Exercices et problèmes}
\exo[Générer]\label{t11s-c12-e01}
Pour $u_n=3n-2$, donner $u_0,u_5,u_{10}$ et les variations. Placer les cinq points $(n;u_n)$ pour $0\le n\le4$ dans un repère; faut-il les relier pour représenter la suite ? Pour $v_0=2$, $v_{n+1}=v_n+3$, donner le terme général.
\begin{corrige}[exercice~\ref{t11s-c12-e01}]\label{sol-t11s-c12-e01}
$-2,13,28$; différence $3>0$, strictement croissante. Les cinq points sont $(0;-2),(1;1),(2;4),(3;7),(4;10)$ : seuls les indices entiers appartiennent au domaine, donc on ne les relie pas. $v_n=2+3n$.
\begin{center}\begin{tikzpicture}\begin{axis}[width=7cm,height=5cm,axis lines=middle,xmin=-0.5,xmax=4.5,ymin=-3,ymax=11,xtick={0,1,2,3,4},xlabel={$n$},ylabel={$u_n$}]\addplot[only marks,mark=*,bmtGrenat] coordinates {(0,-2)(1,1)(2,4)(3,7)(4,10)};\end{axis}\end{tikzpicture}\end{center}
\end{corrige}
\exo[Somme arithmétique]\label{t11s-c12-e02}
Une suite arithmétique vérifie $u_3=7$ et $u_8=22$. Trouver la raison et calculer $u_3+\cdots+u_8$.
\begin{corrige}[exercice~\ref{t11s-c12-e02}]\label{sol-t11s-c12-e02}
$5r=15$, donc $r=3$. Il y a $6$ termes, somme $6(7+22)/2=87$.
\end{corrige}
\exo[Géométrique et somme]\label{t11s-c12-e03}
$u_0=81$ et $u_{n+1}=u_n/3$. Donner $u_n$, sa limite et $\sum_{n=0}^{4}u_n$.
\begin{corrige}[exercice~\ref{t11s-c12-e03}]\label{sol-t11s-c12-e03}
$u_n=81(1/3)^n\to0$. Les cinq termes sont $81,27,9,3,1$, somme $121$.
\end{corrige}
\exo[Variation rationnelle]\label{t11s-c12-e04}
Étudier les variations et la limite de $u_n=(n+1)/(n+2)$.
\begin{corrige}[exercice~\ref{t11s-c12-e04}]\label{sol-t11s-c12-e04}
$u_{n+1}-u_n=1/((n+2)(n+3))>0$. $u_n=1-1/(n+2)\to1$. Les termes restent strictement sous $1$.
\end{corrige}
\exo[Oscillation]\label{t11s-c12-e05}
Étudier $u_n=(-1)^n$ et $v_n=(-1)^n/(n+1)$. La présence de signes alternés empêche-t-elle toujours la convergence ?
\begin{corrige}[exercice~\ref{t11s-c12-e05}]\label{sol-t11s-c12-e05}
$u_{2k}=1$, $u_{2k+1}=-1$, donc divergence. $|v_n|=1/(n+1)\to0$, donc $v_n\to0$. Alternance n’exclut pas une limite nulle.
\end{corrige}
\exo[Un stock]\label{t11s-c12-e06}
Un stock de $100$ unités perd $20\%$ de son contenu chaque jour, puis reçoit $10$ unités. Poser $u_0=100$. Trouver une formule, les variations et la limite.
\begin{corrige}[exercice~\ref{t11s-c12-e06}]\label{sol-t11s-c12-e06}
$u_{n+1}=0{,}8u_n+10$, donc $u_n=50+50(0{,}8)^n$. Différence $-10(0{,}8)^n<0$, décroissante vers $50$; $u_1=90$.
\end{corrige}
\exo[Reconnaître une nature]\label{t11s-c12-e07}
Pour $u_0=1$, $u_{n+1}=2u_n+1$, montrer que la suite n’est ni arithmétique ni géométrique. Expliquer pourquoi la seule forme de la récurrence ne suffirait pas avec un autre terme initial.
\begin{corrige}[exercice~\ref{t11s-c12-e07}]\label{sol-t11s-c12-e07}
Premiers termes $1,3,7$ : différences $2,4$ non constantes, quotients $3,7/3$ non constants. Avec $u_0=-1$, la suite serait constante à $-1$, donc arithmétique (raison $0$) et géométrique (raison $1$).
\end{corrige}
\exo[Algorithme de rang]\label{t11s-c12-e08}
Écrire un algorithme calculant $u_N$ et $u_0+\cdots+u_N$ pour la suite du stock. Tester $N=2$.
\begin{corrige}[exercice~\ref{t11s-c12-e08}]\label{sol-t11s-c12-e08}
Initialiser $u=100,S=100$; pour $k=1$ à $N$, $u\leftarrow0{,}8u+10$, puis $S\leftarrow S+u$; afficher $u,S$. Pour $N=2$, termes $100,90,82$, sortie $82,272$. Ajouter après la mise à jour évite de compter deux fois $u_0$.
\end{corrige}

\exo[Épargne avec versement après intérêts]\label{t11s-c12-p01}
Un compte fictif contient $100\,000$ ariary au début de l’année $0$. Chaque année, le montant augmente de $10\%$, puis on verse $20\,000$ ariary. On note $u_n$ le montant, en milliers d’ariary, après $n$ opérations annuelles, donc $u_0=100$.
\begin{enumerate}
\item Établir la récurrence et calculer $u_1,u_2$.
\item Poser $v_n=u_n+200$, prouver qu’elle est géométrique et en déduire $u_n$.
\item Déterminer variations et limite; expliquer pourquoi l’ordre versement/intérêts compte.
\end{enumerate}
\begin{corrige}[exercice~\ref{t11s-c12-p01}]\label{sol-t11s-c12-p01}
Le modèle est $u_{n+1}=1{,}1u_n+20$, donc $u_1=130,u_2=163$. Puis $v_{n+1}=1{,}1u_n+220=1{,}1v_n$, avec $v_0=300$, d’où $u_n=300(1{,}1)^n-200$. La différence vaut $30(1{,}1)^n>0$ et la limite est $+\infty$. Verser avant les intérêts donnerait $1{,}1(u_n+20)=1{,}1u_n+22$, un autre modèle. Les données sont fictives, pas une recommandation financière.
\end{corrige}
\exo[Même récurrence, sens différents]\label{t11s-c12-p02}
Deux suites vérifient $w_{n+1}=w_n/2+5$, l’une avec $w_0=4$, l’autre avec $w_0=16$. Donner leurs formules, leurs variations et leur limite. Pour chacune, trouver le premier rang $n$ tel que $|w_n-10|<0{,}1$ et vérifier le rang précédent.
\begin{corrige}[exercice~\ref{t11s-c12-p02}]\label{sol-t11s-c12-p02}
Les formules sont $10-6/2^n$ et $10+6/2^n$. La première croît strictement, la seconde décroît strictement, toutes deux vers $10$. La distance dans les deux cas vaut $6/2^n$. La condition est $2^n>60$; $2^5=32$, $2^6=64$, donc le premier rang est $6$. Les distances aux rangs $5,6$ sont $3/16$ et $3/32$; seule la seconde est inférieure à $1/10$.
\end{corrige}

\begin{jemeteste}Je compte les termes, je distingue rang et nombre de tours, et je traite séparément les raisons $0,1,-1$.\end{jemeteste}
\begin{notepourlaclasse}La convergence arithmétique et géométrique se traite dans les exercices. Pour la récurrence affine, faire vérifier le décalage à la valeur fixe plutôt que donner une formule sans origine.\end{notepourlaclasse}

\clearpage\begingroup
\chapter*{Évaluation : algèbre}\markboth{Évaluation : algèbre}{Évaluation : algèbre}\addcontentsline{toc}{chapter}{Évaluation : algèbre}
\renewcommand{\thebmtexo}{DS2.\arabic{bmtexo}}
\renewcommand{\theHbmtexo}{DS2.\arabic{bmtexo}}
\setcounter{bmtexo}{0}
\begin{devoir}[durée indicative : 2 h — 20 points]
Calculatrice autorisée. Rédiger les justifications et indiquer les domaines de validité. Les exercices sont indépendants; à l’intérieur d’un exercice, on peut utiliser les résultats des questions précédentes. Les figures demandées doivent comporter les noms des objets et les points calculés.
\end{devoir}

\exo[Logique et équations avec conditions]\label{t11s-ds2-e01}\bareme{6 points}\par
Les inconnues et les paramètres de cet exercice sont réels.
\begin{enumerate}
\item Écrire la négation de $\forall x\in\R,\exists y\in\R,y>x$. Dire laquelle des deux propositions est vraie et justifier votre réponse (1 point).
\item On considère $(m-1)x^2-(m+1)x+2=0$, d’inconnue $x$ et de paramètre $m\in\R$.
\begin{enumerate}
\item Traiter $m=1$, puis démontrer la factorisation $(x-1)((m-1)x-2)$ (1 point).
\item Donner l’ensemble des solutions dans chaque cas et déterminer la valeur de $m$ pour laquelle les deux racines coïncident (2 points).
\end{enumerate}
\item Résoudre $\sqrt{x+6}=x$ dans $\R$. Indiquer les conditions avant de mettre au carré, puis vérifier toutes les solutions conservées (2 points).
\end{enumerate}
\begin{corrige}[exercice~\ref{t11s-ds2-e01}]\label{sol-t11s-ds2-e01}
\begin{enumerate}
\item La négation est $\exists x\in\R,\forall y\in\R,y\le x$. La proposition initiale est vraie : pour chaque $x$, prendre $y=x+1$. Il n’existe donc pas de réel supérieur ou égal à tous les réels.
\item Pour $m=1$, $-2x+2=0$, donc $S=\{1\}$. Développer le produit annoncé donne bien $(m-1)x^2-(m+1)x+2$. Si $m\ne1$, le produit nul donne $x=1$ ou $x=2/(m-1)$. Elles coïncident si $2/(m-1)=1$, soit $m=3$. Ainsi $S=\{1\}$ pour $m=1$ ou $m=3$, et $S=\{1;2/(m-1)\}$ sinon. Pour $m=3$, il s’agit d’une racine double; pour $m=1$, l’équation est du premier degré.
\item La racine exige $x\ge-6$ et son égalité à $x$ exige $x\ge0$. Sur $[0;+\infty[$, le carré est une transformation équivalente : $x+6=x^2$, soit $(x-3)(x+2)=0$. Seul $3$ satisfait $x\ge0$, et $\sqrt{3+6}=3$. Donc $S=\{3\}$; $-2$ n’est pas une solution de l’équation initiale.
\end{enumerate}
Le barème récompense la distinction perte de degré/racine double, pas seulement la liste des nombres.
\end{corrige}
\exo[Système, divisibilité et numération]\label{t11s-ds2-e02}\bareme{6 points}\par
Les trois parties sont indépendantes.
\begin{enumerate}
\item Résoudre dans $\R^4$ le système
\[
\begin{cases}x+y+z+t=10,\\x-y+z+t=6,\\x+y-z+t=4,\\x+y+z-t=2.\end{cases}
\]
Écrire les opérations d’élimination, donner le quadruplet et le vérifier dans les quatre équations (3 points).
\item Une association dispose de $840$ cahiers et $360$ stylos. Elle souhaite former le plus grand nombre de lots identiques, sans reste et avec chaque sorte d’objet dans chaque lot. Utiliser l’algorithme d’Euclide pour déterminer ce nombre et la composition d’un lot. Calculer aussi le PPCM de $840$ et $360$ (2 points).
\item Convertir l’entier décimal $45$ en base $2$ par divisions successives. Vérifier le résultat en développant les puissances de $2$ (1 point).
\end{enumerate}
\begin{corrige}[exercice~\ref{t11s-ds2-e02}]\label{sol-t11s-ds2-e02}
\begin{enumerate}
\item Avec $L_i\leftarrow L_i-L_1$ pour $i=2,3,4$, on obtient $-2y=-4$, $-2z=-6$, $-2t=-8$. Donc $y=2$, $z=3$, $t=4$, puis $x=10-2-3-4=1$. Le quadruplet $(1;2;3;4)$ donne $10,6,4,2$ dans les membres gauches. L’élimination réversible et les valeurs successivement imposées prouvent aussi l’unicité. Un point pour les opérations, un pour le tuple, un pour la vérification.
\item $840=2\times360+120$ et $360=3\times120+0$ : PGCD $120$. Tout nombre admissible de lots divise les deux effectifs; le maximum est donc $120$, avec $7$ cahiers et $3$ stylos par lot. Le PPCM vaut $840\times360/120=2520$. Un point pour Euclide et son interprétation, un pour composition/PPCM.
\item Les quotients successifs sont $22,11,5,2,1,0$; les restes sont $1,0,1,1,0,1$. Les lire en sens inverse donne $(101101)_2$. Vérification : $2^5+2^3+2^2+1=32+8+4+1=45$.
\end{enumerate}
\end{corrige}
\exo[Suite, seuil et algorithme]\label{t11s-ds2-e03}\bareme{8 points}\par
La suite réelle $(u_n)$ est définie pour $n\in\N$ par $u_0=10$ et $u_{n+1}=\frac12u_n+3$.
\begin{enumerate}
\item Calculer $u_1,u_2,u_3$. Comparer deux différences consécutives pour montrer que la suite n’est pas arithmétique (2 points).
\item Poser $v_n=u_n-6$. Démontrer que $(v_n)$ est géométrique, préciser son premier terme et sa raison, puis en déduire $u_n=6+4(1/2)^n$ (2 points).
\item À l’aide de cette formule, démontrer que $u_n>6$, étudier le signe de $u_{n+1}-u_n$ et déterminer la limite de la suite (2 points).
\item Écrire en pseudo-code un algorithme qui affiche le plus petit rang $n\in\N$ tel que $u_n<6{,}1$. Justifier son arrêt, calculer le rang affiché et vérifier que le rang précédent ne convient pas (2 points).
\end{enumerate}
\begin{corrige}[exercice~\ref{t11s-ds2-e03}]\label{sol-t11s-ds2-e03}
\begin{enumerate}
\item $u_1=8$, $u_2=7$, $u_3=13/2$. Les différences $8-10=-2$ et $7-8=-1$ sont distinctes; une raison arithmétique constante est impossible.
\item $v_{n+1}=u_{n+1}-6=\frac12u_n-3=\frac12(u_n-6)=\frac12v_n$, avec $v_0=4$. Donc $v_n=4(1/2)^n$ et $u_n=6+4(1/2)^n$.
\item Le terme ajouté à $6$ est strictement positif. La différence vaut $-2(1/2)^n<0$ : décroissance stricte. Comme $|1/2|<1$, cette puissance tend vers $0$; la limite est $6$, sans que la suite atteigne $6$.
\item Un algorithme possible est :
\begin{quote}\begin{tabular}{l}
$u\leftarrow10$, $n\leftarrow0$\\
Tant que $u\ge6{,}1$\\
\quad $u\leftarrow u/2+3$\\
\quad $n\leftarrow n+1$\\
Fin tant que\\Afficher $n$
\end{tabular}\end{quote}
Après $n$ tours, $u$ contient $u_n$. La convergence vers $6$ garantit un terme inférieur à $6{,}1$; la boucle s’arrête au premier tel rang. La condition équivaut à $2^n>40$. Puisque $32<40<64$, le rang est $6$. Vérification : $u_5=49/8=6{,}125$ ne convient pas et $u_6=97/16=6{,}0625$ convient. Un point pour pseudo-code/arrêt, un pour rang/minimalité.
\end{enumerate}
\end{corrige}
\endgroup


% ===========================================================================
\bmtpartie{bmtRavenala}{III}{Géométrie}{Plan, espace et transformations}{%
  Le plan d'abord, l'espace ensuite, et entre les deux la trigonométrie qui
  sert aux deux. Vous verrez qu'un même énoncé se traite tantôt par une
  figure, tantôt par un calcul de coordonnées, tantôt par une relation
  vectorielle : savoir choisir vaut mieux que savoir tout faire.}

\clearpage\bmtchapitre{Calcul vectoriel et produit scalaire}{Décomposer un vecteur dans une base et utiliser les expressions du produit scalaire pour démontrer des relations métriques.}{Géométrie}
\label{t11s-c13}
Dans ce chapitre, les calculs de longueurs, d’angles et d’aires à partir de coordonnées se font dans un repère orthonormé.
\begin{revision}Chasles, coordonnées, Pythagore, angles et trigonométrie; repère orthonormé pour les calculs métriques.\end{revision}
\begin{automatismes}\exo\label{t11s-c13-auto}Pour $A(1;2),B(4;6)$, calculer les coordonnées de $\vect{AB}$ et $AB$.\end{automatismes}
\begin{corrige}[exercice~\ref{t11s-c13-auto}]\label{sol-t11s-c13-auto}$\vect{AB}=(3;4)$, $AB=\sqrt{3^2+4^2}=5$.\end{corrige}

\section{Combinaisons linéaires et bases}
\begin{definition}Une combinaison linéaire de $\vecu,\vecv$ est $a\vecu+b\vecv$ avec $a,b\in\R$. Deux vecteurs non colinéaires forment une base du plan : chaque vecteur s’écrit d’une manière unique $a\vecu+b\vecv$.\end{definition}
\begin{propriete}Dans un repère, $\vecu=(x;y)$ et $\vecv=(x';y')$ sont non colinéaires si et seulement si $xy'-yx'\ne0$.\end{propriete}
\begin{demonstration}Le système $a\vecu+b\vecv=\vecw$ a deux équations. L’élimination fournit une solution unique si le déterminant $xy'-yx'$ est non nul. S’il est nul, les colonnes sont proportionnelles (ou l’un des vecteurs est nul), donc colinéaires.\end{demonstration}
\begin{exemple}$\vecu=(1;1)$, $\vecv=(1;-1)$ forment une base. Pour $\vecw=(3;1)$, $a+b=3$ et $a-b=1$, donc $\vecw=2\vecu+\vecv$. Cette base n’est pas un repère orthonormé : ses deux vecteurs ont norme $\sqrt2$.\end{exemple}
\section{Quatre expressions d’un même produit}
\begin{definition}Pour deux vecteurs non nuls, $\vecu\cdot\vecv=\norme{\vecu}\norme{\vecv}\cos\theta$, où $\theta$ est l’angle entre eux. Si l’un est nul, le produit scalaire est nul.\end{definition}
\begin{propriete}
Dans un repère orthonormé, $\vecu\cdot\vecv=xx'+yy'$. Pour $O\ne A$, avec $H$ projeté orthogonal de $B$ sur la droite $(OA)$ et $\vecu=\vect{OA},\vecv=\vect{OB}$, le produit vaut $OA$ multiplié par la longueur algébrique $OH$ orientée vers $A$. Enfin,
\[
\vecu\cdot\vecv=\frac12\big(\norme{\vecu+\vecv}^2-\norme{\vecu}^2-\norme{\vecv}^2\big).
\]
Le produit est symétrique et bilinéaire; $\vecu\cdot\vecu=\norme{\vecu}^2$; deux vecteurs sont orthogonaux si et seulement si leur produit est nul (y compris le vecteur nul).
\end{propriete}
\begin{demonstration}Décomposer dans deux directions orthonormées donne $xx'+yy'$. Développer $(x+x')^2+(y+y')^2$ donne l’identité par les normes. La projection donne $OH=OB\cos\theta$ avec son signe, donc la même expression géométrique. La bilinéarité résulte de la formule des coordonnées.\end{demonstration}
\section{Relations métriques dans un triangle}
\begin{propriete}[Al-Kashi et médiane]Dans un triangle $ABC$,
\[
BC^2=AB^2+AC^2-2AB\,AC\cos\widehat{BAC}.
\]
Si $I$ est milieu de $[BC]$, alors $AB^2+AC^2=2AI^2+BC^2/2$.
\end{propriete}
\begin{demonstration}Al-Kashi résulte du carré de $\vect{BC}=\vect{AC}-\vect{AB}$. Pour la médiane, $\vect{AB}=\vect{AI}+\vect{IB}$ et $\vect{AC}=\vect{AI}-\vect{IB}$; en additionnant les carrés, les produits croisés s’annulent.\end{demonstration}
\begin{propriete}[aire et sinus]Pour un triangle non aplati, $\mathcal A=\dfrac12 AB\,AC\sin\widehat{BAC}$ et
\[
\frac{BC}{\sin\widehat{BAC}}=
\frac{CA}{\sin\widehat{ABC}}=
\frac{AB}{\sin\widehat{ACB}}.
\]
\end{propriete}
\begin{demonstration}La hauteur relative à $AB$ vaut $AC\sin\widehat{BAC}$. Écrire l’aire avec chacun des trois angles puis diviser par les produits des côtés (tous non nuls) donne la relation des sinus.\end{demonstration}

% ENRICHISSEMENT C13

\subsection*{Choisir la bonne expression du produit scalaire}
\begin{methode}{relier angle, projection et coordonnées}
Avec des coordonnées dans un repère orthonormé, calculer d’abord les vecteurs puis leur produit. Pour un angle, diviser par les deux normes non nulles. Pour projeter $C$ sur la droite $(AB)$, $A\ne B$, poser $H=A+t\vect{AB}$; la condition $\vect{HC}\cdot\vect{AB}=0$ impose
\[
t=\frac{\vect{AC}\cdot\vect{AB}}{AB^2}.
\]
Si $t\notin[0;1]$, le pied appartient à la droite mais pas au segment; la hauteur d’un triangle peut être extérieure.
\end{methode}
\begin{exemple}[une projection calculée et vérifiée]
Pour $A(1;1)$, $B(5;3)$, $C(2;4)$, $\vect{AB}=(4;2)$, $\vect{AC}=(1;3)$. Le produit vaut $10$ et $AB^2=20$, donc $t=1/2$, $H(3;2)$. On vérifie $\vect{HC}=(-1;2)$ et son produit avec $(4;2)$ vaut $0$. La distance $CH=\sqrt5$ est la hauteur relative à $AB=2\sqrt5$, d’où l’aire $5$ unités d’aire. Le même produit donne $\cos\widehat{BAC}=10/(\sqrt{20}\sqrt{10})=1/\sqrt2$ : l’angle vaut $\pi/4$.
\end{exemple}
\begin{methode}{contrôler la formule dans une base quelconque}
Les coordonnées d’un vecteur sont les coefficients de sa décomposition dans la base choisie. Si la base n’est pas orthonormée, les produits des vecteurs de base interviennent. On peut revenir aux coordonnées dans un repère orthonormé pour éviter une formule mal appliquée. Une base sert à représenter les vecteurs; elle ne garantit pas des unités de longueur égales.
\end{methode}

\section{Exercices et problèmes}
\exo[Une base oblique]\label{t11s-c13-e01}
Écrire $(5;1)$ dans la base $\vecu=(2;1)$, $\vecv=(1;-1)$.
\begin{corrige}[exercice~\ref{t11s-c13-e01}]\label{sol-t11s-c13-e01}
$2a+b=5$, $a-b=1$; $a=2,b=1$. Déterminant $-3\ne0$, donc base et unicité.
\end{corrige}
\exo[Produit scalaire]\label{t11s-c13-e02}
Calculer le produit de $(2;-1)$ et $(3;4)$, leurs normes et le cosinus de leur angle.
\begin{corrige}[exercice~\ref{t11s-c13-e02}]\label{sol-t11s-c13-e02}
Produit $6-4=2$; normes $\sqrt5$ et $5$; cosinus $2/(5\sqrt5)$. Les deux vecteurs sont non nuls.
\end{corrige}
\exo[Triangle rectangle]\label{t11s-c13-e03}
Pour $A(0;0),B(4;0),C(1;\sqrt3)$, déterminer l’angle en $A$ puis calculer $BC$.
\begin{corrige}[exercice~\ref{t11s-c13-e03}]\label{sol-t11s-c13-e03}
$AB=4$, $AC=2$, produit $4$. Cosinus $4/(4\times2)=1/2$, donc angle $\pi/3$. $BC^2=16+4-2\times4\times2\times1/2=12$, donc $BC=2\sqrt3$.
\end{corrige}
\exo[Médiane]\label{t11s-c13-e04}
Dans un triangle, $AB=5$, $AC=7$, $BC=8$. Calculer la longueur de la médiane issue de $A$.
\begin{corrige}[exercice~\ref{t11s-c13-e04}]\label{sol-t11s-c13-e04}
$25+49=2AI^2+64/2$, donc $AI^2=21$, $AI=\sqrt{21}$. Les longueurs satisfont les inégalités triangulaires.
\end{corrige}
\exo[Une formule mal utilisée]\label{t11s-c13-e05}
Dans la base $\vecu=(1;1)$, $\vecv=(1;-1)$, peut-on calculer le produit scalaire de deux vecteurs en multipliant directement leurs coordonnées dans cette base ?
\begin{corrige}[exercice~\ref{t11s-c13-e05}]\label{sol-t11s-c13-e05}
Pas avec la formule sans coefficients : $\vecu$ a coordonnées $(1;0)$ dans cette base, ce qui donnerait $1$ pour sa norme au carré au lieu de $2$. Il faut une base orthonormée, ou utiliser les produits des vecteurs de base.
\end{corrige}
\exo[Mesurer une aire]\label{t11s-c13-e06}
Deux côtés d’une parcelle triangulaire mesurent $30$ m et $40$ m, avec angle compris $60\degre$. Calculer sa surface et son troisième côté.
\begin{corrige}[exercice~\ref{t11s-c13-e06}]\label{sol-t11s-c13-e06}
Aire $\tfrac12\times30\times40\times\sqrt3/2=300\sqrt3\ \mathrm{m}^2$. Troisième côté au carré $900+1600-1200=1300$, longueur $10\sqrt{13}$ m.
\end{corrige}

\exo[Une hauteur hors du segment]\label{t11s-c13-p01}
Dans un repère orthonormé, $A(0;0)$, $B(4;0)$, $C(-2;3)$. Déterminer le projeté orthogonal $H$ de $C$ sur la droite $(AB)$ et dire s’il appartient au segment $[AB]$. Calculer l’aire du triangle et le cosinus de l’angle en $A$; cet angle est-il aigu, droit ou obtus ?
\begin{corrige}[exercice~\ref{t11s-c13-p01}]\label{sol-t11s-c13-p01}
La droite $(AB)$ est $y=0$, donc $H=(-2;0)$. Son abscisse n’est pas dans $[0;4]$ : pied extérieur. La hauteur $CH=3$ donne l’aire $4\times3/2=6$. Le produit $\vect{AB}\cdot\vect{AC}=4(-2)=-8$, avec normes $4$ et $\sqrt{13}$, donne $\cos\widehat{BAC}=-2/\sqrt{13}<0$. L’angle du triangle est donc obtus. Une hauteur extérieure n’empêche pas la formule base fois hauteur.
\end{corrige}
\exo[Des coordonnées dans une base oblique]\label{t11s-c13-p02}
Dans un repère orthonormé de référence, $\vecu=(1;0)$ et $\vecv=(1;1)$. On définit $\vecw=\vecu+\vecv$ et $\vecz=\vecu-\vecv$. Un élève utilise leurs coordonnées $(1;1)$ et $(1;-1)$ dans la base $(\vecu,\vecv)$ pour conclure qu’ils sont orthogonaux. Contrôler que cette paire est une base, puis calculer le vrai produit scalaire et expliquer l’erreur.
\begin{corrige}[exercice~\ref{t11s-c13-p02}]\label{sol-t11s-c13-p02}
Le déterminant de $\vecu,\vecv$ est $1$, donc ils forment une base. Dans le repère orthonormé de référence, $\vecw=(2;1)$ et $\vecz=(0;-1)$ : produit $-1\ne0$, donc ils ne sont pas orthogonaux. La base donnée est oblique : $\vecu\cdot\vecv=1\ne0$. Multiplier directement les coefficients dans cette base revient à supposer, à tort, des vecteurs unitaires orthogonaux.
\end{corrige}

\begin{jemeteste}Je distingue une base quelconque d’un repère orthonormé; je choisis coordonnées, normes, angle ou projection selon les données.\end{jemeteste}
\begin{notepourlaclasse}Le produit scalaire de vecteurs de coordonnées dans une base oblique est un piège utile. Faire établir les relations métriques avec Chasles et bilinéarité.\end{notepourlaclasse}

\clearpage\bmtchapitre{Barycentres et lignes de niveau}{Construire et calculer un barycentre, regrouper des points pondérés et déterminer des lieux géométriques.}{Géométrie}
\label{t11s-c14}
Dans ce chapitre, les calculs de longueurs, d’angles et d’aires à partir de coordonnées se font dans un repère orthonormé.
\begin{revision}Chasles, produit scalaire, normes, milieux et équations de cercles.\end{revision}
\begin{automatismes}\exo\label{t11s-c14-auto}Calculer le milieu de $A(0;1)$ et $B(6;3)$. Développer $(x-2)^2$.\end{automatismes}
\begin{corrige}[exercice~\ref{t11s-c14-auto}]\label{sol-t11s-c14-auto}Milieu $(3;2)$; $x^2-4x+4$.\end{corrige}

\section{Deux, trois ou quatre points pondérés}
\begin{definition}Pour des points $A_i$ et des poids réels $\alpha_i$ dont la somme $s$ est non nulle, le barycentre $G$ est l’unique point tel que $\sum_i\alpha_i\vect{GA_i}=\vec0$.\end{definition}
\begin{propriete}[existence, coordonnées et réduction]
Pour tout point $O$, $\vect{OG}=\dfrac1s\sum_i\alpha_i\vect{OA_i}$. Pour tout $M$,
\[
\sum_i\alpha_i\vect{MA_i}=s\vect{MG}.
\]
Les coordonnées de $G$ sont les moyennes pondérées des coordonnées des points. Multiplier tous les poids par un même réel non nul ne change pas $G$.
\end{propriete}
\begin{demonstration}Par Chasles, $\vect{GA_i}=\vect{GO}+\vect{OA_i}$. L’équation du barycentre impose $s\vect{OG}=\sum\alpha_i\vect{OA_i}$ : un unique vecteur définit un unique point. La réduction résulte de $\vect{MA_i}=\vect{MG}+\vect{GA_i}$.\end{demonstration}
\begin{exemple}Pour $(A,1),(B,2)$, $\vect{AG}=\tfrac23\vect{AB}$. Avec des poids $1,-2$, la somme vaut $-1$ et $\vect{AG}=2\vect{AB}$ : le barycentre peut être à l’extérieur du segment. Si les poids sont $1,-1$ et $A\ne B$, aucun point ne satisfait la définition, car la somme vectorielle est $\vect{BA}\ne\vec0$.\end{exemple}
\begin{propriete}[associativité]Un groupe de points dont la somme des poids est non nulle peut être remplacé par son barycentre affecté de cette somme. La somme totale doit rester non nulle.\end{propriete}
\begin{demonstration}La réduction de ce groupe remplace exactement sa contribution dans la somme des vecteurs. L’équation finale ne change donc pas.\end{demonstration}
\section{Sommes de distances au carré}
\begin{propriete}[identité de réduction]Pour le barycentre $G$ et $s\ne0$,
\[
\sum_i\alpha_i MA_i^2=sMG^2+\sum_i\alpha_i GA_i^2.
\]
Les poids peuvent être négatifs; on ne doit alors pas supposer que la somme est positive.
\end{propriete}
\begin{demonstration}Développer le carré de $\vect{MA_i}=\vect{MG}+\vect{GA_i}$, multiplier par $\alpha_i$ et sommer. Le terme croisé vaut $2\vect{MG}\cdot\sum\alpha_i\vect{GA_i}=0$.\end{demonstration}
\section{Reconnaître les lieux}
\begin{propriete}[lignes de niveau du plan]
$MA=k$ donne un cercle si $k>0$, le point $A$ si $k=0$, l’ensemble vide si $k<0$.
Pour $\vecv\ne\vec0$, $\vect{MA}\cdot\vecv=0$ est la droite par $A$ perpendiculaire à $\vecv$; pour $\vecv=\vec0$, tout le plan.
Si $I$ est milieu de $[AB]$,
\[
\vect{MA}\cdot\vect{MB}=MI^2-\frac{AB^2}{4}.
\]
Le lieu $\vect{MA}\cdot\vect{MB}=k$ est donc un cercle, un point ou vide selon le signe de $k+AB^2/4$.
\end{propriete}
\begin{demonstration}Écrire $\vect{MA}=\vect{MI}+\vect{IA}$ et $\vect{MB}=\vect{MI}-\vect{IA}$. Le produit vaut $MI^2-IA^2$, et $IA=AB/2$.\end{demonstration}
\begin{exemple}[rapport de distances]
Supposons $A\ne B$ et cherchons $MA=kMB$. Pour $k<0$, aucun point; pour $k=0$, le point $A$; pour $k=1$, la médiatrice de $[AB]$. Pour $k>0$, $k\ne1$, mettre au carré est équivalent; le lieu est un cercle d’Apollonius. Son centre est le barycentre de $(A,1),(B,-k^2)$ et son rayon vaut $kAB/|1-k^2|$.
Pour $A(0;0),B(3;0),k=2$, $x^2+y^2=4((x-3)^2+y^2)$ donne $(x-4)^2+y^2=4$ : centre $(4;0)$ et rayon $2$.
\end{exemple}

% ENRICHISSEMENT C14

\subsection*{Réduire une expression avant de nommer un lieu}
\begin{methode}{une équation de niveau en trois étapes}
Calculer la somme des poids et le barycentre si elle est non nulle. Réduire ensuite l’expression à $sMG^2+c$. Résoudre l’équation pour $MG^2$; une valeur positive donne un cercle, zéro un point, négative l’ensemble vide. Si $s<0$, le sens d’une inégalité serait renversé en divisant; ne pas supposer que l’expression est une somme positive.
\end{methode}
\begin{exemple}[une somme pondérée et son minimum]
Avec $A(0;0)$, $B(6;0)$ et poids $2,1$, le barycentre est $G(2;0)$. La réduction donne
\[
2MA^2+MB^2=3MG^2+2GA^2+GB^2=3MG^2+24.
\]
Le minimum est $24$, atteint uniquement en $G$. Le niveau $12$ est vide; le niveau $24$ est $\{G\}$; le niveau $51$ est le cercle de centre $G$ et de rayon $3$. Il faut calculer la constante de réduction, pas seulement le centre.
\end{exemple}
\begin{methode}{construire en regroupant les poids}
Pour trois points, chercher un groupe dont le barycentre est facile à construire. Si ses poids ont une somme non nulle, remplacer ce groupe par son barycentre affecté de cette somme. Le barycentre final se construit alors sur une droite entre deux points. Vérifier les deux sommes avant les divisions.
\end{methode}
\begin{remarque}[somme nulle ne signifie pas expression inutilisable]
Quand les poids totalisent zéro, la formule du barycentre ne s’applique pas. L’expression initiale peut néanmoins se simplifier directement par Chasles ou en coordonnées. Son lieu n’est pas automatiquement vide; il faut étudier cette expression, sans inventer un centre.
\end{remarque}

\section{Exercices et problèmes}
\exo[Deux poids]\label{t11s-c14-e01}
Pour $A(1;2),B(7;-1)$, calculer le barycentre de $(A,2),(B,1)$.
\begin{corrige}[exercice~\ref{t11s-c14-e01}]\label{sol-t11s-c14-e01}
Somme $3\ne0$. $G=((2+7)/3;(4-1)/3)=(3;1)$.
\end{corrige}
\exo[Regrouper]\label{t11s-c14-e02}
Dans un triangle $ABC$, $I$ est milieu de $[BC]$. Montrer que le barycentre de $(A,2),(B,1),(C,1)$ est le milieu de $[AI]$.
\begin{corrige}[exercice~\ref{t11s-c14-e02}]\label{sol-t11s-c14-e02}
Remplacer $(B,1),(C,1)$ par $(I,2)$. Les deux poids $2,2$ donnent le milieu, somme totale $4$.
\end{corrige}
\exo[Quatre points]\label{t11s-c14-e03}
Pour le rectangle $A(0;0),B(4;0),C(4;2),D(0;2)$, trouver le barycentre des quatre points de poids $1$.
\begin{corrige}[exercice~\ref{t11s-c14-e03}]\label{sol-t11s-c14-e03}
Somme $4$; coordonnées $(2;1)$, centre du rectangle. Regrouper les deux diagonales donne le même résultat.
\end{corrige}
\exo[Lieu avec produit]\label{t11s-c14-e04}
Si $AB=6$, déterminer les lieux $\vect{MA}\cdot\vect{MB}=0$, puis $-9$, puis $-10$.
\begin{corrige}[exercice~\ref{t11s-c14-e04}]\label{sol-t11s-c14-e04}
Avec $I$ milieu, $MI^2=k+9$. Pour $0$, cercle de centre $I$, rayon $3$ (diamètre $AB$); pour $-9$, point $I$; pour $-10$, ensemble vide.
\end{corrige}
\exo[Somme de carrés]\label{t11s-c14-e05}
Pour $AB=4$ et $I$ milieu, déterminer le lieu $MA^2+MB^2=18$.
\begin{corrige}[exercice~\ref{t11s-c14-e05}]\label{sol-t11s-c14-e05}
Théorème de la médiane : $MA^2+MB^2=2MI^2+8$. Donc $MI^2=5$ : cercle de centre $I$, rayon $\sqrt5$.
\end{corrige}
\exo[Apollonius]\label{t11s-c14-e06}
Pour $A(0;0),B(3;0)$, déterminer le lieu $MA=\tfrac12 MB$ et donner le centre par barycentre.
\begin{corrige}[exercice~\ref{t11s-c14-e06}]\label{sol-t11s-c14-e06}
$4(x^2+y^2)=(x-3)^2+y^2$, donc $(x+1)^2+y^2=4$. Centre $(-1;0)$, rayon $2$; barycentre de $(A,1),(B,-1/4)$, somme $3/4$.
\end{corrige}
\exo[Poids nuls au total]\label{t11s-c14-e07}
Pourquoi ne peut-on appliquer la formule du barycentre aux poids $1,-1$ ? Distinguer $A\ne B$ et $A=B$.
\begin{corrige}[exercice~\ref{t11s-c14-e07}]\label{sol-t11s-c14-e07}
Division par somme nulle interdite. Si $A\ne B$, la somme est $\vect{BA}\ne0$, aucun point. Si $A=B$, elle vaut zéro pour tout point : aucune unicité, donc pas de barycentre au sens défini.
\end{corrige}
\exo[Distance et projection]\label{t11s-c14-e08}
Dans un repère orthonormé, $A(1;2)$ et $\vecv=(2;-1)$. Déterminer les lieux $MA=k$ pour $k=2$, $0$, $-1$, puis le lieu $\vect{MA}\cdot\vecv=0$. Que devient ce dernier lieu si $\vecv=\vec0$ ?
\begin{corrige}[exercice~\ref{t11s-c14-e08}]\label{sol-t11s-c14-e08}
Pour $k=2$, cercle de centre $A$ et rayon $2$; pour $k=0$, le seul point $A$; pour $k=-1$, ensemble vide, car une distance est positive ou nulle. Pour $M(x;y)$, le produit vaut $2(1-x)-(2-y)=y-2x$, donc le lieu est la droite $y=2x$, passant par $A$ et perpendiculaire à $(2;-1)$. Si $\vecv=\vec0$, le produit vaut toujours zéro et le lieu est tout le plan.
\end{corrige}

\exo[Des poids négatifs donnent un maximum]\label{t11s-c14-p01}
Dans un repère orthonormé, $A(0;0)$, $B(4;0)$ et $F(M)=MA^2-2MB^2$. Calculer le barycentre des poids $1,-2$, puis montrer que $F(M)=32-MG^2$. En déduire le maximum de $F$ et déterminer les lieux $F(M)=20$, $32$, $33$.
\begin{corrige}[exercice~\ref{t11s-c14-p01}]\label{sol-t11s-c14-p01}
La somme des poids vaut $-1$, et $G=(8;0)$. La constante est $GA^2-2GB^2=64-32=32$, donc $F=32-MG^2$. Le maximum $32$ est atteint seulement en $G$. Pour le niveau $20$, $MG^2=12$ : cercle de rayon $2\sqrt3$; pour $32$, point $G$; pour $33$, $MG^2=-1$, ensemble vide. La fonction peut être négative et n’a pas de minimum, car $MG$ peut croître sans borne.
\end{corrige}
\exo[Une différence de carrés sans barycentre]\label{t11s-c14-p02}
Avec $A(0;0)$, $B(4;0)$ et $M(x;y)$ dans un repère orthonormé, calculer $MA^2-MB^2$. Pour tout réel $k$, déterminer le lieu $MA^2-MB^2=k$. Retrouver la médiatrice pour $k=0$ et expliquer pourquoi les poids $1,-1$ ne définissent pas de barycentre.
\begin{corrige}[exercice~\ref{t11s-c14-p02}]\label{sol-t11s-c14-p02}
La différence vaut $x^2+y^2-[(x-4)^2+y^2]=8x-16$. Le lieu est donc la droite $x=(k+16)/8$, pour chaque $k\in\R$. Pour $k=0$, on obtient $x=2$, médiatrice de $[AB]$. Les poids ont somme nulle et $A\ne B$, donc pas de barycentre; cette absence n’empêche pas les lieux d’exister.
\end{corrige}

\begin{jemeteste}Je vérifie la somme des poids avant de diviser et je classe un lieu par le signe de son rayon au carré.\end{jemeteste}
\begin{notepourlaclasse}Construire d’abord le barycentre de poids positifs, puis faire apparaître un poids négatif. Les cas cercle/point/vide doivent figurer dans les corrections, sans ajouter une figure fictive pour un lieu vide.\end{notepourlaclasse}

\clearpage\bmtchapitre{Transformations et similitudes directes}{Construire les images de points et de figures, utiliser les expressions analytiques et composer deux transformations.}{Géométrie}
\label{t11s-c15}
\begin{revision}Vecteurs, milieux, angles orientés, coordonnées et formules trigonométriques.\end{revision}
\begin{automatismes}\exo\label{t11s-c15-auto}Donner le symétrique de $(2;-3)$ par rapport à l’origine et celui par rapport à l’axe des abscisses.\end{automatismes}
\begin{corrige}[exercice~\ref{t11s-c15-auto}]\label{sol-t11s-c15-auto}$(-2;3)$ et $(2;3)$.\end{corrige}

\section{Translations et symétries}
\begin{definition}La translation de vecteur $\vecu$ envoie $M$ sur $M'$ tel que $\vect{MM'}=\vecu$. La symétrie centrale de centre $O$ fait de $O$ le milieu de $[MM']$. La réflexion d’axe $d$ échange les deux côtés de $d$ et conserve les points de $d$.\end{definition}
\begin{propriete}[repère orthonormé]
Pour une translation de vecteur $(a;b)$, $(x',y')=(x+a,y+b)$. Pour la symétrie centrale de centre $(a;b)$, $(x',y')=(2a-x,2b-y)$. La réflexion d’axe $y=0$ donne $(x,-y)$; d’axe $x=0$, $(-x,y)$; d’axe $y=x$, $(y,x)$.
\end{propriete}
\begin{demonstration}La translation résulte des coordonnées du vecteur; la symétrie centrale résulte de la formule du milieu. Les réflexions se vérifient par perpendicularité et milieu sur l’axe.\end{demonstration}
\section{Homothéties et rotations}
\begin{definition}Une homothétie de centre $O$ et rapport $k\ne0$ envoie $M$ sur $M'$ avec $\vect{OM'}=k\vect{OM}$. Une rotation de centre $O$ et angle $\theta$ conserve $OM$ et ajoute $\theta$ à la direction du vecteur (le centre reste fixe).\end{definition}
\begin{propriete}Pour $O(a;b)$, l’homothétie donne $x'=a+k(x-a)$, $y'=b+k(y-b)$. Dans un repère orthonormé direct, la rotation donne
\[
\begin{cases}x'=a+\cos\theta(x-a)-\sin\theta(y-b),\\
y'=b+\sin\theta(x-a)+\cos\theta(y-b).
\end{cases}
\]
Translations, rotations et symétries conservent les distances. L’homothétie multiplie les distances par $|k|$ et les aires par $k^2$.
\end{propriete}
\begin{demonstration}Les images des deux vecteurs de base sous une rotation ont coordonnées $(\cos\theta;\sin\theta)$ et $(-\sin\theta;\cos\theta)$. Par linéarité, on obtient les formules; l’identité $\cos^2\theta+\sin^2\theta=1$ conserve le carré des normes. Pour l’homothétie, la norme de $k\vecu$ est $|k|\norme{\vecu}$ et chaque longueur d’une aire est multipliée par $|k|$.\end{demonstration}
\section{Composer : l’ordre compte}
\begin{definition}$g\circ f$ applique d’abord $f$, puis $g$. La composée de deux translations est la translation de la somme des vecteurs. Deux rotations de même centre se composent en additionnant les angles; deux homothéties de même centre multiplient les rapports.\end{definition}
\begin{exemple}Si $t(x,y)=(x+1,y)$ et $r(x,y)=(-y,x)$, alors $(r\circ t)(x,y)=(-y,x+1)$ mais $(t\circ r)(x,y)=(1-y,x)$. Les deux transformations ne commutent pas.\end{exemple}
\begin{propriete}La composée de deux symétries centrales, de centres $A$ puis $B$, est la translation de vecteur $2\vect{AB}$.\end{propriete}
\begin{demonstration}Dans un repère, $s_A(M)=2A-M$, puis $s_B(s_A(M))=2B-(2A-M)=M+2(B-A)$.\end{demonstration}
\section{Similitudes planes directes}
\begin{definition}Une similitude directe est une transformation obtenue à partir d’une translation et d’une rotation suivie d’une homothétie de rapport strictement positif. Elle multiplie les distances par un rapport $\rho>0$ et conserve les angles orientés. Dans un repère orthonormé direct,
\[
\begin{cases}x'=\alpha x-\beta y+p,\\y'=\beta x+\alpha y+q,\end{cases}
\qquad \alpha^2+\beta^2>0.
\]
Son rapport est $\rho=\sqrt{\alpha^2+\beta^2}$; son angle vérifie $\cos\theta=\alpha/\rho$, $\sin\theta=\beta/\rho$.
\end{definition}
\begin{propriete}[centre et exception]Si $(\alpha,\beta)\ne(1,0)$, cette similitude possède un unique point fixe, obtenu en résolvant $x'=x,y'=y$. Si $(\alpha,\beta)=(1,0)$, elle est une translation, sans centre unique; pour $p=q=0$, c’est l’identité.\end{propriete}
\begin{demonstration}Le système du point fixe a pour déterminant $(1-\alpha)^2+\beta^2$, strictement positif sauf dans le cas exclu. Autour de ce point fixe, l’expression est une rotation et une multiplication par $\rho$.\end{demonstration}

\section{Exercices et problèmes}
\exo[Images]\label{t11s-c15-e01}
Pour $M(2;1)$, calculer son image par la translation $(3;-2)$, la symétrie de centre $(1;1)$ et la rotation de centre origine d’angle $\pi/2$.
\begin{corrige}[exercice~\ref{t11s-c15-e01}]\label{sol-t11s-c15-e01}
$(5;-1)$; $(0;1)$; $(-1;2)$. Chaque calcul utilise la transformation indiquée, sur le point initial.
\end{corrige}
\exo[Homothétie négative]\label{t11s-c15-e02}
Pour $O(1;2)$, $A(3;2)$ et rapport $-2$, déterminer $A'$ et comparer $OA,OA'$.
\begin{corrige}[exercice~\ref{t11s-c15-e02}]\label{sol-t11s-c15-e02}
$\vect{OA}=(2;0)$, $\vect{OA'}=(-4;0)$, donc $A'=(-3;2)$. $OA=2$, $OA'=4$; le rapport des distances est $2$, pas $-2$.
\end{corrige}
\exo[Deux symétries]\label{t11s-c15-e03}
Pour $A(1;0),B(3;2)$, déterminer $s_B\circ s_A$ et l’image de $(0;1)$.
\begin{corrige}[exercice~\ref{t11s-c15-e03}]\label{sol-t11s-c15-e03}
Translation de $2\vect{AB}=(4;4)$, image $(4;5)$. Un calcul successif donne $s_A(0;1)=(2;-1)$ puis $s_B=(4;5)$.
\end{corrige}
\exo[Un ordre à vérifier]\label{t11s-c15-e04}
Pour la translation $(1;0)$ et la rotation de centre origine d’angle $\pi/2$, calculer les deux composées sur l’origine.
\begin{corrige}[exercice~\ref{t11s-c15-e04}]\label{sol-t11s-c15-e04}
Rotation après translation donne $(0;1)$; translation après rotation donne $(1;0)$. L’ordre change le résultat.
\end{corrige}
\exo[Reconnaître une similitude]\label{t11s-c15-e05}
Étudier $x'=2x-y+1$, $y'=x+2y-2$ : rapport, angle et centre.
\begin{corrige}[exercice~\ref{t11s-c15-e05}]\label{sol-t11s-c15-e05}
$\alpha=2,\beta=1$, rapport $\sqrt5$; angle de cosinus $2/\sqrt5$ et sinus $1/\sqrt5$. Point fixe : $-x+y=1$, $-x-y=-2$, d’où $(1/2;3/2)$.
\end{corrige}
\exo[Construire une figure]\label{t11s-c15-e06}
Un triangle a sommets $A(0;0),B(2;0),C(0;1)$. Construire et calculer son image par la rotation d’angle $\pi/2$ suivie de l’homothétie de rapport $3$ de même centre origine. Comparer les aires.
\begin{corrige}[exercice~\ref{t11s-c15-e06}]\label{sol-t11s-c15-e06}
Images $A'=(0;0),B'=(0;6),C'=(-3;0)$. Aire initiale $1$, aire finale $9$; rapport de similitude $3$ et angle $\pi/2$.
\end{corrige}
\begin{jemeteste}Je calcule les images dans l’ordre prescrit; je distingue rapport signé d’homothétie et rapport positif des distances.\end{jemeteste}
\begin{notepourlaclasse}Les expressions analytiques se prêtent aux travaux dirigés. Les nombres complexes ne sont pas nécessaires. Pour une similitude, faire résoudre le point fixe et contrôler l’exception translation.\end{notepourlaclasse}

\clearpage\bmtchapitre{Trigonométrie}{Convertir les unités d’angle, transformer des expressions et résoudre des équations et inéquations trigonométriques.}{Géométrie}
\label{t11s-c16}
\begin{revision}Cercle trigonométrique, coordonnées et angles remarquables de seconde; produit scalaire.\end{revision}
\begin{automatismes}\exo\label{t11s-c16-auto}Calculer $\sin0$, $\cos\pi$ et $\sin(\pi/6)$. Convertir $90\degre$ en radians.\end{automatismes}
\begin{corrige}[exercice~\ref{t11s-c16-auto}]\label{sol-t11s-c16-auto}$0,-1,1/2$; $\pi/2$ radians.\end{corrige}

\section{Mesures et cercle orienté}
\begin{definition}Un tour correspond à $360\degre=2\pi$ radians $=400$ grades. Les mesures d’un angle orienté diffèrent de $2k\pi$ ($k\in\Z$). Sa mesure principale est l’unique représentant dans $]-\pi;\pi]$.\end{definition}
\begin{exemple}$7\pi/3$ a mesure principale $\pi/3$. $270\degre=3\pi/2=300$ grades; la mesure principale en radians est $-\pi/2$.\end{exemple}
\begin{visualisation}[coordonnées sur le cercle unité]
\begin{center}\begin{tikzpicture}[scale=2.1]
\draw[->](-1.25,0)--(1.35,0)node[right]{$x$};\draw[->](0,-1.2)--(0,1.3)node[above]{$y$};
\draw[bmtSarcelle,thick](0,0)circle(1);
\draw[bmtOcre,thick](0,0)--(60:1)node[above right]{$M$};
\draw[dashed](60:1)--(0.5,0)node[below]{$\cos x$};
\draw[dashed](60:1)--(0,0.866)node[left]{$\sin x$};
\draw[->](0.3,0)arc(0:60:0.3)node[midway,right]{$x$};
\node[below left]at(0,0){$O$};
\end{tikzpicture}\end{center}
Le point représenté a pour angle $x=\pi/3$; les deux projections donnent ses coordonnées.
\end{visualisation}
\section{Formules à utiliser et à justifier}
\begin{propriete}
$\cos(-x)=\cos x$, $\sin(-x)=-\sin x$;
$\cos(\pi-x)=-\cos x$, $\sin(\pi-x)=\sin x$;
$\cos(x+\pi)=-\cos x$, $\sin(x+\pi)=-\sin x$.
\[
\begin{aligned}
\cos(a+b)&=\cos a\cos b-\sin a\sin b,\\
\sin(a+b)&=\sin a\cos b+\cos a\sin b,\\
\cos(a-b)&=\cos a\cos b+\sin a\sin b,\\
\sin(a-b)&=\sin a\cos b-\cos a\sin b.
\end{aligned}
\]
\end{propriete}
\begin{demonstration}Les symétries du cercle donnent les angles associés. Pour l’addition, appliquer la matrice de rotation d’angle $a$ au point $(\cos b;\sin b)$ : les nouvelles coordonnées sont celles de l’angle $a+b$. Les formules de différence suivent en remplaçant $b$ par $-b$.\end{demonstration}
\begin{propriete}[duplication et transformation]
$\sin2x=2\sin x\cos x$;
$\cos2x=\cos^2x-\sin^2x=2\cos^2x-1=1-2\sin^2x$.
\[
\begin{aligned}
\cos a+\cos b&=2\cos\tfrac{a+b}{2}\cos\tfrac{a-b}{2},\\
\sin a+\sin b&=2\sin\tfrac{a+b}{2}\cos\tfrac{a-b}{2},\\
2\cos a\cos b&=\cos(a+b)+\cos(a-b),\\
2\sin a\sin b&=\cos(a-b)-\cos(a+b).
\end{aligned}
\]
\end{propriete}
\begin{demonstration}La duplication prend $a=b=x$. Pour les sommes, poser $u=(a+b)/2,v=(a-b)/2$, puis additionner les formules en $u+v$ et $u-v$. Les produits résultent de l’addition ou de la soustraction des formules d’addition.\end{demonstration}
\section{Équations et inéquations}
\begin{propriete}Pour $k\in\Z$,
\[
\cos x=\cos a\Leftrightarrow x=2k\pi\pm a,
\]
\[
\sin x=\sin a\Leftrightarrow x=a+2k\pi\ \text{ou}\ x=\pi-a+2k\pi.
\]
Pour les tangentes définies, $\tan x=\tan a\Leftrightarrow x=a+k\pi$.
\end{propriete}
\begin{demonstration}Une abscisse sur le cercle correspond aux deux points symétriques par rapport à l’axe horizontal; une ordonnée correspond aux points symétriques par rapport à l’axe vertical. Ajouter les tours complets donne toutes les mesures. La tangente a période $\pi$ et est strictement croissante entre deux pôles.\end{demonstration}
\begin{methode}{une inéquation sur un intervalle}Dessiner les intersections de la droite de niveau avec le cercle; identifier les arcs qui satisfont l’inégalité; traduire les arcs en intervalles; conserver les bornes selon le signe strict ou large. Pour $\tan$, exclure les zéros de $\cos$.\end{methode}
\begin{exemple}Sur $[0;2\pi[$, $\cos x\ge1/2$ donne $[0;\pi/3]\cup[5\pi/3;2\pi[$. Dans $\R$, la solution est la réunion des intervalles $[-\pi/3+2k\pi;\pi/3+2k\pi]$.\end{exemple}

\section{Exercices et problèmes}
\exo[Mesure principale]\label{t11s-c16-e01}
Convertir $150\degre$ en radians et grades; donner la mesure principale de $-17\pi/6$.
\begin{corrige}[exercice~\ref{t11s-c16-e01}]\label{sol-t11s-c16-e01}
$5\pi/6$ radians, $500/3$ grades; $-17\pi/6+2\pi=-5\pi/6$ appartient à $]-\pi;\pi]$.
\end{corrige}
\exo[Valeur exacte]\label{t11s-c16-e02}
Calculer $\cos(\pi/12)$ par une formule d’addition ou de différence.
\begin{corrige}[exercice~\ref{t11s-c16-e02}]\label{sol-t11s-c16-e02}
$\pi/12=\pi/4-\pi/6$. Cosinus $\sqrt2\sqrt3/4+\sqrt2/4=(\sqrt6+\sqrt2)/4$.
\end{corrige}
\exo[Équations]\label{t11s-c16-e03}
Résoudre dans $\R$ : $\sin x=1/2$ et $\cos(2x)=0$.
\begin{corrige}[exercice~\ref{t11s-c16-e03}]\label{sol-t11s-c16-e03}
$x=\pi/6+2k\pi$ ou $5\pi/6+2k\pi$. Pour le cosinus, $2x=\pi/2+k\pi$, donc $x=\pi/4+k\pi/2$, $k\in\Z$.
\end{corrige}
\exo[Inéquation]\label{t11s-c16-e04}
Résoudre $\sin x>\sqrt2/2$ sur $[-\pi;\pi]$.
\begin{corrige}[exercice~\ref{t11s-c16-e04}]\label{sol-t11s-c16-e04}
Arc supérieur entre $\pi/4$ et $3\pi/4$, bornes exclues : $]\pi/4;3\pi/4[$. Aucune solution sur le demi-cercle inférieur.
\end{corrige}
\exo[Transformer]\label{t11s-c16-e05}
Montrer $\sin3x+\sin x=2\sin2x\cos x$ puis résoudre $\sin3x+\sin x=0$ dans $[0;2\pi[$.
\begin{corrige}[exercice~\ref{t11s-c16-e05}]\label{sol-t11s-c16-e05}
Formule somme-produit. $\sin2x=0$ donne $x=k\pi/2$, ce qui inclut déjà les zéros de $\cos x$. Réponse $0,\pi/2,\pi,3\pi/2$.
\end{corrige}
\exo[Tangente]\label{t11s-c16-e06}
Résoudre $\tan x\le1$ sur $]-\pi/2;\pi/2[$. Pourquoi les extrémités sont-elles exclues ?
\begin{corrige}[exercice~\ref{t11s-c16-e06}]\label{sol-t11s-c16-e06}
La tangente est strictement croissante sur cet intervalle; solution $]-\pi/2;\pi/4]$. Les extrémités annulent $\cos x$, donc la tangente n’y existe pas.
\end{corrige}
\exo[Diviser sans perdre]\label{t11s-c16-e07}
Pour résoudre $\sin x\cos x=\sin x$, un élève divise par $\sin x$. Quelles solutions risque-t-il de perdre ? Résoudre dans $\R$.
\begin{corrige}[exercice~\ref{t11s-c16-e07}]\label{sol-t11s-c16-e07}
Factoriser $\sin x(\cos x-1)=0$. Les zéros de $\sin x$ sont $k\pi$, et ceux de $\cos x-1$ sont $2k\pi$, déjà inclus. La division perdrait les multiples impairs de $\pi$.
\end{corrige}
\begin{jemeteste}Je conserve toutes les familles de solutions, je précise l’intervalle et je ne divise pas par une expression sans traiter ses zéros.\end{jemeteste}
\begin{notepourlaclasse}Faire lire le cercle avant le calcul. Choisir des angles exacts; pour la tangente, séparer chaque intervalle entre pôles. Les graphiques sinusoïdaux se relient au chapitre fonctions.\end{notepourlaclasse}

\clearpage\bmtchapitre{Géométrie dans l’espace}{Utiliser un repère direct, calculer produits scalaire et vectoriel et établir les équations des droites et des plans.}{Géométrie}
\label{t11s-c17}
Dans ce chapitre, les calculs de longueurs, d’angles et d’aires à partir de coordonnées se font dans un repère orthonormé. Il est direct pour les rotations et le produit vectoriel.
\begin{revision}Coordonnées et produit scalaire du plan; systèmes linéaires; trigonométrie et déterminants à deux dimensions.\end{revision}
\begin{automatismes}\exo\label{t11s-c17-auto}Calculer $1^2+2^2+2^2$. Résoudre $x+y=3$, $x-y=1$.\end{automatismes}
\begin{corrige}[exercice~\ref{t11s-c17-auto}]\label{sol-t11s-c17-auto}$9$; $(x;y)=(2;1)$.\end{corrige}

\section{Repérage et orientation}
\begin{definition}Un repère de l’espace comporte une origine et trois vecteurs de base non coplanaires. Un repère orthonormé utilise trois vecteurs unitaires orthogonaux. Un repère orthonormé est direct lorsque le produit vectoriel du premier et du deuxième vecteur de base donne le troisième. Dans une base quelconque, être direct exprime seulement une orientation, pas cette égalité.\end{definition}
\begin{exemple}[un cube unité]On repère un cube par $O(0;0;0)$, $A(1;0;0)$, $B(0;1;0)$, $C(0;0;1)$. Le sommet opposé à $O$ a coordonnées $(1;1;1)$ et la grande diagonale vaut $\sqrt3$. Une perspective cavalière aide à lire la configuration, mais ses longueurs dessinées ne donnent pas les longueurs de l’espace.\end{exemple}
\begin{propriete}[produit scalaire]Dans un repère orthonormé, pour $\vecu=(u_1;u_2;u_3)$ et $\vecv=(v_1;v_2;v_3)$, $\vecu\cdot\vecv=u_1v_1+u_2v_2+u_3v_3$. La norme vaut $\sqrt{u_1^2+u_2^2+u_3^2}$, et, pour deux vecteurs non nuls, le produit géométrique vaut $\norme{\vecu}\norme{\vecv}\cos\theta$; si l’un est nul, le produit vaut zéro.\end{propriete}
\section{Produit vectoriel}
\begin{definition}[repère orthonormé direct]
\[
\vecu\wedge\vecv=(u_2v_3-u_3v_2\,;\,
 u_3v_1-u_1v_3\,;\,u_1v_2-u_2v_1).
\]
Ce vecteur est orthogonal aux deux vecteurs, et son sens dépend de l’ordre. Si les deux vecteurs sont non colinéaires, le triplet $(\vecu,\vecv,\vecu\wedge\vecv)$ est direct.
\end{definition}
\begin{propriete}$\vecv\wedge\vecu=-\vecu\wedge\vecv$; pour deux vecteurs non nuls, sa norme vaut $\norme{\vecu}\norme{\vecv}\sin\theta$ avec $\theta\in[0;\pi]$; si l’un est nul, le produit vectoriel est nul. Elle est l’aire du parallélogramme, donc la moitié donne l’aire du triangle. Le produit est nul exactement lorsque les vecteurs sont colinéaires.\end{propriete}
\begin{demonstration}Le produit scalaire avec chaque facteur s’annule en développant les coordonnées. Développer son carré donne $\norme{\vecu\wedge\vecv}^2=\norme{\vecu}^2\norme{\vecv}^2-(\vecu\cdot\vecv)^2$, soit la formule avec $\sin^2\theta$. Le sinus est positif ou nul dans $[0;\pi]$. L’orientation relève de la convention du repère direct.\end{demonstration}
\begin{exemple}Pour $\vecu=(1;2;0)$ et $\vecv=(0;1;3)$, $\vecu\wedge\vecv=(6;-3;1)$. Produits scalaires $6-6=0$ et $-3+3=0$. Norme $\sqrt{46}$, aire du triangle $\sqrt{46}/2$.\end{exemple}
\section{Droites et plans}
\begin{definition}La droite passant par $A(a;b;c)$ de vecteur directeur non nul $(u;v;w)$ a pour équations paramétriques $x=a+ut$, $y=b+vt$, $z=c+wt$, $t\in\R$. Dans l’espace, une seule équation cartésienne linéaire ne décrit pas une droite; on utilise un système de deux équations de plans non parallèles.\end{definition}
\begin{definition}Un plan passant par $A$ et dirigé par deux vecteurs non colinéaires $\vecu,\vecv$ est décrit par $\vect{AM}=s\vecu+t\vecv$, $s,t\in\R$. Pour un vecteur normal non nul $\vecn=(\alpha;\beta;\gamma)$, il a une équation $\alpha x+\beta y+\gamma z+\delta=0$.\end{definition}
\begin{demonstration}L’équation normale résulte de $\vecn\cdot\vect{AM}=0$, avec $\delta=-\alpha a-\beta b-\gamma c$. Pour trois points non alignés, leur produit vectoriel fournit un vecteur normal.\end{demonstration}
\begin{methode}{droite et plan}Substituer la représentation paramétrique de la droite dans l’équation du plan. Une équation linéaire non constante donne un point unique; une égalité impossible donne aucun point (parallélisme strict); une identité donne la droite incluse dans le plan.\end{methode}
\begin{propriete}[distance à un plan]Pour $P:\alpha x+\beta y+\gamma z+\delta=0$, la distance de $A(a;b;c)$ à $P$ vaut $|\alpha a+\beta b+\gamma c+\delta|/\sqrt{\alpha^2+\beta^2+\gamma^2}$.\end{propriete}
\begin{demonstration}La projection orthogonale $H$ de $A$ vérifie $\vect{HA}=\lambda\vecn$. Remplacer dans l’équation du plan donne la longueur $AH$ égale à la formule. Le vecteur normal non nul permet la division.\end{demonstration}

% ENRICHISSEMENT C17

\subsection*{Un normal et un directeur ont des rôles différents}
\begin{methode}{écrire puis vérifier un plan}
Trouver une normale non nulle, par exemple un produit vectoriel de deux directions non colinéaires. Écrire $\vecn\cdot\vect{AM}=0$, puis développer. Vérifier les points de départ dans l’équation finale. Pour une droite, utiliser au contraire un vecteur directeur dans la représentation paramétrique; le normal d’un plan n’est pas une direction contenue dans ce plan.
\end{methode}
\begin{exemple}[projection et distance par le même calcul]
Pour $P:x+2y-z=3$ et $A(1;0;0)$, la normale est $\vecn=(1;2;-1)$, de norme au carré $6$. La valeur du membre $x+2y-z-3$ en $A$ est $-2$. Donc
\[
H=A-\frac{-2}{6}\vecn=(4/3;2/3;-1/3).
\]
Sa substitution donne $3$; le vecteur $\vect{AH}$ est colinéaire à $\vecn$. Le point est bien le projeté orthogonal, et $AH=\sqrt6/3$, ce que confirme la formule de distance $2/\sqrt6$.
\end{exemple}
\begin{methode}{distinguer équations proportionnelles et plans distincts}
Des normales proportionnelles indiquent des plans parallèles, éventuellement confondus. Comparer aussi la constante : multiplier toute l’équation par un même réel non nul ne change pas le plan. Si les normales ne sont pas colinéaires, résoudre les deux équations pour obtenir leur droite d’intersection.
\end{methode}
\begin{remarque}Une substitution donnant $0=0$ signifie que tous les paramètres conviennent, pas que le paramètre vaut zéro. Une substitution impossible interdit tous les points de la droite. Ces trois cas sont les analogues géométriques des issues d’un système linéaire.\end{remarque}

\section{Exercices et problèmes}
\exo[Cube]\label{t11s-c17-e01}
Dans le cube unité, donner les coordonnées du milieu de la grande diagonale issue de $O$ et sa distance à $O$.
\begin{corrige}[exercice~\ref{t11s-c17-e01}]\label{sol-t11s-c17-e01}
Milieu $(1/2;1/2;1/2)$, distance $\sqrt3/2$.
\end{corrige}
\exo[Produits]\label{t11s-c17-e02}
Pour $\vecu=(1;0;1)$, $\vecv=(0;2;1)$, calculer le produit scalaire, le produit vectoriel et l’aire du triangle engendré.
\begin{corrige}[exercice~\ref{t11s-c17-e02}]\label{sol-t11s-c17-e02}
Scalaire $1$; vectoriel $(-2;-1;2)$, norme $3$; aire $3/2$. Les deux produits scalaires de contrôle avec le vectoriel valent $0$.
\end{corrige}
\exo[Plan par trois points]\label{t11s-c17-e03}
Déterminer une équation du plan passant par $A(1;0;0),B(0;1;0),C(0;0;1)$ et une représentation paramétrique.
\begin{corrige}[exercice~\ref{t11s-c17-e03}]\label{sol-t11s-c17-e03}
$\vect{AB}=(-1;1;0)$, $\vect{AC}=(-1;0;1)$, produit vectoriel $(1;1;1)\ne0$. Plan $x+y+z=1$. Paramétrique $(x;y;z)=(1-s-t;s;t)$, $s,t\in\R$.
\end{corrige}
\exo[Intersection]\label{t11s-c17-e04}
La droite $x=t,y=1+t,z=2-t$ coupe-t-elle le plan $x+y+z=4$ ? Calculer le point.
\begin{corrige}[exercice~\ref{t11s-c17-e04}]\label{sol-t11s-c17-e04}
Somme $3+t$, donc $t=1$. Point $(1;2;1)$. Le coefficient en $t$ est non nul, intersection unique.
\end{corrige}
\exo[Droite cartésienne]\label{t11s-c17-e05}
Montrer que le système $x+y=1$, $z=2x$ décrit une droite et en donner un point et un vecteur directeur.
\begin{corrige}[exercice~\ref{t11s-c17-e05}]\label{sol-t11s-c17-e05}
Poser $x=t$ : $(t;1-t;2t)=(0;1;0)+t(1;-1;2)$. Le directeur est non nul; les normales $(1;1;0)$ et $(-2;0;1)$ ne sont pas colinéaires.
\end{corrige}
\exo[Parallèle ou incluse]\label{t11s-c17-e06}
Étudier l’intersection de $x=t,y=1-t,z=2$ avec les plans $x+y+z=3$ et $x+y+z=4$.
\begin{corrige}[exercice~\ref{t11s-c17-e06}]\label{sol-t11s-c17-e06}
La somme vaut constamment $3$. La droite est incluse dans le premier; aucune intersection avec le second, auquel elle est strictement parallèle.
\end{corrige}
\exo[Distance]\label{t11s-c17-e07}
Calculer la distance de $A(1;2;3)$ au plan $x+2y+2z=2$.
\begin{corrige}[exercice~\ref{t11s-c17-e07}]\label{sol-t11s-c17-e07}
Valeur du membre $1+4+6-2=9$, norme du normal $3$, donc distance $3$. Projection $H=A-(9/9)(1;2;2)=(0;0;1)$, qui appartient au plan et vérifie $AH=3$.
\end{corrige}

\exo[Une équation, trois types d’intersection]\label{t11s-c17-p01}
Dans un repère orthonormé, on considère $P:2x-y+z=4$ et les droites $d_1:(t;0;0)$, $d_2:(t;2t;0)$, $d_3:(t;2t;4)$, $t\in\R$. Déterminer chaque intersection avec $P$ et classer la droite comme sécante, strictement parallèle ou incluse. Dans chaque cas, écrire l’équation obtenue après substitution.
\begin{corrige}[exercice~\ref{t11s-c17-p01}]\label{sol-t11s-c17-p01}
Pour $d_1$, $2t=4$ impose $t=2$ : intersection unique $(2;0;0)$, droite sécante. Pour $d_2$, on obtient $0=4$, impossible : droite strictement parallèle. Pour $d_3$, l’égalité est $4=4$ pour tout $t$ : droite incluse. Les directions sont non nulles, et les conclusions portent sur toutes les valeurs réelles du paramètre.
\end{corrige}
\exo[Plans confondus, parallèles ou sécants]\label{t11s-c17-p02}
Pour $P:x+2y-z=3$, comparer les plans $Q:2x+4y-2z=6$, $R:2x+4y-2z=7$ et $S:x-y+z=0$. Justifier les positions relatives. Pour $P\cap S$, donner une représentation paramétrique de la droite et vérifier un point et sa direction dans les deux plans.
\begin{corrige}[exercice~\ref{t11s-c17-p02}]\label{sol-t11s-c17-p02}
L’équation de $Q$ est deux fois celle de $P$, donc ils sont confondus. Celle de $R$ a la même normale, mais impose à deux fois le membre de $P$ la valeur $7$ au lieu de $6$ : parallèles distincts. Les normales de $P,S$ ne sont pas colinéaires. Poser $z=3t$ donne $y=1+2t$, $x=1-t$, donc $(x;y;z)=(1;1;0)+t(-1;2;3)$. Le point $(1;1;0)$ vérifie $P,S$, et la direction a produits scalaires $-1+4-3=0$ et $-1-2+3=0$ avec leurs normales. La paramétrisation décrit toutes les solutions du système.
\end{corrige}

\begin{jemeteste}Je précise le repère direct pour le produit vectoriel, je distingue un vecteur directeur d’un normal et je vérifie les équations sur les points.\end{jemeteste}
\begin{notepourlaclasse}Manipuler un cube avant les coordonnées. Deux équations cartésiennes pour une droite répondent au programme sans entretenir l’ambiguïté d’une « équation unique ». La formule de distance complète les problèmes métriques.\end{notepourlaclasse}

\clearpage\begingroup
\chapter*{Évaluation : géométrie}\markboth{Évaluation : géométrie}{Évaluation : géométrie}\addcontentsline{toc}{chapter}{Évaluation : géométrie}
\renewcommand{\thebmtexo}{DS3.\arabic{bmtexo}}
\renewcommand{\theHbmtexo}{DS3.\arabic{bmtexo}}
\setcounter{bmtexo}{0}
\begin{devoir}[2 h — 20 points]
Calculatrice autorisée. Justifier les réponses et préciser les domaines utiles. Les questions sont indépendantes sauf indication. 
\end{devoir}
\exo[Produit scalaire et barycentre]\label{t11s-ds3-e01}\bareme{6 points}
Dans un repère orthonormé, $A(0;0),B(6;0),C(2;4)$.
\begin{enumerate}\item Calculer $\vect{AB}\cdot\vect{AC}$ et l’aire du triangle (2 points).
\item Calculer le barycentre $G$ de $(A,1),(B,1),(C,2)$ (2 points).
\item Déterminer le lieu $MA^2+MB^2=26$ (2 points).\end{enumerate}
\begin{corrige}[exercice~\ref{t11s-ds3-e01}]\label{sol-t11s-ds3-e01}
1. Produit $12$; aire $6\times4/2=12$ (1+1).
2. Somme $4$, $G=((0+6+4)/4;(0+0+8)/4)=(5/2;2)$ (condition 1, coordonnées 1).
3. $I=(3;0)$, somme $2MI^2+18=26$, donc cercle de centre $I$ et rayon $2$ (réduction 1, lieu 1).
\end{corrige}
\exo[Trigonométrie]\label{t11s-ds3-e02}\bareme{4 points}
\begin{enumerate}\item Déterminer la mesure principale de $19\pi/6$ (1 point).
\item Résoudre $2\cos^2x-1=0$ sur $[0;2\pi[$ (2 points).
\item Résoudre $\sin x\ge1/2$ sur le même intervalle (1 point).\end{enumerate}
\begin{corrige}[exercice~\ref{t11s-ds3-e02}]\label{sol-t11s-ds3-e02}
1. $19\pi/6-4\pi=-5\pi/6$.
2. $\cos2x=0$, donc $x=\pi/4+k\pi/2$, soit $\pi/4,3\pi/4,5\pi/4,7\pi/4$ (famille 1, intervalle 1).
3. $[\pi/6;5\pi/6]$.
\end{corrige}
\exo[Similitude]\label{t11s-ds3-e03}\bareme{4 points}
Étudier $x'=-2y+1$, $y'=2x+2$ : rapport et angle, centre, image de $(1;0)$ et facteur d’évolution des aires.
\begin{corrige}[exercice~\ref{t11s-ds3-e03}]\label{sol-t11s-ds3-e03}
Rapport $2$, angle $\pi/2$ (1). Point fixe $x=-2y+1$, $y=2x+2$, donnant $x=-3/5,y=4/5$ (1). Image $(1;4)$ (1); facteur des aires $4$ (1).
\end{corrige}
\exo[Un plan dans l’espace]\label{t11s-ds3-e04}\bareme{6 points}
En repère orthonormé direct, $A(1;0;0),B(0;1;0),C(0;0;1)$.
\begin{enumerate}\item Calculer $\vect{AB}\wedge\vect{AC}$ et donner une équation de $(ABC)$ (2 points).
\item Déterminer l’intersection avec la droite $(x;y;z)=(t;2t;3t)$ (2 points).
\item Calculer la distance de l’origine au plan et l’aire du triangle $ABC$ (2 points).\end{enumerate}
\begin{corrige}[exercice~\ref{t11s-ds3-e04}]\label{sol-t11s-ds3-e04}
1. Vectoriel $(1;1;1)$; plan $x+y+z=1$ (1+1).
2. $6t=1$, donc $(1/6;1/3;1/2)$, unique (1+1).
3. Distance $1/\sqrt3$; aire $\sqrt3/2$ (1+1).
\end{corrige}
\endgroup


% ===========================================================================
\bmtpartie{bmtPalissandre}{IV}{Traitement des données et probabilités}{Décrire, compter, calculer}{%
  Trois chapitres : c'est ici que les mathématiques
  rejoignent le plus directement ce que vous lirez dans un journal ou un
  rapport. Décrire une série, compter sans énumérer, mesurer une chance : les
  trois gestes tiennent ensemble.}

\clearpage\bmtchapitre{Organisation des données et statistique}{Calculer proportions et évolutions, résumer une série par ses indicateurs et interpréter une boîte à moustaches.}{Données et probabilités}
\label{t11s-c18}
\begin{revision}Pourcentages, effectifs, fréquences, ordre des valeurs et moyenne pondérée.\end{revision}
\begin{automatismes}\exo\label{t11s-c18-auto}Calculer $20\%$ de $150$; ranger $8,3,5,3$; calculer leur moyenne.\end{automatismes}
\begin{corrige}[exercice~\ref{t11s-c18-auto}]\label{sol-t11s-c18-auto}$30$; $3,3,5,8$; moyenne $19/4=4{,}75$.\end{corrige}

\section{Proportion et évolution}
\begin{definition}Une proportion est un rapport partie/ensemble, avec effectif total non nul. Le taux d’évolution d’une valeur initiale $V_i>0$ à une valeur finale $V_f$ est $t=(V_f-V_i)/V_i$. Le coefficient multiplicateur est $1+t$; un taux de $p\%$ vaut $p/100$.\end{definition}
\begin{propriete}Deux évolutions successives de taux $t_1,t_2$ ont pour coefficient $(1+t_1)(1+t_2)$; leur taux global est ce produit moins $1$. Si $1+t\ne0$, le taux réciproque est $1/(1+t)-1$.\end{propriete}
\begin{demonstration}La valeur finale après deux évolutions vaut $V_i(1+t_1)(1+t_2)$. Pour revenir en arrière, multiplier par l’inverse du coefficient.\end{demonstration}
\begin{exemple}Hausse de $20\%$ puis baisse de $20\%$ : coefficient $1{,}2\times0{,}8=0{,}96$, baisse globale $4\%$. Passer de $30\%$ à $40\%$ est une hausse de $10$ points de pourcentage, mais une hausse relative de $1/3$, soit environ $33{,}3\%$.\end{exemple}
\section{Indicateurs de position}
\begin{definition}Pour une série ordonnée $x_1\le\cdots\le x_N$, $N\ge1$, la moyenne est $\bar x=\frac1N\sum x_i$; avec des effectifs $n_j$, elle vaut $\sum n_jx_j/\sum n_j$. La médiane est la valeur centrale si $N$ est impair, la moyenne des deux centrales si $N$ est pair.\end{definition}
\begin{definition}[convention de rang]
$Q_1$ est la valeur de rang $\lceil N/4\rceil$, $Q_3$ celle de rang $\lceil3N/4\rceil$. Le décile $D_k$ est la valeur de rang $\lceil kN/10\rceil$, $k=1,\ldots,9$. $\lceil a\rceil$ désigne le plus petit entier supérieur ou égal à $a$. Ces conventions sont précisées car elles peuvent différer selon les logiciels.
\end{definition}
\section{Dispersion et représentation}
\begin{definition}L’étendue vaut maximum moins minimum; l’écart interquartile vaut $Q_3-Q_1$. La variance descriptive vaut $V=\frac1N\sum(x_i-\bar x)^2$; l’écart type est $\sigma=\sqrt V$. Une boîte à moustaches représente minimum, $Q_1$, médiane, $Q_3$, maximum, selon la convention utilisée ici.\end{definition}
\begin{propriete}$V=\frac1N\sum x_i^2-\bar x^2$. La variance est positive ou nulle, et vaut zéro exactement lorsque toutes les valeurs sont égales.\end{propriete}
\begin{demonstration}Développer les carrés et utiliser $\sum x_i=N\bar x$. La première définition est une moyenne de carrés positifs ou nuls; elle est nulle exactement si chaque carré est nul.\end{demonstration}
\begin{exemple}[huit temps en minutes]
Série $4,6,6,8,10,10,12,14$. $N=8$, somme $70$, moyenne $8{,}75$, médiane $9$; rangs $2$ et $6$, donc $Q_1=6,Q_3=10$. Étendue $10$, écart interquartile $4$. Somme des carrés $792$; variance $99-(8{,}75)^2=22{,}4375$; écart type $\sqrt{22{,}4375}\simeq4{,}737$ minutes. La variance s’exprime en minutes carrées.
\end{exemple}
\begin{visualisation}[la boîte ne montre pas chaque donnée]
\begin{center}\begin{tikzpicture}[x=0.55cm,y=1cm]
\draw[->](3,0)--(15,0)node[right]{minutes};
\foreach \x in {4,6,8,10,12,14}{\draw(\x,0.06)--(\x,-0.06)node[below]{$\x$};}
\draw[bmtPalissandre,thick](4,0.9)--(6,0.9);\draw[bmtPalissandre,thick](10,0.9)--(14,0.9);
\draw[bmtPalissandre,thick](6,0.55)rectangle(10,1.25);\draw[bmtOcre,thick](9,0.55)--(9,1.25);
\draw(4,0.7)--(4,1.1);\draw(14,0.7)--(14,1.1);
\end{tikzpicture}\end{center}
Le trait intérieur est la médiane $9$, et non la moyenne $8{,}75$.
\end{visualisation}
\section{Choisir un résumé sans surinterpréter}
Une moyenne sensible aux valeurs extrêmes et une médiane n’apportent pas la même information. Comparer des séries avec mêmes unités et conventions. Pour des données groupées en classes, remplacer chaque valeur par le centre de sa classe donne une moyenne et une variance approchées; les données originales ne peuvent être reconstruites exactement.

\section{Exercices et problèmes}
\exo[Taux successifs]\label{t11s-c18-e01}
Un prix de $50000$ ariary augmente de $10\%$, puis diminue de $5\%$. Calculer le prix final et le taux global.
\begin{corrige}[exercice~\ref{t11s-c18-e01}]\label{sol-t11s-c18-e01}
$50000\times1{,}1\times0{,}95=52250$ ariary. Coefficient $1{,}045$, hausse globale $4{,}5\%$.
\end{corrige}
\exo[Retour au départ]\label{t11s-c18-e02}
Après une hausse de $25\%$, quelle baisse ramène à la valeur initiale ?
\begin{corrige}[exercice~\ref{t11s-c18-e02}]\label{sol-t11s-c18-e02}
Coefficient retour $1/1{,}25=0{,}8$, donc baisse $20\%$. Les pourcentages portent sur des bases différentes.
\end{corrige}
\exo[Effectifs]\label{t11s-c18-e03}
Les valeurs $5,10,15$ ont respectivement les effectifs $2,5,3$. Calculer moyenne, médiane, quartiles, $D_1,D_9$, variance et écart type.
\begin{corrige}[exercice~\ref{t11s-c18-e03}]\label{sol-t11s-c18-e03}
$N=10$, somme $105$, moyenne $10{,}5$. Série ordonnée : deux $5$, cinq $10$, trois $15$; médiane $10$, $Q_1=10$ (rang $3$), $Q_3=15$ (rang $8$), $D_1=5$, $D_9=15$. Moyenne des carrés $(50+500+675)/10=122{,}5$; variance $12{,}25$, écart type $3{,}5$.
\end{corrige}
\exo[Mêmes moyennes]\label{t11s-c18-e04}
Comparer les séries $(8,10,12)$ et $(0,10,20)$ : moyenne, médiane, étendue et variance. Que montre la comparaison ?
\begin{corrige}[exercice~\ref{t11s-c18-e04}]\label{sol-t11s-c18-e04}
Moyenne et médiane $10$ pour les deux. Étendues $4$ et $20$; variances $8/3$ et $200/3$. Une même moyenne ne garantit pas une même dispersion.
\end{corrige}
\exo[Décalage]\label{t11s-c18-e05}
Si on ajoute $3$ à toutes les valeurs d’une série, que deviennent moyenne, médiane, variance et écart type ?
\begin{corrige}[exercice~\ref{t11s-c18-e05}]\label{sol-t11s-c18-e05}
Moyenne et médiane augmentent de $3$. Les écarts à la moyenne restent identiques : variance et écart type inchangés. L’ordre des données est conservé.
\end{corrige}
\exo[Boîte à moustaches]\label{t11s-c18-e06}
Pour la série $2,3,4,4,5,7,8,9,10,12$, calculer les cinq valeurs nécessaires à la boîte et la tracer sur un axe gradué.
\begin{corrige}[exercice~\ref{t11s-c18-e06}]\label{sol-t11s-c18-e06}
Minimum $2$, $Q_1=x_3=4$, médiane $(5+7)/2=6$, $Q_3=x_8=9$, maximum $12$. Boîte de $4$ à $9$, trait à $6$, moustaches jusqu’à $2$ et $12$.
\end{corrige}
\exo[Des classes]\label{t11s-c18-e07}
Dix trajets durent entre $0$ et $10$ min, vingt entre $10$ et $20$ min. Estimer la moyenne par les centres. Cette valeur est-elle exacte ?
\begin{corrige}[exercice~\ref{t11s-c18-e07}]\label{sol-t11s-c18-e07}
Centres $5$ et $15$; moyenne estimée $(10\times5+20\times15)/30=35/3\simeq11{,}67$ min. Ce n’est pas une moyenne exacte sans les temps individuels; les intervalles doivent être disjoints, par exemple $[0;10[$ et $[10;20]$.
\end{corrige}
\begin{jemeteste}Je vérifie l’effectif total, les rangs, les unités et la distinction entre pourcentage et points de pourcentage.\end{jemeteste}
\begin{notepourlaclasse}Les données sont fictives pour l’apprentissage. La boîte utilise minimum et maximum aux moustaches; ne pas mêler à une convention de détection de valeurs atypiques sans l’annoncer.\end{notepourlaclasse}

\clearpage\bmtchapitre{Ensembles, applications et dénombrement}{Choisir entre produit, permutation, arrangement et combinaison; reconnaître injection, surjection et bijection.}{Données et probabilités}
\label{t11s-c19}
\begin{revision}Ensembles finis, union, intersection, complémentaire; multiplication et factorielle.\end{revision}
\begin{automatismes}\exo\label{t11s-c19-auto}Lister les couples de $\{a,b\}\times\{1,2\}$. Calculer $4\times3\times2\times1$.\end{automatismes}
\begin{corrige}[exercice~\ref{t11s-c19-auto}]\label{sol-t11s-c19-auto}$(a,1),(a,2),(b,1),(b,2)$; $24$.\end{corrige}

\section{Cardinal et opérations}
\begin{definition}Le cardinal $|E|$ d’un ensemble fini est son nombre d’éléments. L’intersection contient les éléments communs; l’union ceux de l’un ou de l’autre; le complémentaire $\overline A$ est pris dans un ensemble de référence $E$ annoncé.\end{definition}
\begin{propriete}Pour $A,B$ finis, $|A\cup B|=|A|+|B|-|A\cap B|$ et $|\overline A|=|E|-|A|$. Le principe du produit multiplie les nombres de choix successifs lorsque le nombre de possibilités à chaque étape est fixé pour tous les choix précédents.\end{propriete}
\begin{demonstration}La somme des deux cardinaux compte deux fois l’intersection. Pour le produit, grouper les résultats par premier choix : chaque groupe a le même nombre de possibilités suivantes; répéter le raisonnement.\end{demonstration}
\section{Applications entre ensembles finis}
\begin{definition}Une application $f:E\to F$ associe à chaque élément de $E$ un unique élément de $F$. Elle est injective si deux éléments distincts ont des images distinctes; surjective si chaque élément de $F$ possède au moins un antécédent; bijective si elle est injective et surjective.\end{definition}
\begin{propriete}Si $|E|=p,|F|=n$, avec $n\ge1$, il existe $n^p$ applications. Une injection exige $p\le n$; leur nombre est $n(n-1)\cdots(n-p+1)$ pour $p\ge1$, et $1$ pour $p=0$. Pour $p=n$, les bijections sont $n!$. Une surjection exige $p\ge n$ si $F$ est non vide.\end{propriete}
\begin{demonstration}Une application choisit indépendamment une image parmi $n$ pour chacun des $p$ éléments. Une injection retire à chaque étape les images déjà utilisées. Une injection entre deux ensembles de même cardinal utilise toutes les images, donc elle est aussi surjective. La condition pour la surjection résulte du nombre d’images disponibles.\end{demonstration}
\begin{remarque}La formule $n^p$ compte toutes les applications, pas seulement les surjections. Le programme demande de reconnaître la surjection, sans formule générale de comptage; les petits exemples peuvent être énumérés.\end{remarque}
\section{Ordre et répétition}
\begin{definition}$n!=1\times2\times\cdots\times n$, avec $0!=1$. Une permutation ordonne tous les éléments distincts : $n!$ choix. Un arrangement de $p$ éléments distincts parmi $n$ est une liste ordonnée sans répétition : $A_n^p=n!/(n-p)!$. Une combinaison est un sous-ensemble de $p$ éléments : $\binom np=n!/[p!(n-p)!]$, pour $0\le p\le n$.\end{definition}
\begin{demonstration}Chaque sous-ensemble de $p$ éléments donne exactement $p!$ ordres; les arrangements comptent chacun $p!$ fois, donc on divise par $p!$.\end{demonstration}
\begin{exemple}Parmi $8$ élèves, un podium de trois places se compte par $8\times7\times6=336$; un groupe de trois élèves se compte par $\binom83=56$. Un code de trois chiffres avec répétition autorisée a $10^3=1000$ possibilités, si le zéro initial est accepté.\end{exemple}
\begin{methode}{poser trois questions}L’ordre distingue-t-il deux résultats ? La répétition est-elle autorisée ? Choisit-on tous les éléments ou seulement certains ? Donner un exemple de résultat permet de choisir le bon univers de comptage.\end{methode}

\section{Exercices et problèmes}
\exo[Union]\label{t11s-c19-e01}
Dans un groupe de $40$ élèves, $25$ pratiquent le football, $18$ le basket et $10$ les deux. Combien pratiquent au moins un de ces sports ? Aucun ?
\begin{corrige}[exercice~\ref{t11s-c19-e01}]\label{sol-t11s-c19-e01}
Union $25+18-10=33$; complément $40-33=7$. Football seulement $15$, basket seulement $8$, ce qui vérifie le total.
\end{corrige}
\exo[Codes]\label{t11s-c19-e02}
Compter les codes de quatre chiffres avec répétition autorisée, puis sans répétition, si le premier chiffre peut être zéro. Recommencer si zéro est interdit au premier rang.
\begin{corrige}[exercice~\ref{t11s-c19-e02}]\label{sol-t11s-c19-e02}
Zéro initial permis : $10^4=10000$; sans répétition $10\times9\times8\times7=5040$. Zéro initial interdit : $9\times10^3=9000$; sans répétition $9\times9\times8\times7=4536$.
\end{corrige}
\exo[Responsabilités]\label{t11s-c19-e03}
Parmi $12$ personnes, choisir un président et un secrétaire distincts, puis un comité de trois personnes sans fonctions distinctes.
\begin{corrige}[exercice~\ref{t11s-c19-e03}]\label{sol-t11s-c19-e03}
Président/secrétaire : $12\times11=132$. Comité : $\binom{12}3=220$. Les deux choix sont deux questions séparées.
\end{corrige}
\exo[Applications]\label{t11s-c19-e04}
Pour $E=\{1,2,3\}$ et $F=\{a,b\}$, compter les applications et les injections; compter les surjections par retrait des applications constantes.
\begin{corrige}[exercice~\ref{t11s-c19-e04}]\label{sol-t11s-c19-e04}
$2^3=8$ applications, aucune injection. Deux applications constantes ne sont pas surjectives; toutes les autres atteignent les deux éléments : $6$ surjections.
\end{corrige}
\exo[Frontières]\label{t11s-c19-e05}
Expliquer $\binom n0=\binom nn=1$ et $\binom np=\binom n{n-p}$.
\begin{corrige}[exercice~\ref{t11s-c19-e05}]\label{sol-t11s-c19-e05}
Un seul sous-ensemble vide et un seul égal à l’ensemble. Le complémentaire associe bijectivement les sous-ensembles de taille $p$ à ceux de taille $n-p$.
\end{corrige}
\exo[Choix mixtes]\label{t11s-c19-e06}
Un groupe contient $5$ filles et $4$ garçons. Compter les comités de trois élèves comprenant exactement deux filles, puis au moins une fille.
\begin{corrige}[exercice~\ref{t11s-c19-e06}]\label{sol-t11s-c19-e06}
Exactement deux : $\binom52\binom41=40$. Au moins une : $\binom93-\binom43=84-4=80$. Les élèves sont distincts, et l’ordre ne compte pas.
\end{corrige}
\exo[Listes ou groupes]\label{t11s-c19-e07}
Un élève calcule un comité de trois parmi huit par $8\times7\times6$. Expliquer précisément le surcomptage et réparer.
\begin{corrige}[exercice~\ref{t11s-c19-e07}]\label{sol-t11s-c19-e07}
Chaque comité apparaît sous $3!=6$ ordres différents; diviser $336$ par $6$ donne $56$. Il ne faut pas diviser par $3$, car trois personnes ont six ordres.
\end{corrige}
\begin{jemeteste}Je définis un résultat avant de compter et je contrôle les frontières $p=0$ et $p=n$.\end{jemeteste}
\begin{notepourlaclasse}Faire représenter d’abord les petites situations par arbres ou tableaux. Éviter les formules de surjection générales et les variables aléatoires qui ne font pas partie de ce chapitre.\end{notepourlaclasse}

\clearpage\bmtchapitre{Probabilités sur un univers fini}{Décrire l’univers, calculer la probabilité d’un événement et exploiter unions, intersections et événements contraires.}{Données et probabilités}
\label{t11s-c20}
\begin{revision}Dénombrement, fractions et ensembles finis du chapitre précédent.\end{revision}
\begin{automatismes}\exo\label{t11s-c20-auto}Calculer $1-2/5$; simplifier $6/15$; combien de couples ordonnés donne le lancer de deux dés à six faces ?\end{automatismes}
\begin{corrige}[exercice~\ref{t11s-c20-auto}]\label{sol-t11s-c20-auto}$3/5$, $2/5$, $36$ couples.\end{corrige}

\section{Univers et événements}
\begin{definition}L’univers $\Omega$ est l’ensemble des issues d’une expérience aléatoire. Un événement est un sous-ensemble de $\Omega$; un événement élémentaire contient une issue. $\varnothing$ est impossible, $\Omega$ certain. Deux événements sont incompatibles lorsque leur intersection est vide.\end{definition}
\begin{definition}[probabilité finie]À chaque issue $\omega$ on associe un nombre $p_\omega\ge0$ avec $\sum_{\omega\in\Omega}p_\omega=1$. Pour $A\subseteq\Omega$, $\Prob(A)=\sum_{\omega\in A}p_\omega$. Dans un univers fini non vide équiprobable, chaque poids vaut $1/|\Omega|$, donc $\Prob(A)=|A|/|\Omega|$.\end{definition}
\begin{avertissement}L’équiprobabilité doit venir du modèle, par exemple un dé équilibré ou un tirage uniforme. Les valeurs d’une somme de deux dés ne sont pas équiprobables : la somme $7$ a plus de couples que la somme $2$.\end{avertissement}
\section{Propriétés}
\begin{propriete}
$0\le\Prob(A)\le1$, $\Prob(\varnothing)=0$, $\Prob(\Omega)=1$;
$\Prob(\overline A)=1-\Prob(A)$;
\[
\Prob(A\cup B)=\Prob(A)+\Prob(B)-\Prob(A\cap B).
\]
Si $A\subseteq B$, $\Prob(A)\le\Prob(B)$.
\end{propriete}
\begin{demonstration}Les sommes portent sur des poids positifs ou nuls. Le complément partage l’univers avec $A$. La formule de réunion retire les poids de l’intersection comptés deux fois. Pour l’inclusion, les poids supplémentaires de $B\setminus A$ sont positifs ou nuls.\end{demonstration}
\begin{exemple}[un dé équilibré]Pour $A=\{2,4,6\}$ et $B=\{4,5,6\}$, $\Prob(A)=\Prob(B)=1/2$, $A\cap B=\{4,6\}$ donc $\Prob(A\cap B)=1/3$; la réunion a probabilité $2/3$. Les événements ne sont pas incompatibles.\end{exemple}
\section{Compter un tirage}
\begin{methode}{uniformiser le modèle}Pour deux lancers successifs équilibrés et indépendants, les couples ordonnés sont équiprobables. Pour un tirage simultané uniforme de $p$ objets distincts parmi $n$, les sous-ensembles sont équiprobables. Compter favorable et total avec la même notion de résultat; vérifier que la probabilité est comprise entre zéro et un.\end{methode}
\begin{exemple}[tirage sans remise]
Une urne contient $3$ boules rouges et $2$ bleues, toutes distinguables. Tirer simultanément deux boules de façon uniforme. Il y a $\binom52=10$ paires; les deux rouges donnent $\binom32=3$ paires, probabilité $3/10$. Au moins une rouge est le contraire de deux bleues : $1-\binom22/\binom52=9/10$.
\end{exemple}
\begin{remarque}Dans un petit univers, un tableau ou une liste complète évite une formule inadaptée. Les probabilités conditionnelles, les lois de variables aléatoires et la loi binomiale ne sont pas nécessaires pour ces exercices.\end{remarque}
\section{Fréquence observée et probabilité du modèle}
Une fréquence mesure la proportion observée dans une série d’expériences; une probabilité appartient au modèle. Une fréquence de $0{,}47$ pour « pile » n’invalide pas automatiquement un modèle à $1/2$. Le lien entre fréquences et probabilités se découvre par répétition; une expérience courte ne démontre pas l’équilibre d’une pièce.

\begin{situation}[fréquence et modèle]
\exo\label{t11s-c20-act}Une série de $50$ lancers d’une pièce a donné $23$ fois pile. Calculer la fréquence de pile et la comparer au modèle d’une pièce équilibrée. Proposer un protocole de répétition et un tableau pour comparer plusieurs séries, sans prétendre qu’une seule série prouve l’équilibre.
\end{situation}
\begin{corrige}[exercice~\ref{t11s-c20-act}]\label{sol-t11s-c20-act}
La fréquence est $23/50=0{,}46$, contre une probabilité de $1/2$ dans le modèle équilibré. Répéter le même nombre de lancers dans les mêmes conditions, noter pour chaque série son numéro, le nombre de piles et sa fréquence; ajouter les effectifs pour une fréquence globale. Les résultats fluctuent : le seul écart $0{,}04$ ne prouve ni l’équilibre ni le déséquilibre. Les $23$ piles sont des données d’exercice, pas le compte rendu d’une expérience effectuée.
\end{corrige}
\section{Exercices et problèmes}
\exo[Deux dés]\label{t11s-c20-e01}
Deux dés équilibrés sont lancés indépendamment. Énumérer les couples de somme $7$ et calculer sa probabilité, puis celle d’une somme $2$.
\begin{corrige}[exercice~\ref{t11s-c20-e01}]\label{sol-t11s-c20-e01}
Univers de $36$ couples équiprobables. Somme $7$ : $(1,6),(2,5),(3,4),(4,3),(5,2),(6,1)$, donc $1/6$. Somme $2$ : $(1,1)$, donc $1/36$.
\end{corrige}
\exo[Union]\label{t11s-c20-e02}
Un dé équilibré est lancé. $A$ : résultat pair; $B$ : résultat au moins $5$. Calculer intersection, réunion et contraire de la réunion.
\begin{corrige}[exercice~\ref{t11s-c20-e02}]\label{sol-t11s-c20-e02}
Intersection $\{6\}$ : $1/6$. Réunion $\{2,4,5,6\}$ : $2/3$. Contraire $\{1,3\}$ : $1/3$.
\end{corrige}
\exo[Urne]\label{t11s-c20-e03}
Une urne contient $4$ rouges et $3$ bleues distinguables. Tirer simultanément trois boules de façon uniforme. Calculer les probabilités de trois rouges, exactement deux rouges et au moins une bleue.
\begin{corrige}[exercice~\ref{t11s-c20-e03}]\label{sol-t11s-c20-e03}
Total $\binom73=35$. Trois rouges $\binom43=4$, donc $4/35$. Exactement deux : $\binom42\binom31=18$, donc $18/35$. Au moins une bleue : $1-4/35=31/35$.
\end{corrige}
\exo[Avec remise]\label{t11s-c20-e04}
Dans la même urne, tirer successivement deux boules avec remise, indépendamment et uniformément à chaque fois. Calculer la probabilité de deux rouges puis de couleurs différentes.
\begin{corrige}[exercice~\ref{t11s-c20-e04}]\label{sol-t11s-c20-e04}
Les $7^2=49$ couples de boules physiques sont équiprobables. Deux rouges : $4^2/49=16/49$. Couleurs différentes : $(4\times3+3\times4)/49=24/49$.
\end{corrige}
\exo[Probabilités imposées]\label{t11s-c20-e05}
On donne $\Prob(A)=0{,}6$, $\Prob(B)=0{,}5$, $\Prob(A\cap B)=0{,}2$. Calculer la réunion, puis sa probabilité contraire. Ces données peuvent-elles correspondre à des événements incompatibles ?
\begin{corrige}[exercice~\ref{t11s-c20-e05}]\label{sol-t11s-c20-e05}
Réunion $0{,}9$, contraire $0{,}1$. Non, incompatibilité imposerait intersection vide et probabilité $0$. Les quatre régions ont masses $0{,}2,0{,}4,0{,}3,0{,}1$, positives et de somme $1$, donc données cohérentes.
\end{corrige}
\exo[Modèle non uniforme]\label{t11s-c20-e06}
Une roue porte trois secteurs de tailles différentes. Peut-on affirmer que chacun a probabilité $1/3$ ? Proposer un modèle si les angles sont $90\degre,90\degre,180\degre$ et le point d’arrêt est uniforme en angle.
\begin{corrige}[exercice~\ref{t11s-c20-e06}]\label{sol-t11s-c20-e06}
Non, le nombre de secteurs ne suffit pas. Dans le modèle uniforme en angle, les probabilités sont $1/4,1/4,1/2$. Elles somment à $1$.
\end{corrige}
\exo[Un tirage de comité]\label{t11s-c20-e07}
On choisit uniformément un comité de trois parmi neuf élèves, dont cinq filles. Quelle est la probabilité d’au moins une fille ?
\begin{corrige}[exercice~\ref{t11s-c20-e07}]\label{sol-t11s-c20-e07}
Univers $\binom93=84$ comités. Sans fille : $\binom43=4$. Probabilité $1-4/84=20/21$. Le choix uniforme des comités doit être explicitement supposé.
\end{corrige}
\begin{jemeteste}Je justifie le modèle équiprobable, je compte des issues cohérentes et je vérifie les probabilités par leur complément.\end{jemeteste}
\begin{notepourlaclasse}Ne pas confondre incompatibilité et indépendance. L’hypothèse d’indépendance des deux dés établit le modèle des couples; aucune règle générale de probabilité conditionnelle n’est requise.\end{notepourlaclasse}

\clearpage\begingroup
\chapter*{Évaluation : données et probabilités}\markboth{Évaluation : données et probabilités}{Évaluation : données et probabilités}\addcontentsline{toc}{chapter}{Évaluation : données et probabilités}
\renewcommand{\thebmtexo}{DS4.\arabic{bmtexo}}
\renewcommand{\theHbmtexo}{DS4.\arabic{bmtexo}}
\setcounter{bmtexo}{0}
\begin{devoir}[2 h — 20 points]
Calculatrice autorisée. Justifier les réponses et préciser les domaines utiles. Les questions sont indépendantes sauf indication. 
\end{devoir}
\exo[Évolutions]\label{t11s-ds4-e01}\bareme{4 points}
Un prix vaut $80000$ ariary; il subit une baisse de $15\%$ puis une hausse de $10\%$. Calculer prix final, taux global et taux qui ramènerait le prix final à $80000$ ariary.
\begin{corrige}[exercice~\ref{t11s-ds4-e01}]\label{sol-t11s-ds4-e01}
Prix $80000\times0{,}85\times1{,}1=74800$ (1). Coefficient $0{,}935$, baisse globale $6{,}5\%$ (1). Coefficient retour $80000/74800=200/187$, hausse $13/187\simeq6{,}95\%$ (2).
\end{corrige}
\exo[Statistique]\label{t11s-ds4-e02}\bareme{6 points}
Une série est $4,6,6,8,10,10,12,14$. Calculer moyenne, médiane, quartiles, étendue, variance et écart type. Tracer la boîte à moustaches et interpréter l’écart interquartile.
\begin{corrige}[exercice~\ref{t11s-ds4-e02}]\label{sol-t11s-ds4-e02}
Moyenne $8{,}75$, médiane $9$ (1). Quartiles $6,10$ (1), étendue $10$ (1). Somme des carrés $692$, variance $692/8-(8{,}75)^2=9{,}9375$, écart type environ $3{,}152$ (1). Boîte entre $6,10$, trait $9$, moustaches $4,14$ (1). Écart interquartile $4$ : largeur entre les deux quartiles; avec la convention de rang, au moins la moitié des observations se situe dans $[6;10]$ (1).
\end{corrige}
\exo[Comités et probabilités]\label{t11s-ds4-e03}\bareme{6 points}
Parmi $6$ filles et $4$ garçons, on choisit uniformément un comité de trois élèves.
\begin{enumerate}\item Compter les comités et les comités avec exactement deux filles (2 points).
\item Calculer la probabilité d’au moins une fille (2 points).
\item Choisir un président et un secrétaire distincts parmi les dix élèves donne-t-il le même nombre de résultats que choisir deux élèves sans fonction ? Justifier les deux comptes (2 points).\end{enumerate}
\begin{corrige}[exercice~\ref{t11s-ds4-e03}]\label{sol-t11s-ds4-e03}
1. Total $\binom{10}3=120$; exactement deux filles $\binom62\binom41=60$ (1+1).
2. Sans fille $\binom43=4$, donc $1-4/120=29/30$ (compte 1, complément 1).
3. Fonctions distinctes : $10\times9=90$; groupe de deux : $\binom{10}2=45$, car chaque groupe a deux ordres (1+1).
\end{corrige}
\exo[Réunion d’événements]\label{t11s-ds4-e04}\bareme{4 points}
Pour un dé équilibré, $A$ signifie « résultat pair » et $B$ « résultat supérieur ou égal à $4$ ». Calculer les probabilités de $A$, $B$, $A\cap B$ et $A\cup B$, puis vérifier la formule de réunion.
\begin{corrige}[exercice~\ref{t11s-ds4-e04}]\label{sol-t11s-ds4-e04}
$A=\{2,4,6\}$, $B=\{4,5,6\}$ : $1/2,1/2,1/3,2/3$, puis $1/2+1/2-1/3=2/3$ (1 par étape, formule incluse avec la réunion).
\end{corrige}
\endgroup


% ===========================================================================
\appendix
\ifbmtprof\clearpage\bmtchapitre{Progression indicative}{Organiser la dépendance des notions et adapter les séances au calendrier de l’établissement.}{Édition du professeur}
Les quatre parties du livre sont des regroupements éditoriaux. Elles n’imposent pas de commencer par l’analyse. La référence PE T11 conservée au dépôt attribue $65$ h à sa composante Analyse (avec équations, arithmétique et boucles), $20$ h à Algèbre (logique et suites), $45$ h à Géométrie et $20$ h aux données et probabilités : $150$ h au total. Les heures d’évaluation et de remédiation se répartissent dans ces enveloppes.
\begin{center}\begin{tabular}{L{2cm}L{9cm}}\toprule
Séquence&Chapitres et dépendances\\\midrule
1&Logique (\ref{t11s-c06}), boucles (\ref{t11s-c07}), équations et systèmes (\ref{t11s-c08}–\ref{t11s-c09}).\\
2&Arithmétique et bases (\ref{t11s-c10}–\ref{t11s-c11}); vecteurs et produit scalaire (\ref{t11s-c13}), puis trigonométrie (\ref{t11s-c16}).\\
3&Fonctions, limites, branches, dérivation et étude (\ref{t11s-c01}–\ref{t11s-c05}).\\
4&Suites (\ref{t11s-c12}); barycentres et transformations (\ref{t11s-c14}–\ref{t11s-c15}).\\
5&Dénombrement (\ref{t11s-c19}), espace (\ref{t11s-c17}), probabilités (\ref{t11s-c20}) et statistique (\ref{t11s-c18}).\\\bottomrule
\end{tabular}\end{center}
La RAPE archivée concerne 2025–2026; les dates de l’année 2026–2027 ne sont pas reprises sans document actuel vérifié. Introduire les fonctions trigonométriques après les rappels du cercle. Faire les DS de partie quand tous leurs prérequis sont installés, même si les chapitres ont été enseignés dans un autre ordre.
\section{Remédiation}
Après une erreur, identifier sa nature : calcul, domaine, choix de méthode, justification ou lecture. Reprendre un exemple court puis demander un exercice de structure différente. Pour une classe nombreuse, comparer deux solutions au tableau et faire justifier la différence; éviter une correction réduite au résultat final.
\fi
\clearpage\bmtchapitre{Formulaire}{Retrouver une formule avec ses conditions d’utilisation.}{Aide à la révision}
\section{Analyse}
Tangente : $y=f'(a)(x-a)+f(a)$, si $f$ est dérivable en $a$.
Pour $u,v$ dérivables : $(uv)'=u'v+uv'$, $(u/v)'=(u'v-uv')/v^2$ si $v\ne0$.
$\sqrt{x}$ a pour dérivée $1/(2\sqrt{x})$ sur $]0;+\infty[$.
Asymptote $y=ax+b$ : $f(x)-(ax+b)\to0$ en l’infini considéré.
\section{Algèbre}
Second degré ($a\ne0$) : $\Delta=b^2-4ac$, racines $(-b\pm\sqrt\Delta)/(2a)$ si $\Delta\ge0$.
Arithmétique : $u_n=u_p+(n-p)r$; somme de $k$ termes $k(\text{premier}+\text{dernier})/2$.
Géométrique ($q\ne0$) : $u_n=u_pq^{n-p}$; somme $u_p(1-q^k)/(1-q)$ si $q\ne1$, $ku_p$ si $q=1$. Si $q=0$, traiter le terme initial séparément.
Pour deux entiers positifs : $\pgcd(a,b)\ppcm(a,b)=ab$.
\section{Géométrie}
Barycentre : $\sum\alpha_i\vect{MA_i}=s\vect{MG}$ avec $s=\sum\alpha_i\ne0$.
$\sum\alpha_i MA_i^2=sMG^2+\sum\alpha_i GA_i^2$.
Al-Kashi : $BC^2=AB^2+AC^2-2AB\,AC\cos\widehat A$.
Produit vectoriel en repère orthonormé direct :
\[
(u_2v_3-u_3v_2\,;\,u_3v_1-u_1v_3\,;\,u_1v_2-u_2v_1).
\]
$\cos(a+b)=\cos a\cos b-\sin a\sin b$;
$\sin(a+b)=\sin a\cos b+\cos a\sin b$.
$\sin2x=2\sin x\cos x$, $\cos2x=2\cos^2x-1$.
\section{Données et probabilités}
Taux : $(V_f-V_i)/V_i$, avec $V_i>0$. Coefficients successifs : produit des $1+t$.
Pour $N=\sum n_i>0$ : $\bar x=\sum n_ix_i/N$, $V=\sum n_i(x_i-\bar x)^2/N$, $\sigma=\sqrt V$.
$A_n^p=n!/(n-p)!$, $\binom np=n!/[p!(n-p)!]$, $0\le p\le n$.
Équiprobabilité finie : $\Prob(A)=|A|/|\Omega|$, $\Omega\ne\varnothing$.
$\Prob(\overline A)=1-\Prob(A)$ et $\Prob(A\cup B)=\Prob(A)+\Prob(B)-\Prob(A\cap B)$.

\clearpage\bmtchapitre{Glossaire}{Relier un terme à son sens précis.}{Repères de vocabulaire}
\begin{description}[style=nextline,leftmargin=1em]
\item[Antécédent]Élément $x$ du domaine tel que $f(x)=y$; un même $y$ peut en avoir plusieurs.
\item[Asymptote]Droite dont l’écart à la courbe vérifie la limite donnée au chapitre~\ref{t11s-c03}.
\item[Barycentre]Point unique annulant une somme de vecteurs pondérés, lorsque la somme des poids est non nulle.
\item[Bijection]Application injective et surjective.
\item[Contre-exemple]Cas appartenant au domaine et réfutant une proposition universelle.
\item[Contraposée]Pour $P\Rightarrow Q$, l’implication $\neg Q\Rightarrow\neg P$, logiquement équivalente.
\item[Dérivée]Fonction des nombres dérivés, sur les points où ils existent.
\item[Équiprobabilité]Égalité des probabilités de toutes les issues du modèle fini considéré.
\item[Invariant]Propriété conservée à chaque tour d’un algorithme.
\item[Médiane]Valeur centrale de la série ordonnée selon la convention précisée.
\item[Monotone]Croissante ou décroissante sur l’intervalle, ou sur les rangs, précisés.
\item[Paramètre]Valeur fixée pendant la résolution, dont on étudie ensuite les différents cas.
\item[Point fixe]Point dont l’image par une transformation est lui-même.
\item[Produit vectoriel]Vecteur orthogonal aux deux facteurs; son sens dépend de l’ordre et de l’orientation du repère.
\item[Réciproque]Pour $P\Rightarrow Q$, l’implication $Q\Rightarrow P$, dont la vérité est à établir séparément.
\item[Variance]Moyenne des carrés des écarts à la moyenne, exprimée en unités au carré.
\end{description}

\clearpage\begingroup
\chapter*{Sujet de synthèse A}\markboth{Sujet de synthèse A}{Sujet de synthèse A}\addcontentsline{toc}{chapter}{Sujet de synthèse A}
\renewcommand{\thebmtexo}{SYN1.\arabic{bmtexo}}
\renewcommand{\theHbmtexo}{SYN1.\arabic{bmtexo}}
\setcounter{bmtexo}{0}
\begin{devoir}[3 h — 20 points]
Calculatrice autorisée. Justifier les réponses et préciser les domaines utiles. Les questions sont indépendantes sauf indication. 
\end{devoir}
\exo[Fonction et équation]\label{t11s-syn1-e01}\bareme{8 points}
$f(x)=x+2+4/x$, pour $x\ne0$.
\begin{enumerate}\item Limites et asymptotes (2 points).
\item Variations et extrema (2 points).
\item Résoudre $f(x)=7$ (2 points).
\item Donner le centre de symétrie et la tangente en $1$ (2 points).\end{enumerate}
\begin{corrige}[exercice~\ref{t11s-syn1-e01}]\label{sol-t11s-syn1-e01}
1. Asymptotes $x=0,y=x+2$; limites en $0^-:-\infty$, $0^+:+\infty$, aux infinis $-\infty,+\infty$.
2. $f'=1-4/x^2$, positive pour $|x|>2$, négative pour $0<|x|<2$. Maximum local $(-2;-2)$, minimum local $(2;6)$.
3. $x+2+4/x=7\Leftrightarrow x^2-5x+4=0$, donc $1,4$.
4. Centre $(0;2)$ car $f(x)+f(-x)=4$. $f(1)=7,f'(1)=-3$, tangente $y=-3x+10$. Deux points par sous-question, avec justification.
\end{corrige}
\exo[Suites]\label{t11s-syn1-e02}\bareme{6 points}
$u_0=2$, $u_{n+1}=3u_n-4$. Poser $v_n=u_n-2$ et étudier la suite. Puis refaire l’étude avec $u_0=3$ et calculer $u_0+\cdots+u_4$.
\begin{corrige}[exercice~\ref{t11s-syn1-e02}]\label{sol-t11s-syn1-e02}
$v_{n+1}=3v_n$. Si $u_0=2$, $v_0=0$, donc suite constante $u_n=2$ (2). Si $u_0=3$, $v_0=1$, $u_n=2+3^n$; croissante vers $+\infty$ (2). Somme des cinq termes $10+(1+3+9+27+81)=131$ (2).
\end{corrige}
\exo[Géométrie et probabilité]\label{t11s-syn1-e03}\bareme{6 points}
\begin{enumerate}\item Pour $A(0;0),B(4;0)$, déterminer le lieu $\vect{MA}\cdot\vect{MB}=5$ (2 points).
\item Donner l’image de ce lieu par l’homothétie de centre $A$ et rapport $2$ (2 points).
\item Choisir uniformément deux sommets parmi les quatre sommets d’un carré. Quelle est la probabilité qu’ils soient opposés ? (2 points)\end{enumerate}
\begin{corrige}[exercice~\ref{t11s-syn1-e03}]\label{sol-t11s-syn1-e03}
1. $I(2;0)$, $MI^2-4=5$ : cercle de rayon $3$.
2. Centre image $(4;0)$, rayon $6$.
3. Six paires de sommets, deux diagonales, donc probabilité $1/3$. Deux points par sous-question.
\end{corrige}
\endgroup

\begingroup
\chapter*{Sujet de synthèse B}\markboth{Sujet de synthèse B}{Sujet de synthèse B}\addcontentsline{toc}{chapter}{Sujet de synthèse B}
\renewcommand{\thebmtexo}{SYN2.\arabic{bmtexo}}
\renewcommand{\theHbmtexo}{SYN2.\arabic{bmtexo}}
\setcounter{bmtexo}{0}
\begin{devoir}[3 h — 20 points]
Calculatrice autorisée. Les repères utilisés pour les calculs métriques sont orthonormés; ils sont directs pour le produit vectoriel. Justifier les réponses et préciser les domaines utiles. Les questions sont indépendantes sauf indication. 
\end{devoir}
\exo[Algèbre et numération]\label{t11s-syn2-e01}\bareme{6 points}
\begin{enumerate}\item Résoudre $x^3-2x^2-x+2=0$ en utilisant une racine vérifiée (2 points).
\item Calculer PGCD et PPCM de $180$ et $252$ (2 points).
\item Convertir $(110101)_2$ en base dix puis en base cinq (2 points).\end{enumerate}
\begin{corrige}[exercice~\ref{t11s-syn2-e01}]\label{sol-t11s-syn2-e01}
1. $1$ est racine; factorisation $(x-1)(x-2)(x+1)$, solutions $-1,1,2$.
2. Décompositions $2^2 3^2 5$ et $2^2 3^2 7$, PGCD $36$, PPCM $1260$.
3. Valeur $53$, soit $(203)_5$. Deux points par sous-question.
\end{corrige}
\exo[Trigonométrie et espace]\label{t11s-syn2-e02}\bareme{8 points}
\begin{enumerate}\item Résoudre $\cos x=\sqrt2/2$ sur $[-\pi;\pi]$ (2 points).
\item Calculer $\vecu\wedge\vecv$ pour $\vecu=(1;0;1),\vecv=(0;2;1)$ et l’aire du triangle engendré (2 points).
\item Déterminer le plan passant par l’origine et dirigé par ces vecteurs (2 points).
\item Trouver son intersection avec la droite $(x;y;z)=(1;0;t)$ (2 points).\end{enumerate}
\begin{corrige}[exercice~\ref{t11s-syn2-e02}]\label{sol-t11s-syn2-e02}
1. $x=\pm\pi/4$.
2. Vectoriel $(-2;-1;2)$, norme $3$, aire $3/2$.
3. Normal non nul, équation $-2x-y+2z=0$.
4. $-2+2t=0$, $t=1$, point $(1;0;1)$, unique. Deux points par sous-question.
\end{corrige}
\exo[Données et tirage]\label{t11s-syn2-e03}\bareme{6 points}
\begin{enumerate}\item Pour les valeurs $5,10,15$ d’effectifs $2,5,3$, calculer moyenne et écart type (2 points).
\item Un prix augmente de $10\%$ deux fois. Calculer le taux global (1 point).
\item Tirer uniformément trois boules simultanément parmi quatre rouges et trois bleues; calculer la probabilité d’exactement deux rouges (2 points).
\item Donner celle d’au moins une rouge (1 point).\end{enumerate}
\begin{corrige}[exercice~\ref{t11s-syn2-e03}]\label{sol-t11s-syn2-e03}
1. Moyenne $10{,}5$, variance $12{,}25$, écart type $3{,}5$.
2. Coefficient $1{,}1^2=1{,}21$, hausse $21\%$.
3. Total $35$, favorables $18$, probabilité $18/35$.
4. Contraire : trois bleues, une combinaison; donc $1-1/35=34/35$.
\end{corrige}
\endgroup



\backmatter
\end{document}
